% =====================================================================
%  Anatomía de un LLM — Fase 0 — Módulo 0.5
%  Cálculo diferencial de una variable
%  Compilar: latexmk -r ../latexmkrc P05-calculo-diferencial.tex
% =====================================================================
\documentclass[11pt,a4paper,twoside,openany]{book}
\usepackage{../estilo-apuntes}

\newcommand{\ApuntesCurso}{Anatomía de un LLM}
\newcommand{\ApuntesModulo}{Módulo 0.5}

% Notación propia del módulo
\newcommand{\N}{\mathbb{N}}
\newcommand{\Z}{\mathbb{Z}}
\newcommand{\Q}{\mathbb{Q}}
\newcommand{\dom}{\operatorname{dom}}
\newcommand{\ReLU}{\operatorname{ReLU}}
\newcommand{\softplus}{\operatorname{softplus}}
\newcommand{\logit}{\operatorname{logit}}
\newcommand{\arsinh}{\operatorname{arsinh}}
\newcommand{\artanh}{\operatorname{artanh}}
\newcommand{\sgn}{\operatorname{sgn}}
\newcommand{\GELU}{\operatorname{GELU}}
\newcommand{\dd}{\mathrm{d}}
\DeclarePairedDelimiter{\floor}{\lfloor}{\rfloor}

% Anclas de hiperenlace únicas para los enunciados numerados por tema
\makeatletter
\@for\p@entorno:=definicion,teorema,proposicion,lema,corolario,ejemplo,observacion\do{%
  \expandafter\def\csname theH\p@entorno\endcsname{\thechapter.\arabic{definicion}}}
\makeatother

\hypersetup{pdftitle={Módulo 0.5. Cálculo diferencial de una variable},pdfauthor={Anatomía de un LLM}}

\begin{document}
\frontmatter
% ------------------------------------------------------------------ Portada
\begin{titlepage}
\centering
\vspace*{2.5cm}
{\large\scshape Anatomía de un LLM}\par
\vspace{0.4cm}
{\normalsize Fundamentos matemáticos de los modelos de lenguaje}\par
\vspace{2.2cm}
\rule{\textwidth}{0.8pt}\par\vspace{0.5cm}
{\small\scshape Fase 0 \textperiodcentered\ Preliminares matemáticos}\par\vspace{0.3cm}
{\Huge\bfseries Módulo 0.5\par}\vspace{0.3cm}
{\LARGE Cálculo diferencial de una variable\par}
\vspace{0.5cm}\rule{\textwidth}{0.4pt}\par
\vspace{1.6cm}
{\large Apuntes del curso}\par
\vfill
{\small Edición de octubre de 2026}\par
\end{titlepage}

% ------------------------------------------------------- Presentación
\chapter*{Presentación del módulo}
\addcontentsline{toc}{chapter}{Presentación del módulo}
\markboth{Presentación del módulo}{}

Entrenar un modelo de lenguaje consiste en ajustar millones de parámetros para que una función de pérdida disminuya, y casi todas las herramientas que lo hacen posible proceden del cálculo diferencial: la derivada indica en qué dirección cambiar un parámetro, el teorema del valor medio acota cuánto cambia la pérdida en un paso, el desarrollo de Taylor justifica las aproximaciones lineales y cuadráticas, y los límites describen el comportamiento de la softmax cuando la temperatura tiende a cero o a infinito. Este módulo desarrolla el cálculo diferencial de una variable desde el nivel de bachillerato hasta el nivel de un primer curso universitario de análisis: con definiciones precisas de límite y continuidad, demostraciones de los teoremas fundamentales y una colección muy extensa de ejercicios, tanto de cálculo mecánico como de demostración. Las funciones de activación de las redes neuronales (sigmoide, tangente hiperbólica, ReLU, softplus) aparecen a lo largo del módulo como ejemplos recurrentes.

\section*{Objetivos}
Al finalizar el módulo, el estudiante debe ser capaz de:
\begin{enumerate}[label=\arabic*.]
  \item manejar las funciones reales elementales, incluidas las hiperbólicas y las funciones de activación, y razonar sobre su dominio, monotonía, paridad e inversa;
  \item demostrar límites de sucesiones y de funciones con las definiciones $\varepsilon$--$N$ y $\varepsilon$--$\delta$, y calcular límites con el álgebra de límites, el teorema del sándwich y los límites notables;
  \item demostrar y aplicar los teoremas de Bolzano, de los valores intermedios y de Weierstrass;
  \item calcular derivadas con las reglas de la suma, el producto, el cociente, la cadena y la función inversa, y mediante derivación implícita y logarítmica, y demostrar esas reglas a partir de la definición;
  \item demostrar los teoremas de Rolle, del valor medio y de Cauchy, y usarlos para probar desigualdades, estudiar la monotonía, calcular límites con la regla de L'Hôpital y resolver problemas de optimización y de convexidad;
  \item construir polinomios de Taylor, acotar el error con el resto de Lagrange y usar los desarrollos para aproximar valores y calcular límites.
\end{enumerate}

\section*{Requisitos previos}
Matemáticas de bachillerato: funciones elementales y sus gráficas, cálculo intuitivo de límites y derivadas. Se usan además los módulos anteriores de esta fase: el módulo~0.1 (cuantificadores, contrarrecíproco, reducción al absurdo e inducción), el módulo~0.2 (funciones inyectivas, sobreyectivas e inversas, composición), el módulo~0.3 (axiomas de los números reales y supremo, valor absoluto y desigualdad triangular, desigualdad de Bernoulli, exponenciales y logaritmos) y el módulo~0.4 (funciones trigonométricas y sus inversas).

\section*{Organización}
\begin{center}
\begin{tabular}{@{}cll@{}}
\toprule
Tema & Contenido & Fuentes principales \\
\midrule
1 & Funciones reales elementales y sus gráficas & \cite{stewart}, \cite{osc1}, \cite{spivak} \\
2 & Sucesiones y límites & \cite{spivak}, \cite{cambA1}, \cite{osc1} \\
3 & Continuidad & \cite{spivak}, \cite{apostol}, \cite{cambA1} \\
4 & La derivada y las reglas de derivación & \cite{stewart}, \cite{mit1801}, \cite{spivak} \\
5 & Teoremas del valor medio y aplicaciones & \cite{spivak}, \cite{mit1801}, \cite{osc1} \\
6 & Polinomio de Taylor & \cite{spivak}, \cite{apostol}, \cite{mit1801} \\
\bottomrule
\end{tabular}
\end{center}

Cada tema termina con una lista de lecturas y con una colección extensa de ejercicios graduados en tres niveles: $\star$ comprensión, $\star\star$ cálculo o práctica mecánica y $\star\star\star$ demostración. Los ejercicios marcados con \emph{Fuente} están adaptados de cursos y libros de texto universitarios; el enunciado se ha redactado de nuevo y, en algunos casos, se han cambiado los datos. Las soluciones completas, paso a paso, están en el \cref{ap:soluciones}. En los ejercicios de demostración, el apartado «Desarrollo» de la solución contiene la demostración redactada tal como debería entregarse.

\section*{Dónde se usa en el curso}
\begin{center}
\small
\begin{tabularx}{\textwidth}{@{}>{\raggedright\arraybackslash}p{5.0cm}>{\raggedright\arraybackslash}X@{}}
\toprule
Concepto de este módulo & Módulos del curso que lo emplean \\
\midrule
Sigmoide, tangente hiperbólica, ReLU, softplus & M08 (regresión logística), M09 (funciones de activación), M11 (saturación y gradientes que se desvanecen), M14 y M18 (bloques del Transformer) \\
Límites y definición $\varepsilon$--$\delta$ & M01 (límites de la softmax con temperatura), M04 (derivadas parciales), M07 (convergencia de algoritmos) \\
Sucesiones monótonas y acotadas, el número $e$ & M07 (medias exponenciales, convergencia del descenso de gradiente), M21 (retorno descontado) \\
Continuidad, Bolzano y Weierstrass & M04 y M07 (existencia de mínimos), M09 (aproximación universal en una dimensión) \\
Reglas de derivación, regla de la cadena & M04 (gradiente y jacobiano), M09 y M10 (retropropagación y diferenciación automática) \\
$\sigma'=\sigma(1-\sigma)$, $\tanh'=1-\tanh^{2}$ & M08, M09, M11 (gradientes de las activaciones) \\
Teorema del valor medio, funciones lipschitzianas & M07 (lema de descenso), M11 (recorte del gradiente), M17 (cotas de error) \\
Monotonía, extremos y convexidad & M04 y M07 (convexidad, mínimos), M06 (desigualdad de Gibbs) \\
Regla de L'Hôpital y órdenes de crecimiento & M01 (temperatura), M06 ($0\ln 0=0$ en la entropía), M16 (leyes de potencia) \\
Taylor con resto de Lagrange & M04 (Taylor de segundo orden, diferencias finitas), M07 (análisis del descenso de gradiente), M11 (aproximaciones numéricas, $\ln(1+x)$) \\
\bottomrule
\end{tabularx}
\end{center}

\section*{Sobre las fuentes}
El texto es de redacción propia. La organización sigue la de los capítulos~3 a~12, 18, 20 y~22 de \emph{Calculus}, de Spivak~\cite{spivak}, el texto de referencia para un primer curso riguroso de cálculo, y la de los capítulos~3, 4, 6 y~7 del primer volumen de Apostol~\cite{apostol}. El nivel de las demostraciones de los temas~2, 3, 5 y~6 es el del curso \emph{Analysis~I} de la Universidad de Cambridge~\cite{cambA1}, cuyos exámenes~\cite{tripos2013,tripos2019} se usan como fuente de varios ejercicios. La práctica de cálculo mecánico sigue a Stewart~\cite{stewart}, al libro libre \emph{Calculus} de OpenStax~\cite{osc1,osc2} y a los problemas y exámenes del curso 18.01SC del MIT~\cite{mit1801,mit1801ex}.

\section*{Notación}
\begin{center}
\small
\begin{tabular}{@{}ll@{}}
\toprule
Símbolo & Significado \\
\midrule
$\N=\{1,2,3,\dots\}$, $\Z$, $\Q$, $\R$ & naturales (sin el cero), enteros, racionales y reales \\
$\dom f$ & dominio de la función $f$ \\
$(a_n)_{n\ge1}$, $a_n\to L$ & sucesión y convergencia de una sucesión \\
$\lim_{x\to a}f(x)$, $\lim_{x\to a^{\pm}}f(x)$ & límite y límites laterales \\
$f'(a)$, $\dfrac{\dd f}{\dd x}(a)$, $f^{(n)}$ & derivada en $a$ y derivada $n$-ésima \\
$\sigma(x)=1/(1+e^{-x})$ & función sigmoide (logística) \\
$\ReLU(x)=\max(0,x)$ & unidad lineal rectificada \\
$\softplus(x)=\ln(1+e^{x})$ & softplus, versión suave de la ReLU \\
$\sinh$, $\cosh$, $\tanh$ & funciones hiperbólicas \\
$\ln$ & logaritmo natural (neperiano) \\
$P_{n,a}f$, $R_{n,a}f$ & polinomio de Taylor de orden $n$ en $a$ y su resto \\
$o(g(x))$ & función $h$ con $h(x)/g(x)\to0$ \\
\bottomrule
\end{tabular}
\end{center}

\tableofcontents

\mainmatter
% =====================================================================
\chapter{Funciones reales elementales y sus gráficas}
\label{tema:1}
% =====================================================================

Todo el cálculo trata de funciones reales de una variable real. Este tema fija el vocabulario (dominio, imagen, gráfica, monotonía, paridad, inversa), repasa el catálogo de funciones elementales que se usarán como ejemplos en el resto del módulo y presenta las funciones de activación de las redes neuronales como funciones elementales más. Se siguen las §1.1--1.5 de Stewart~\cite{stewart}, las §1.1--1.5 de OpenStax~\cite{osc1} y los capítulos~3 y~4 de Spivak~\cite{spivak}; la parte conjuntista (inyectividad, composición e inversa) se estudió en el módulo~0.2.

\section{Funciones reales de variable real}

\begin{definicion}[Función real, dominio, imagen y gráfica]\label{def:funcion}
Una \emph{función real de variable real} es una función $f\colon D\to\R$ con $D\subseteq\R$. El conjunto $D=\dom f$ es su \emph{dominio}; $f(D)=\{f(x):x\in D\}$ es su \emph{imagen} o recorrido, y su \emph{gráfica} es el conjunto $\{(x,f(x)):x\in D\}\subseteq\R^{2}$. Cuando una función se da por una fórmula sin indicar el dominio, se sobrentiende el \emph{dominio natural}: el mayor conjunto de reales en el que la fórmula tiene sentido.
\end{definicion}

\begin{ejemplo}[Dominios naturales]\label{ej:dominios}
(a) $f(x)=\sqrt{4-x^{2}}$: se necesita $4-x^{2}\ge0$, luego $\dom f=[-2,2]$ y la imagen es $[0,2]$ (semicircunferencia superior de radio~2). (b) $g(x)=\ln(x-1)/(x-3)$: se necesita $x>1$ y $x\ne3$, luego $\dom g=(1,3)\cup(3,\infty)$. (c) $h(x)=1/(1+e^{-x})$: el denominador es siempre mayor que~1, luego $\dom h=\R$.
\end{ejemplo}

Las funciones se suman, se multiplican y se dividen punto a punto, en la intersección de los dominios (y, para el cociente, donde el denominador no se anula). La \emph{composición} $(g\circ f)(x)=g(f(x))$ está definida en $\{x\in\dom f: f(x)\in\dom g\}$.

\begin{definicion}[Acotación, monotonía, paridad y periodicidad]\label{def:propiedades}
Sea $f\colon D\to\R$.
\begin{enumerate}
  \item $f$ está \emph{acotada superiormente} si existe $M$ con $f(x)\le M$ para todo $x\in D$; \emph{inferiormente}, si existe $m$ con $f(x)\ge m$; \emph{acotada}, si lo está de ambas formas, es decir, si existe $K$ con $\abs{f(x)}\le K$ para todo $x\in D$.
  \item $f$ es \emph{creciente} si $x<y\Rightarrow f(x)\le f(y)$, y \emph{estrictamente creciente} si $x<y\Rightarrow f(x)<f(y)$, para todos $x,y\in D$. Análogamente, (estrictamente) \emph{decreciente}. Se llama \emph{monótona} a una función creciente o decreciente.
  \item Si $D$ es simétrico ($x\in D\Rightarrow -x\in D$), $f$ es \emph{par} si $f(-x)=f(x)$ e \emph{impar} si $f(-x)=-f(x)$ para todo $x\in D$.
  \item $f$ es \emph{periódica} de periodo $T>0$ si $x+T\in D$ y $f(x+T)=f(x)$ para todo $x\in D$.
\end{enumerate}
\end{definicion}

La gráfica de una función par es simétrica respecto del eje de ordenadas, y la de una impar, respecto del origen. Una función impar definida en~$0$ cumple $f(0)=-f(0)$, luego $f(0)=0$.

\begin{proposicion}[Transformaciones de gráficas]\label{prop:transformaciones}
Sea $f$ una función, $c\in\R$ y $k>0$. La gráfica de
\begin{enumerate}
  \item $y=f(x)+c$ es la de $f$ trasladada $c$ unidades en vertical;
  \item $y=f(x-c)$ es la de $f$ trasladada $c$ unidades en horizontal (hacia la derecha si $c>0$);
  \item $y=kf(x)$ es la de $f$ dilatada verticalmente por el factor $k$;
  \item $y=f(kx)$ es la de $f$ comprimida horizontalmente por el factor $k$ (dilatada si $k<1$);
  \item $y=-f(x)$ e $y=f(-x)$ son sus simétricas respecto del eje de abscisas y del de ordenadas.
\end{enumerate}
\end{proposicion}

\begin{proof}
Se prueba (ii); las demás son análogas. $(x,y)$ está en la gráfica de $y=f(x-c)$ si y solo si $y=f(x-c)$, es decir, si y solo si $(x-c,y)$ está en la gráfica de $f$. Luego la nueva gráfica se obtiene sumando $c$ a la abscisa de cada punto de la gráfica de~$f$.
\end{proof}

\section{El catálogo de funciones elementales}

\begin{definicion}[Funciones elementales básicas]\label{def:elementales}
Se llaman \emph{elementales} las funciones que se obtienen, mediante sumas, productos, cocientes, composiciones e inversas, a partir de:
\begin{enumerate}
  \item los \emph{polinomios} $p(x)=a_nx^{n}+\dots+a_1x+a_0$ y las \emph{funciones racionales} $p/q$ (cocientes de polinomios, en los puntos donde $q\ne0$);
  \item las \emph{potencias} $x^{a}$ ($x>0$, $a\in\R$) y las raíces $\sqrt[n]{x}$;
  \item la \emph{exponencial} $\exp(x)=e^{x}$, biyección creciente de $\R$ en $(0,\infty)$, y su inversa, el \emph{logaritmo} $\ln\colon(0,\infty)\to\R$;
  \item las funciones trigonométricas $\sin$, $\cos$, $\tan$ y sus inversas $\arcsin$, $\arccos$, $\arctan$ (módulo~0.4).
\end{enumerate}
\end{definicion}

Se recuerdan las propiedades de la exponencial y del logaritmo (módulo~0.3): $e^{x+y}=e^{x}e^{y}$, $e^{0}=1$, $\ln(xy)=\ln x+\ln y$, $\ln(x^{a})=a\ln x$, $e^{\ln x}=x$ para $x>0$ y $\ln(e^{x})=x$ para todo $x$; además $a^{x}=e^{x\ln a}$ para $a>0$. El valor absoluto $\abs{x}=\max(x,-x)$ y la parte entera $\floor{x}$ (el mayor entero $\le x$) completan el repertorio de ejemplos.

\begin{definicion}[Funciones hiperbólicas]\label{def:hiperbolicas}
Para $x\in\R$ se definen
\[
\sinh x=\frac{e^{x}-e^{-x}}{2},\qquad \cosh x=\frac{e^{x}+e^{-x}}{2},\qquad \tanh x=\frac{\sinh x}{\cosh x}=\frac{e^{x}-e^{-x}}{e^{x}+e^{-x}}.
\]
\end{definicion}

\begin{proposicion}[Propiedades de las funciones hiperbólicas]\label{prop:hiperbolicas}
Para todo $x\in\R$:
\begin{enumerate}
  \item $\cosh$ es par, $\cosh x\ge1$, y $\sinh$ y $\tanh$ son impares.
  \item $\cosh^{2}x-\sinh^{2}x=1$ para todo $x$.
  \item $\cosh x+\sinh x=e^{x}$ y $\cosh x-\sinh x=e^{-x}$.
  \item $\abs{\tanh x}<1$, y $\tanh$ es estrictamente creciente de $\R$ en $(-1,1)$.
\end{enumerate}
\end{proposicion}

\begin{proof}
(i) Al cambiar $x$ por $-x$ se intercambian $e^{x}$ y $e^{-x}$; la suma no cambia y la diferencia cambia de signo. Por la desigualdad entre las medias aritmética y geométrica, $\cosh x\ge\sqrt{e^{x}e^{-x}}=1$. (iii) Es inmediato sumando y restando las definiciones. (ii) Por (iii), $\cosh^{2}x-\sinh^{2}x=(\cosh x+\sinh x)(\cosh x-\sinh x)=e^{x}e^{-x}=1$. (iv) Multiplicando numerador y denominador por $e^{x}$, $\tanh x=\dfrac{e^{2x}-1}{e^{2x}+1}=1-\dfrac{2}{e^{2x}+1}$. Como $e^{2x}+1>1$, se tiene $-1<\tanh x<1$; y como $x\mapsto e^{2x}+1$ es estrictamente creciente y positiva, $2/(e^{2x}+1)$ es estrictamente decreciente y $\tanh$ estrictamente creciente. Que la imagen es todo $(-1,1)$ se ve resolviendo $\tanh x=y$ (\cref{eje:1-arsinh}).
\end{proof}

El nombre procede de (ii): el punto $(\cosh t,\sinh t)$ recorre la rama derecha de la hipérbola $x^{2}-y^{2}=1$, como $(\cos t,\sin t)$ recorre la circunferencia $x^{2}+y^{2}=1$.

\section{Funciones de activación}

Una neurona artificial calcula una combinación lineal de sus entradas y le aplica una función no lineal fija, la \emph{función de activación}. Las más usadas son funciones elementales (o casi) y sirven como ejemplos recurrentes en este módulo.

\begin{definicion}[Funciones de activación]\label{def:activaciones}
Para $x\in\R$ se definen
\begin{enumerate}
  \item la \emph{sigmoide} o función logística: $\sigma(x)=\dfrac{1}{1+e^{-x}}$;
  \item la \emph{tangente hiperbólica} $\tanh x$ (\cref{def:hiperbolicas});
  \item la \emph{unidad lineal rectificada}: $\ReLU(x)=\max(0,x)$;
  \item la \emph{softplus}: $\softplus(x)=\ln(1+e^{x})$.
\end{enumerate}
\end{definicion}

\begin{figure}[h]
\centering
\begin{tikzpicture}[x=0.75cm,y=1.1cm]
  \foreach \y in {-1,1} {\draw[dashed,color=regla] (-4.2,\y)--(-0.55,\y); \draw[dashed,color=regla] (0.1,\y)--(4.2,\y);}
  \draw[->,color=regla] (-4.3,0)--(4.4,0) node[right,color=tinta]{\small $x$};
  \draw[->,color=regla] (0,-1.25)--(0,2.6);
  \foreach \t in {-4,-2,2,4} \draw[color=regla] (\t,0.04)--(\t,-0.04) node[below,color=tinta]{\scriptsize $\t$};
  \foreach \t in {-1,1,2} \draw[color=regla] (0.06,\t)--(-0.06,\t) node[left,color=tinta]{\scriptsize $\t$};
  \draw[thick,domain=-4.2:4.2,samples=80] plot (\x,{1/(1+exp(-\x))});
  \draw[thick,densely dashed,domain=-4.2:4.2,samples=80] plot (\x,{(exp(\x)-exp(-\x))/(exp(\x)+exp(-\x))});
  \draw[very thick,color=enlace] (-4.2,0)--(0,0)--(2.5,2.5);
  \draw[thick,densely dotted,domain=-4.2:2.0,samples=60] plot (\x,{ln(1+exp(\x))});
  \node[right] at (4.2,1.0) {\small $\sigma$};
  \node[left] at (-4.2,-1.0) {\small $\tanh$};
  \node[right,color=enlace] at (2.5,2.5) {\small $\ReLU$};
  \node[left] at (1.9,2.15) {\small $\softplus$};
\end{tikzpicture}
\caption{Gráficas de $\sigma$ (trazo continuo), $\tanh$ (discontinuo), $\ReLU$ (trazo grueso) y $\softplus$ (punteado).}\label{fig:activaciones}
\end{figure}

\begin{proposicion}[Propiedades de la sigmoide]\label{prop:sigmoide}
La función sigmoide cumple:
\begin{enumerate}
  \item $0<\sigma(x)<1$ para todo $x$, y $\sigma(0)=\tfrac12$.
  \item $\sigma(-x)=1-\sigma(x)$; en consecuencia, la gráfica es simétrica respecto del punto $(0,\tfrac12)$.
  \item $\sigma$ es estrictamente creciente y biyectiva de $\R$ en $(0,1)$, con inversa $\logit(p)=\ln\dfrac{p}{1-p}$.
  \item $\tanh x=2\sigma(2x)-1$.
\end{enumerate}
\end{proposicion}

\begin{proof}
(i) $1+e^{-x}>1$. (iii) $x\mapsto 1+e^{-x}$ es estrictamente decreciente y positiva, luego su inversa multiplicativa es estrictamente creciente. Dado $p\in(0,1)$, la ecuación $\sigma(x)=p$ equivale a $e^{-x}=\frac1p-1=\frac{1-p}{p}>0$, que tiene la solución única $x=\ln\frac{p}{1-p}$. Las demostraciones de (ii) y (iv) son el \cref{eje:1-sigmoide-simetria}.
\end{proof}

La propiedad (iv) dice que $\tanh$ es una sigmoide reescalada: misma forma, imagen $(-1,1)$ en lugar de $(0,1)$ y pendiente cuatro veces mayor en el origen (\cref{tema:4}). En la regresión logística (módulo~8), $\sigma$ transforma una puntuación real en una probabilidad, y $\logit$ deshace esa transformación: el logit de una probabilidad es el logaritmo de la razón de probabilidades $p/(1-p)$, de donde procede el término \emph{logits} con que se designan las salidas de una red antes de la softmax.

\begin{proposicion}[ReLU y softplus]\label{prop:relu}
Para todo $x\in\R$:
\begin{enumerate}
  \item $\ReLU(x)=\dfrac{x+\abs{x}}{2}$ y $\ReLU(x)-\ReLU(-x)=x$;
  \item $\softplus(x)-\softplus(-x)=x$;
  \item $\ReLU(x)<\softplus(x)\le\ReLU(x)+\ln 2$, con igualdad a la derecha solo en $x=0$.
\end{enumerate}
\end{proposicion}

\begin{proof}
(i) Si $x\ge0$, $(x+\abs{x})/2=x$; si $x<0$, vale~$0$. La segunda igualdad se sigue de la primera, porque $\abs{-x}=\abs{x}$. (ii) $\softplus(x)-\softplus(-x)=\ln\dfrac{1+e^{x}}{1+e^{-x}}=\ln\dfrac{e^{x}(e^{-x}+1)}{1+e^{-x}}=\ln e^{x}=x$. (iii) Si $x\le0$: $\softplus(x)=\ln(1+e^{x})$ con $0<e^{x}\le1$, luego $0<\softplus(x)\le\ln2$, con igualdad solo si $x=0$. Si $x>0$: por (ii), $\softplus(x)=x+\softplus(-x)$ y $0<\softplus(-x)<\ln 2$.
\end{proof}

La softplus es, por tanto, una versión «suavizada» de la ReLU que nunca se separa de ella más de $\ln 2\approx0{,}693$ (\cref{fig:activaciones}). Otras activaciones modernas, como la GELU, $\GELU(x)=x\,\Phi(x)$ con $\Phi$ la función de distribución de la normal estándar, no son elementales; se estudian en el módulo~9.

\section{Función inversa}

\begin{proposicion}[Monotonía estricta e inversa]\label{prop:inversa-monotona}
Sea $f\colon D\to\R$ estrictamente creciente. Entonces $f$ es inyectiva, su inversa $f^{-1}\colon f(D)\to D$ es estrictamente creciente, y la gráfica de $f^{-1}$ es la simétrica de la de $f$ respecto de la recta $y=x$.
\end{proposicion}

\begin{proof}
Si $x\ne y$, por ejemplo $x<y$, entonces $f(x)<f(y)$ y, en particular, $f(x)\ne f(y)$: $f$ es inyectiva y, por tanto, biyectiva sobre su imagen. Sean $u<v$ en $f(D)$, con $u=f(x)$ y $v=f(y)$. Si fuera $x\ge y$, el crecimiento de $f$ daría $u=f(x)\ge f(y)=v$, contradicción; luego $f^{-1}(u)=x<y=f^{-1}(v)$. Por último, $(a,b)$ está en la gráfica de $f$ si y solo si $b=f(a)$, si y solo si $a=f^{-1}(b)$, si y solo si $(b,a)$ está en la gráfica de $f^{-1}$; y $(a,b)\mapsto(b,a)$ es la simetría respecto de $y=x$.
\end{proof}

\begin{ejemplo}[Inversas de funciones elementales]\label{ej:inversas}
(a) $\exp\colon\R\to(0,\infty)$ y $\ln$ son inversas, ambas estrictamente crecientes. (b) $x\mapsto x^{2}$ no es inyectiva en $\R$, pero su restricción a $[0,\infty)$ es estrictamente creciente y su inversa es $\sqrt{\phantom{x}}$. (c) $\sigma$ y $\logit$ son inversas (\cref{prop:sigmoide}). (d) Para invertir $y=\frac{2x+1}{x-3}$ ($x\ne3$) se despeja: $y(x-3)=2x+1$, $x(y-2)=3y+1$, $x=\frac{3y+1}{y-2}$, válido para $y\ne2$; la imagen es $\R\setminus\{2\}$.
\end{ejemplo}

\begin{observacion}[Errores frecuentes]\label{obs:1-errores}
(a) Confundir $f^{-1}(x)$ (función inversa) con $1/f(x)=f(x)^{-1}$ (inverso multiplicativo): $\sin^{-1}$ suele denotar $\arcsin$, no $1/\sin$.
(b) Usar $\ln(a+b)=\ln a+\ln b$ o $e^{a+b}=e^{a}+e^{b}$, que son falsas; las correctas son $\ln(ab)=\ln a+\ln b$ y $e^{a+b}=e^{a}e^{b}$.
(c) Olvidar el dominio al componer: $\ln(\ln x)$ exige $x>1$, no solo $x>0$.
(d) Pensar que una función que no es par es impar: la mayoría no son ni lo uno ni lo otro (por ejemplo, $e^{x}$), aunque toda función de dominio simétrico es suma de una par y una impar (\cref{eje:1-par-impar}).
(e) En $y=f(x-c)$, trasladar la gráfica hacia la izquierda cuando $c>0$; el desplazamiento es hacia la derecha.
\end{observacion}

\begin{observacion}[Conexión con el curso]
La sigmoide y la tangente hiperbólica fueron las activaciones clásicas de las redes neuronales; su saturación (para $\abs{x}$ grande se aplanan, \cref{fig:activaciones}) provoca los gradientes que se desvanecen del módulo~11. La ReLU, que no satura para $x>0$, y sus variantes suaves (softplus, GELU, SiLU $x\sigma(x)$) son las activaciones de los perceptrones multicapa del módulo~9 y de los bloques del Transformer de los módulos~14 y~18. La función $\logit=\sigma^{-1}$ explica el nombre de las salidas de un LLM antes de la softmax (módulo~1).
\end{observacion}

\begin{lecturas}
\item Stewart~\cite{stewart}, §1.1 «Four Ways to Represent a Function», §1.3 «New Functions from Old Functions» y §1.5 «Inverse Functions and Logarithms».
\item OpenStax~\cite{osc1}, §1.1 «Review of Functions», §1.2 «Basic Classes of Functions», §1.4 «Inverse Functions» y §1.5 «Exponential and Logarithmic Functions» (incluye las funciones hiperbólicas).
\item Spivak~\cite{spivak}, cap.~3 «Functions» y cap.~4 «Graphs».
\item MIT 18.01SC~\cite{mit1801}, sesión~20 «Hyperbolic Trig Functions».
\end{lecturas}

\section*{Ejercicios del tema 1}
\addcontentsline{toc}{section}{Ejercicios del tema 1}

\begin{ejercicio}{1}{eje:1-dominios}
Determine el dominio natural de: (a) $f(x)=\sqrt{x^{2}-5x+6}$; (b) $g(x)=\ln\dfrac{1+x}{1-x}$; (c) $h(x)=\dfrac{1}{\ln x}$; (d) $k(x)=\sqrt{\ln x}$.
\end{ejercicio}

\begin{ejercicio}{1}{eje:1-transformaciones}
La gráfica de $f$ pasa por los puntos $(0,0)$, $(1,2)$ y $(3,-1)$. ¿Por qué puntos pasa la gráfica de: (a) $y=f(x-2)+1$; (b) $y=3f(x)$; (c) $y=f(2x)$; (d) $y=-f(-x)$? Describa en cada caso la transformación con la \cref{prop:transformaciones}.
\end{ejercicio}

\begin{ejercicio}{2}{eje:1-paridad}
Estudie si son pares, impares o ninguna de las dos cosas: (a) $x^{3}-x$; (b) $\abs{x}+\cos x$; (c) $e^{x}$; (d) $\ln\dfrac{1-x}{1+x}$ en $(-1,1)$; (e) $x\,\sigma(x)$; (f) $\softplus(x)-\dfrac{x}{2}$.
\end{ejercicio}

\begin{ejercicio}{2}{eje:1-composicion}
Sean $f(x)=\sqrt{x}$ y $g(x)=1-x^{2}$. Calcule $f\circ g$, $g\circ f$, $f\circ f$ y $g\circ g$, y el dominio de cada una.
\end{ejercicio}

\begin{ejercicio}{2}{eje:1-ecuaciones}
Resuelva las ecuaciones: (a) $e^{3x}-15=0$; (b) $7^{3x-2}=11$; (c) $\ln x+\ln(x-2)=\ln 4$; (d) $4^{x+1}-32=0$.
\fuente{adaptado de OpenStax~\cite{osc1}, §1.5, ejercicios de la sección (ecuaciones exponenciales y logarítmicas).}
\end{ejercicio}

\begin{ejercicio}{2}{eje:1-hiperbolicas-calculo}
Escriba con exponenciales y simplifique: (a) $2\cosh(\ln x)$ para $x>0$; (b) $\cosh 4x+\sinh 4x$; (c) $\cosh 2x-\sinh 2x$; (d) $\ln(\cosh x+\sinh x)+\ln(\cosh x-\sinh x)$.
\fuente{adaptado de OpenStax~\cite{osc1}, §1.5, ejercicio de la sección sobre funciones hiperbólicas.}
\end{ejercicio}

\begin{ejercicio}{2}{eje:1-sigmoide-valores}
(a) Calcule $\sigma(0)$, $\sigma(\ln 3)$ y $\sigma(-\ln 3)$. (b) Resuelva $\sigma(x)=0{,}9$ y $\sigma(x)=0{,}01$. (c) ¿Para qué valores de $x$ es $\sigma(x)>0{,}99$? (d) Calcule $\tanh(\ln 2)$.
\end{ejercicio}

\begin{ejercicio}{2}{eje:1-inversa-racional}
Sea $f(x)=\dfrac{x+1}{x-1}$ en $\R\setminus\{1\}$. (a) Calcule $f\circ f$ y deduzca que $f$ es su propia inversa. (b) Halle la imagen de $f$. (c) Explique geométricamente por qué la gráfica de $f$ es simétrica respecto de la recta $y=x$.
\end{ejercicio}

\begin{ejercicio}{3}{eje:1-sigmoide-simetria}
Demuestre las propiedades (ii) y (iv) de la \cref{prop:sigmoide}: $\sigma(-x)=1-\sigma(x)$ y $\tanh x=2\sigma(2x)-1$ para todo $x\in\R$. Deduzca que $\sigma(x)=\tfrac12\bigl(1+\tanh\tfrac{x}{2}\bigr)$.
\end{ejercicio}

\begin{ejercicio}{3}{eje:1-hiperbolicas-suma}
Demuestre que, para todos $x,y\in\R$:
(a) $\sinh(x+y)=\sinh x\cosh y+\cosh x\sinh y$;
(b) $\cosh(x+y)=\cosh x\cosh y+\sinh x\sinh y$;
(c) $\tanh(x+y)=\dfrac{\tanh x+\tanh y}{1+\tanh x\tanh y}$.
Compare con las fórmulas trigonométricas del módulo~0.4.
\fuente{adaptado de MIT 18.01SC~\cite{mit1801}, sesión~20 «Hyperbolic Trig Functions».}
\end{ejercicio}

\begin{ejercicio}{3}{eje:1-arsinh}
(a) Demuestre que $\sinh\colon\R\to\R$ es biyectiva y que su inversa es $\arsinh y=\ln\bigl(y+\sqrt{y^{2}+1}\bigr)$. (b) Demuestre que $\tanh\colon\R\to(-1,1)$ es biyectiva y que $\artanh y=\tfrac12\ln\dfrac{1+y}{1-y}$. (c) Compruebe que $\artanh y=\tfrac12\logit\bigl(\tfrac{1+y}{2}\bigr)$.
\end{ejercicio}

\begin{ejercicio}{3}{eje:1-estricta}
Sea $f\colon[a,b]\to\R$ estrictamente creciente. Demuestre que $f$ es inyectiva y que, si $f(x)<f(y)$, entonces $x<y$. Muestre con un ejemplo que una función inyectiva en $[0,2]$ no tiene por qué ser monótona.
\fuente{adaptado de Cambridge, \emph{Tripos} IA 2013, Paper~1~\cite{tripos2013}, pregunta~10E(b)(i).}
\end{ejercicio}

\begin{ejercicio}{3}{eje:1-monotonia-composicion}
Sean $f$ y $g$ funciones de $\R$ en $\R$. Demuestre: (a) si $f$ y $g$ son crecientes, $g\circ f$ es creciente; (b) si $f$ es creciente y $g$ decreciente, $g\circ f$ es decreciente; (c) si $f$ y $g$ son crecientes, $f+g$ es creciente; (d) el producto de dos funciones crecientes no es necesariamente creciente. Deduzca de (a) y (b) que $x\mapsto\sigma(-x^{3})$ es decreciente.
\end{ejercicio}

\begin{ejercicio}{3}{eje:1-par-impar}
Sea $D\subseteq\R$ simétrico y $f\colon D\to\R$. Demuestre que existen una única función par $p$ y una única función impar $q$ con $f=p+q$, y obtenga fórmulas para ellas. Identifíquelas cuando $f(x)=e^{x}$ y cuando $f(x)=\softplus(x)$.
\end{ejercicio}

\begin{ejercicio}{3}{eje:1-softplus}
Demuestre que $\softplus(x)=x+\softplus(-x)$ y que $\softplus(x)=-\ln\sigma(-x)$ para todo $x$. Deduzca que $\softplus$ es estrictamente creciente y que $\ln\sigma(x)=-\softplus(-x)$. Explique por qué esta última identidad es útil para calcular $\ln\sigma(x)$ con $x$ muy negativo.
\end{ejercicio}

% =====================================================================
\chapter{Sucesiones y límites}
\label{tema:2}
% =====================================================================

El concepto central del análisis es el de límite. Este tema lo introduce primero para sucesiones, donde la definición es más sencilla, y después para funciones, con la definición $\varepsilon$--$\delta$. Se demuestran las reglas de cálculo de límites, el teorema del sándwich y el teorema de la convergencia monótona; se obtienen los límites notables que se usarán en el resto del módulo y se estudian las asíntotas. Se siguen los capítulos~5 y~22 de Spivak~\cite{spivak}, el capítulo~2 y la §5.1 de los apuntes de \emph{Analysis~I} de Cambridge~\cite{cambA1}, y las §2.2--2.5 y~4.6 de OpenStax~\cite{osc1}.

\section{Sucesiones convergentes}

\begin{definicion}[Sucesión y límite]\label{def:limite-sucesion}
Una \emph{sucesión} de números reales es una función $\N\to\R$, $n\mapsto a_n$; se escribe $(a_n)_{n\ge1}$ o $(a_n)$. Se dice que $(a_n)$ \emph{converge} a $L\in\R$, y se escribe $a_n\to L$ o $\lim_{n\to\infty}a_n=L$, si
\[
\forall\varepsilon>0\ \ \exists N\in\N\ \ \forall n\ge N:\quad \abs{a_n-L}<\varepsilon.
\]
Una sucesión que no converge a ningún número real es \emph{divergente}.
\end{definicion}

La definición se lee como un juego: un adversario fija una tolerancia $\varepsilon>0$, tan pequeña como quiera, y hay que responder con un índice $N$ a partir del cual todos los términos están a distancia menor que $\varepsilon$ de $L$. El índice $N$ puede depender de $\varepsilon$, y casi siempre depende.

\begin{ejemplo}[$1/n\to0$]\label{ej:1n}
Sea $\varepsilon>0$. Por la propiedad arquimediana de $\R$ (módulo~0.3) existe $N\in\N$ con $N>1/\varepsilon$. Para $n\ge N$, $\abs{1/n-0}=1/n\le1/N<\varepsilon$. Luego $1/n\to0$.
\end{ejemplo}

\begin{ejemplo}[Una demostración $\varepsilon$--$N$ con acotación previa]\label{ej:eN}
Se demuestra que $a_n=\frac{3n+1}{n+2}\to3$. \emph{Trabajo previo:}
\[
\abs{a_n-3}=\Bigl|\frac{3n+1-3n-6}{n+2}\Bigr|=\frac{5}{n+2}<\frac5n.
\]
Basta, pues, $5/n<\varepsilon$, es decir, $n>5/\varepsilon$. \emph{Demostración:} dado $\varepsilon>0$, sea $N\in\N$ con $N>5/\varepsilon$. Para $n\ge N$, $\abs{a_n-3}<5/n\le5/N<\varepsilon$. \hfill$\blacksquare$
\end{ejemplo}

En las demostraciones $\varepsilon$--$N$ conviene separar el trabajo previo (acotar $\abs{a_n-L}$ por una expresión sencilla que tienda a~0) de la demostración en sí, que se redacta en el orden de los cuantificadores: «sea $\varepsilon>0$; sea $N=\dots$; para $n\ge N$, \dots».

\begin{proposicion}[Unicidad y acotación]\label{prop:unicidad}
(i) El límite de una sucesión convergente es único. (ii) Toda sucesión convergente está acotada.
\end{proposicion}

\begin{proof}
(i) Supóngase $a_n\to L$ y $a_n\to M$ con $L\ne M$, y sea $\varepsilon=\abs{L-M}/2>0$. Existen $N_1$ y $N_2$ tales que $\abs{a_n-L}<\varepsilon$ para $n\ge N_1$ y $\abs{a_n-M}<\varepsilon$ para $n\ge N_2$. Para $n=\max(N_1,N_2)$, la desigualdad triangular da $\abs{L-M}\le\abs{L-a_n}+\abs{a_n-M}<2\varepsilon=\abs{L-M}$, contradicción.
(ii) Sea $a_n\to L$. Con $\varepsilon=1$ existe $N$ tal que $\abs{a_n}\le\abs{a_n-L}+\abs{L}<1+\abs{L}$ para $n\ge N$. Entonces $\abs{a_n}\le K=\max\bigl(\abs{a_1},\dots,\abs{a_{N-1}},1+\abs{L}\bigr)$ para todo $n$.
\end{proof}

\begin{teorema}[Álgebra de límites de sucesiones]\label{teo:algebra-sucesiones}
Sean $a_n\to L$ y $b_n\to M$. Entonces
(i) $a_n+b_n\to L+M$; (ii) $ca_n\to cL$ para todo $c\in\R$; (iii) $a_nb_n\to LM$; (iv) si $M\ne0$ y $b_n\ne0$ para todo $n$, $a_n/b_n\to L/M$.
\end{teorema}

\begin{proof}
(i) Dado $\varepsilon>0$, existen $N_1$, $N_2$ con $\abs{a_n-L}<\varepsilon/2$ para $n\ge N_1$ y $\abs{b_n-M}<\varepsilon/2$ para $n\ge N_2$. Para $n\ge\max(N_1,N_2)$, $\abs{(a_n+b_n)-(L+M)}\le\abs{a_n-L}+\abs{b_n-M}<\varepsilon$.
(iii) Se escribe $a_nb_n-LM=a_n(b_n-M)+M(a_n-L)$. Por la \cref{prop:unicidad}, $\abs{a_n}\le K$ para algún $K>0$. Dado $\varepsilon>0$, existen $N_1,N_2$ con $\abs{b_n-M}<\varepsilon/(2K)$ para $n\ge N_1$ y $\abs{a_n-L}<\varepsilon/(2\abs{M}+1)$ para $n\ge N_2$. Para $n\ge\max(N_1,N_2)$, $\abs{a_nb_n-LM}\le K\abs{b_n-M}+\abs{M}\abs{a_n-L}<\varepsilon/2+\varepsilon/2$.
(ii) Es (iii) con la sucesión constante $b_n=c$. (iv) Basta ver que $1/b_n\to1/M$ y aplicar (iii); es el \cref{eje:2-cociente}.
\end{proof}

\begin{teorema}[Teorema del sándwich]\label{teo:sandwich}
Si $a_n\le b_n\le c_n$ para todo $n\ge n_0$ y $a_n\to L$, $c_n\to L$, entonces $b_n\to L$.
\end{teorema}

\begin{proof}
Dado $\varepsilon>0$, para $n$ suficientemente grande se tiene $L-\varepsilon<a_n$ y $c_n<L+\varepsilon$, luego $L-\varepsilon<a_n\le b_n\le c_n<L+\varepsilon$, es decir, $\abs{b_n-L}<\varepsilon$.
\end{proof}

\begin{definicion}[Divergencia a infinito]\label{def:div-infinito}
$a_n\to+\infty$ si para todo $K\in\R$ existe $N$ tal que $a_n>K$ para todo $n\ge N$. Análogamente, $a_n\to-\infty$.
\end{definicion}

\begin{proposicion}[Sucesión geométrica]\label{prop:geometrica}
Sea $r\in\R$. Si $\abs{r}<1$, $r^{n}\to0$; si $r>1$, $r^{n}\to+\infty$; si $r=1$, $r^{n}\to1$; si $r\le-1$, $(r^{n})$ diverge.
\end{proposicion}

\begin{proof}
Si $r>1$, sea $h=r-1>0$. Por la desigualdad de Bernoulli (módulo~0.3), $r^{n}=(1+h)^{n}\ge1+nh$, que supera cualquier $K$ para $n$ grande. Si $0<\abs{r}<1$, entonces $1/\abs{r}>1$ y, dado $\varepsilon>0$, existe $N$ con $(1/\abs{r})^{n}>1/\varepsilon$ para $n\ge N$, es decir, $\abs{r^{n}-0}<\varepsilon$. El caso $r=0$ es trivial, y $r=1$ también. Si $r\le-1$, $\abs{r^{n}}\ge1$ y los términos alternan de signo, luego dos términos consecutivos distan al menos~2 y la sucesión no puede converger (\cref{eje:2-alternada}).
\end{proof}

\section{Sucesiones monótonas y el número \texorpdfstring{$e$}{e}}

\begin{teorema}[Convergencia monótona]\label{teo:monotona}
Toda sucesión creciente y acotada superiormente converge a $\sup\{a_n:n\in\N\}$. Toda sucesión decreciente y acotada inferiormente converge a su ínfimo.
\end{teorema}

\begin{proof}
Sea $(a_n)$ creciente y acotada superiormente. Por el axioma del supremo (módulo~0.3), existe $L=\sup_n a_n$. Dado $\varepsilon>0$, $L-\varepsilon$ no es cota superior, luego existe $N$ con $a_N>L-\varepsilon$. Para $n\ge N$, el crecimiento da $L-\varepsilon<a_N\le a_n\le L$, luego $\abs{a_n-L}<\varepsilon$. El caso decreciente se obtiene aplicando lo anterior a $(-a_n)$.
\end{proof}

Este teorema es la primera consecuencia importante de la completitud de $\R$: en $\Q$ es falso, porque la sucesión creciente y acotada $1;\ 1{,}4;\ 1{,}41;\ 1{,}414;\dots$ de aproximaciones decimales de $\sqrt2$ no tiene límite racional.

\begin{teorema}[El número $e$]\label{teo:e}
La sucesión $e_n=\bigl(1+\tfrac1n\bigr)^{n}$ es creciente y está acotada superiormente por~3. Su límite se denota $e$ y vale $e=2{,}71828\dots$
\end{teorema}

\begin{proof}
Por el binomio de Newton,
\[
e_n=\sum_{k=0}^{n}\binom{n}{k}\frac{1}{n^{k}}=\sum_{k=0}^{n}\frac{1}{k!}\cdot\frac{n(n-1)\cdots(n-k+1)}{n^{k}}=\sum_{k=0}^{n}\frac{1}{k!}\prod_{j=0}^{k-1}\Bigl(1-\frac{j}{n}\Bigr).
\]
Al pasar de $n$ a $n+1$, cada factor $1-j/n$ aumenta a $1-j/(n+1)$ y se añade un término positivo más, luego $e_n<e_{n+1}$. Además, cada producto es $\le1$ y $k!\ge2^{k-1}$ para $k\ge1$, de modo que $e_n\le1+\sum_{k=1}^{n}2^{-(k-1)}<1+2=3$. Por el \cref{teo:monotona}, $(e_n)$ converge.
\end{proof}

\section{Límite de una función}

\begin{definicion}[Límite $\varepsilon$--$\delta$]\label{def:limite-funcion}
Sea $f$ definida en un intervalo abierto que contiene a $a$, salvo quizá en el propio $a$. Se dice que $f$ tiene límite $L$ en $a$, y se escribe $\lim_{x\to a}f(x)=L$, si
\[
\forall\varepsilon>0\ \ \exists\delta>0\ \ \forall x:\quad 0<\abs{x-a}<\delta\ \Rightarrow\ \abs{f(x)-L}<\varepsilon.
\]
\end{definicion}

La condición $0<\abs{x-a}$ excluye el punto $a$: el límite describe el comportamiento de $f$ \emph{cerca} de $a$, no \emph{en} $a$, y no depende de $f(a)$ ni de que esté definido.

\begin{ejemplo}[Un límite lineal]\label{ej:lineal}
$\lim_{x\to1}(2x+1)=3$. \emph{Trabajo previo:} $\abs{(2x+1)-3}=2\abs{x-1}$, que es $<\varepsilon$ si $\abs{x-1}<\varepsilon/2$. \emph{Demostración:} dado $\varepsilon>0$, sea $\delta=\varepsilon/2$. Si $0<\abs{x-1}<\delta$, entonces $\abs{(2x+1)-3}=2\abs{x-1}<2\delta=\varepsilon$. \hfill$\blacksquare$
\end{ejemplo}

\begin{ejemplo}[Un límite cuadrático: el truco del mínimo]\label{ej:cuadratico}
$\lim_{x\to2}x^{2}=4$. \emph{Trabajo previo:} $\abs{x^{2}-4}=\abs{x+2}\,\abs{x-2}$. El factor $\abs{x+2}$ no es constante, pero se acota si se restringe $x$: si $\abs{x-2}<1$, entonces $1<x<3$ y $\abs{x+2}<5$. \emph{Demostración:} dado $\varepsilon>0$, sea $\delta=\min(1,\varepsilon/5)$. Si $0<\abs{x-2}<\delta$, entonces $\abs{x-2}<1$, luego $\abs{x+2}<5$, y $\abs{x^{2}-4}<5\abs{x-2}<5\cdot\varepsilon/5=\varepsilon$. \hfill$\blacksquare$
\end{ejemplo}

\begin{teorema}[Caracterización por sucesiones]\label{teo:heine}
$\lim_{x\to a}f(x)=L$ si y solo si, para toda sucesión $(x_n)$ con $x_n\ne a$ y $x_n\to a$, se tiene $f(x_n)\to L$.
\end{teorema}

\begin{proof}
($\Rightarrow$) Sea $\varepsilon>0$ y $\delta$ el de la definición. Como $x_n\to a$, existe $N$ con $\abs{x_n-a}<\delta$ para $n\ge N$; como $x_n\ne a$, también $0<\abs{x_n-a}$, luego $\abs{f(x_n)-L}<\varepsilon$.
($\Leftarrow$) Por contrarrecíproco. Si el límite no es $L$, existe $\varepsilon>0$ tal que para todo $\delta>0$ hay un $x$ con $0<\abs{x-a}<\delta$ y $\abs{f(x)-L}\ge\varepsilon$. Tomando $\delta=1/n$ se obtiene $x_n$ con $0<\abs{x_n-a}<1/n$ y $\abs{f(x_n)-L}\ge\varepsilon$: $x_n\to a$, $x_n\ne a$ y $f(x_n)\not\to L$.
\end{proof}

El \cref{teo:heine} traslada a las funciones todo lo demostrado para sucesiones: el límite es único, y valen el álgebra de límites (\cref{teo:algebra-sucesiones}) y el teorema del sándwich para funciones. También sirve para probar que un límite \emph{no} existe: basta encontrar dos sucesiones $x_n\to a$ y $y_n\to a$ con $\lim f(x_n)\ne\lim f(y_n)$.

\begin{ejemplo}[Un límite que no existe]\label{ej:sin1x}
$\lim_{x\to0}\sin(1/x)$ no existe. Con $x_n=1/(2\pi n)$, $\sin(1/x_n)=0$; con $y_n=1/(2\pi n+\pi/2)$, $\sin(1/y_n)=1$. Ambas sucesiones tienden a~0 y las imágenes tienen límites distintos. En cambio, $\lim_{x\to0}x\sin(1/x)=0$ por el sándwich: $-\abs{x}\le x\sin(1/x)\le\abs{x}$.
\end{ejemplo}

\begin{definicion}[Límites laterales, infinitos y en el infinito]\label{def:limites-variantes}
Sea $f$ una función real y $a\in\R$.
\begin{enumerate}
  \item $\lim_{x\to a^{+}}f(x)=L$ si para todo $\varepsilon>0$ existe $\delta>0$ con $\abs{f(x)-L}<\varepsilon$ siempre que $a<x<a+\delta$. Análogamente, $\lim_{x\to a^{-}}$ con $a-\delta<x<a$.
  \item $\lim_{x\to a}f(x)=+\infty$ si para todo $K$ existe $\delta>0$ con $f(x)>K$ siempre que $0<\abs{x-a}<\delta$.
  \item $\lim_{x\to+\infty}f(x)=L$ si para todo $\varepsilon>0$ existe $M$ con $\abs{f(x)-L}<\varepsilon$ siempre que $x>M$.
\end{enumerate}
\end{definicion}

El límite existe si y solo si existen los dos laterales y coinciden (\cref{eje:2-laterales}). Por ejemplo, para $f(x)=\abs{x}/x$ los laterales en~0 son $1$ y $-1$, y el límite no existe.

\section{Límites notables}

\begin{proposicion}[$\sin x/x$]\label{prop:senx}
Para $0<\abs{x}<\pi/2$ se cumple $\cos x<\dfrac{\sin x}{x}<1$. En consecuencia, $\displaystyle\lim_{x\to0}\frac{\sin x}{x}=1$ y $\displaystyle\lim_{x\to0}\frac{1-\cos x}{x^{2}}=\frac12$.
\end{proposicion}

\begin{proof}
Sea $0<x<\pi/2$ y considérense, en la circunferencia unidad, el triángulo de vértices $O=(0,0)$, $A=(1,0)$ y $P=(\cos x,\sin x)$, el sector circular $OAP$ y el triángulo $OAT$ con $T=(1,\tan x)$. Están contenidos uno en otro, luego sus áreas cumplen $\tfrac12\sin x<\tfrac12x<\tfrac12\tan x$. Dividiendo por $\tfrac12\sin x>0$ y tomando inversos, $\cos x<\sin x/x<1$. Las tres expresiones son pares, luego la desigualdad vale también para $-\pi/2<x<0$. Como $\cos x\to1$ (continuidad del coseno, \cref{tema:3}), el sándwich da $\sin x/x\to1$. Por último,
\[
\frac{1-\cos x}{x^{2}}=\frac{1-\cos^{2}x}{x^{2}(1+\cos x)}=\Bigl(\frac{\sin x}{x}\Bigr)^{2}\frac{1}{1+\cos x}\longrightarrow 1\cdot\frac12.\qedhere
\]
\end{proof}

Para la exponencial se adopta el punto de vista del curso 18.01 del MIT~\cite{mit1801} (sesiones~16 a~19): el número $e$ es la única base $b>1$ para la que la gráfica de $b^{x}$ tiene pendiente~1 en el origen. Se admite, por tanto, que la exponencial es continua y que
\begin{equation}\label{eq:lim-exp}
\lim_{h\to0}\frac{e^{h}-1}{h}=1,
\end{equation}
y que esta base coincide con el límite del \cref{teo:e}. En Spivak~\cite[cap.~18]{spivak} se construyen $\ln$ y $\exp$ a partir de la integral y todo lo anterior se demuestra.

\begin{proposicion}[Límites notables de la exponencial y el logaritmo]\label{prop:notables-exp}
\[
\text{(i)}\ \lim_{x\to0}\frac{\ln(1+x)}{x}=1;\qquad
\text{(ii)}\ \lim_{x\to0}(1+x)^{1/x}=e;\qquad
\text{(iii)}\ \lim_{n\to\infty}\Bigl(1+\frac{t}{n}\Bigr)^{n}=e^{t}\ \ (t\in\R).
\]
\end{proposicion}

\begin{proof}
(i) Sea $u=\ln(1+x)$, que tiende a~0 cuando $x\to0$ (continuidad del logaritmo) y es distinto de~0 si $x\ne0$. Como $x=e^{u}-1$, $\dfrac{\ln(1+x)}{x}=\dfrac{u}{e^{u}-1}\to1$ por~\eqref{eq:lim-exp}. (ii) $(1+x)^{1/x}=\exp\bigl(\ln(1+x)/x\bigr)\to e^{1}$, por (i) y la continuidad de la exponencial. (iii) Para $t=0$ es trivial. Si $t\ne0$, con $x_n=t/n\to0$, $(1+t/n)^{n}=\bigl((1+x_n)^{1/x_n}\bigr)^{t}\to e^{t}$ por (ii) y el \cref{teo:heine}.
\end{proof}

\begin{table}[h]
\centering\small
\caption{Límites notables en $x\to0$. Se demuestran en este tema o en sus ejercicios.}\label{tab:notables}
\begin{tabular}{@{}lll@{}}
\toprule
$\dfrac{\sin x}{x}\to1$ & $\dfrac{1-\cos x}{x^{2}}\to\dfrac12$ & $\dfrac{\tan x}{x}\to1$ \\[2mm]
$\dfrac{e^{x}-1}{x}\to1$ & $\dfrac{\ln(1+x)}{x}\to1$ & $(1+x)^{1/x}\to e$ \\[2mm]
$\dfrac{a^{x}-1}{x}\to\ln a$ & $\dfrac{(1+x)^{p}-1}{x}\to p$ & $\dfrac{\tanh x}{x}\to1$ \\
\bottomrule
\end{tabular}
\end{table}

\section{Asíntotas}

\begin{definicion}[Asíntotas]\label{def:asintotas}
Sea $f$ una función.
\begin{enumerate}
  \item La recta $x=a$ es una \emph{asíntota vertical} si alguno de los límites laterales de $f$ en $a$ es $+\infty$ o $-\infty$.
  \item La recta $y=L$ es una \emph{asíntota horizontal} en $+\infty$ si $\lim_{x\to+\infty}f(x)=L$ (análogamente en $-\infty$).
  \item La recta $y=mx+b$ con $m\ne0$ es una \emph{asíntota oblicua} en $+\infty$ si $\lim_{x\to+\infty}\bigl(f(x)-mx-b\bigr)=0$.
\end{enumerate}
\end{definicion}

\begin{proposicion}[Cálculo de la asíntota oblicua]\label{prop:oblicua}
$y=mx+b$ es asíntota oblicua de $f$ en $+\infty$ si y solo si $m=\lim_{x\to+\infty}f(x)/x$ (finito y no nulo) y $b=\lim_{x\to+\infty}\bigl(f(x)-mx\bigr)$ (finito).
\end{proposicion}

\begin{proof}
Si $f(x)-mx-b\to0$, entonces $f(x)-mx\to b$, y dividiendo por $x$, $f(x)/x-m=(f(x)-mx)/x\to0$. Recíprocamente, si existen ambos límites, $f(x)-mx-b\to b-b=0$.
\end{proof}

\begin{ejemplo}[Asíntotas de una función racional y de la softplus]\label{ej:asintotas}
(a) $f(x)=\dfrac{x^{2}+1}{x-1}=x+1+\dfrac{2}{x-1}$ (división de polinomios). Tiene asíntota vertical $x=1$, con $f\to+\infty$ por la derecha y $f\to-\infty$ por la izquierda, y asíntota oblicua $y=x+1$ en $\pm\infty$, porque $2/(x-1)\to0$.
(b) $\softplus(x)\to0$ cuando $x\to-\infty$ (asíntota horizontal $y=0$) y $\softplus(x)-x=\softplus(-x)\to0$ cuando $x\to+\infty$ (asíntota oblicua $y=x$). Es lo que muestra la \cref{fig:activaciones}: la softplus se pega a la ReLU lejos del origen.
(c) $\sigma$ tiene asíntotas horizontales $y=0$ en $-\infty$ e $y=1$ en $+\infty$.
\end{ejemplo}

\begin{observacion}[Errores frecuentes]\label{obs:2-errores}
(a) En una demostración $\varepsilon$--$\delta$, elegir $\delta$ dependiente de $x$: $\delta$ solo puede depender de $\varepsilon$ (y de $a$).
(b) Escribir $\infty-\infty=0$, $0\cdot\infty=0$ o $1^{\infty}=1$: son \emph{formas indeterminadas}, cuyo valor depende de las funciones concretas; por ejemplo, $(1+1/n)^{n}\to e\ne1$.
(c) Concluir que $\lim_{x\to a}f(x)=f(a)$ sin justificar la continuidad; el límite no depende del valor $f(a)$.
(d) Aplicar el álgebra de límites cuando alguno no existe: $\lim(f+g)$ puede existir aunque no existan $\lim f$ ni $\lim g$.
(e) Creer que una sucesión acotada converge: $(-1)^{n}$ es acotada y divergente; hace falta además la monotonía (\cref{teo:monotona}).
\end{observacion}

\begin{observacion}[Conexión con el curso]
En el módulo~1 se demuestra que la softmax con temperatura tiende al vector indicador del máximo cuando $T\to0^{+}$ y a la distribución uniforme cuando $T\to\infty$; son límites de funciones como los de este tema, y la demostración usa que $e^{-c/T}\to0$ para $c>0$. El \cref{teo:monotona} es la herramienta con la que, en el módulo~7, se prueba que la pérdida decrece y converge a lo largo del descenso de gradiente, y la \cref{prop:notables-exp}(iii) reaparece en el módulo~21 como el factor $(1-1/n)^{n}\to e^{-1}$.
\end{observacion}

\begin{lecturas}
\item Spivak~\cite{spivak}, cap.~5 «Limits» y cap.~22 «Infinite Sequences».
\item Gowers y Chua~\cite{cambA1}, cap.~2 «Convergence of sequences» (§2.1--2.3) y §5.1 «Limits».
\item OpenStax~\cite{osc1}, §2.2 «The Limit of a Function», §2.3 «The Limit Laws», §2.5 «The Precise Definition of a Limit» y §4.6 «Limits at Infinity and Asymptotes».
\item Stewart~\cite{stewart}, §2.4 «The Precise Definition of a Limit» y §2.6 «Limits at Infinity; Horizontal Asymptotes».
\end{lecturas}

\section*{Ejercicios del tema 2}
\addcontentsline{toc}{section}{Ejercicios del tema 2}

\begin{ejercicio}{1}{eje:2-negacion}
Escriba con cuantificadores la negación de «$a_n\to L$» y la de «$\lim_{x\to a}f(x)=L$». Use la primera para explicar qué hay que exhibir para probar que $a_n\not\to L$.
\end{ejercicio}

\begin{ejercicio}{1}{eje:2-N}
Para $a_n=\dfrac{2n}{n+3}$, que tiende a~2, halle el menor $N$ tal que $\abs{a_n-2}<0{,}01$ para todo $n\ge N$. Repita con $\varepsilon=10^{-6}$.
\end{ejercicio}

\begin{ejercicio}{2}{eje:2-sucesiones-calculo}
Calcule los límites: (a) $\dfrac{3n^{2}-n+4}{5n^{2}+2}$; (b) $\sqrt{n^{2}+n}-n$; (c) $\dfrac{2^{n}+3^{n}}{3^{n+1}}$; (d) $\dfrac{\sin n}{n}$; (e) $\Bigl(1-\dfrac{2}{n}\Bigr)^{n}$; (f) $\sqrt[n]{2^{n}+3^{n}}$.
\end{ejercicio}

\begin{ejercicio}{2}{eje:2-funciones-calculo}
Calcule, si existen, los límites siguientes (indique $\pm\infty$ cuando proceda):
\begin{enumerate}[label=(\alph*)]
  \item $\lim_{x\to2}\dfrac{x^{2}-4}{x^{2}-3x+2}$;
  \item $\lim_{x\to1^{+}}\dfrac{x+1}{x^{2}-1}$ y $\lim_{x\to1^{-}}\dfrac{x+1}{x^{2}-1}$;
  \item $\lim_{x\to\infty}\dfrac{2x^{3}-x}{x^{3}+4x^{2}}$;
  \item $\lim_{x\to0}\dfrac{\sqrt{1+x}-1}{x}$;
  \item $\lim_{x\to-2}\dfrac{x^{3}+8}{x+2}$.
\end{enumerate}
\fuente{adaptado de MIT 18.01SC, Jerison y Mattuck~\cite{mit1801ex}, ejercicio 1D-1 (límites de funciones racionales).}
\end{ejercicio}

\begin{ejercicio}{2}{eje:2-notables}
Usando la \cref{tab:notables} (y demostrando de paso las entradas que no se han probado en el texto), calcule: (a) $\lim_{x\to0}\dfrac{\sin 5x}{3x}$; (b) $\lim_{x\to0}\dfrac{\tan x}{x}$; (c) $\lim_{x\to0}\dfrac{e^{2x}-1}{\sin x}$; (d) $\lim_{x\to0}\dfrac{2^{x}-1}{x}$; (e) $\lim_{x\to0}\dfrac{\tanh x}{x}$; (f) $\lim_{x\to\infty}x\sin(1/x)$.
\end{ejercicio}

\begin{ejercicio}{2}{eje:2-asintotas}
Halle todas las asíntotas de: (a) $f(x)=\dfrac{2x^{2}-3x}{x+1}$; (b) $g(x)=\dfrac{x}{\sqrt{x^{2}+1}}$; (c) $h(x)=xe^{1/x}$ (para la oblicua, use la \cref{tab:notables}); (d) $k(x)=\ln(1+e^{2x})$.
\end{ejercicio}

\begin{ejercicio}{3}{eje:2-eN}
Demuestre con la definición $\varepsilon$--$N$: (a) $\dfrac{n^{2}}{n^{2}+1}\to1$; (b) $\dfrac{1}{\sqrt n}\to0$; (c) $\dfrac{(-1)^{n}}{n^{2}}\to0$; (d) $\dfrac{n+5}{2n-1}\to\dfrac12$.
\end{ejercicio}

\begin{ejercicio}{3}{eje:2-ed}
Demuestre con la definición $\varepsilon$--$\delta$: (a) $\lim_{x\to2}(5x+8)=18$; (b) $\lim_{x\to2}(x^{2}+2x)=8$; (c) $\lim_{x\to3}\dfrac{x^{2}-9}{x-3}=6$; (d) $\lim_{x\to0}x^{4}=0$.
\fuente{adaptado de OpenStax~\cite{osc1}, §2.5, ejercicios de la sección.}
\end{ejercicio}

\begin{ejercicio}{3}{eje:2-ed-dificiles}
Demuestre con la definición $\varepsilon$--$\delta$: (a) $\lim_{x\to4}\sqrt{x}=2$ (indicación: $\sqrt x-2=\frac{x-4}{\sqrt x+2}$); (b) $\lim_{x\to1}\dfrac1x=1$; (c) $\lim_{x\to0}\dfrac{1}{x^{2}}=+\infty$, con la \cref{def:limites-variantes}(ii).
\fuente{adaptado de OpenStax~\cite{osc1}, §2.5, ejercicios de la sección (límites infinitos).}
\end{ejercicio}

\begin{ejercicio}{3}{eje:2-no-existe}
Demuestre que no existen: (a) $\lim_{x\to1}\dfrac{\abs{x-1}}{x-1}$, directamente con la definición $\varepsilon$--$\delta$ (tome $\varepsilon=1$); (b) $\lim_{x\to0}f(x)$, donde $f(x)=1$ si $x\in\Q$ y $f(x)=0$ si $x\notin\Q$, con el \cref{teo:heine}.
\fuente{adaptado de OpenStax~\cite{osc1}, §2.5, ejercicios de la sección.}
\end{ejercicio}

\begin{ejercicio}{3}{eje:2-producto}
Sean $\lim_{x\to a}f(x)=L$ y $\lim_{x\to a}g(x)=M$. Demuestre directamente con la definición $\varepsilon$--$\delta$, sin usar sucesiones, que $\lim_{x\to a}f(x)g(x)=LM$. (Indicación: siga el esquema de la demostración del \cref{teo:algebra-sucesiones}(iii) y acote primero $f$ cerca de $a$.)
\fuente{adaptado de OpenStax~\cite{osc1}, §2.5, ejercicio de la sección sobre la ley del producto.}
\end{ejercicio}

\begin{ejercicio}{3}{eje:2-cociente}
(a) Sea $b_n\to M\ne0$ con $b_n\ne0$ para todo $n$. Demuestre que existe $N$ con $\abs{b_n}>\abs{M}/2$ para $n\ge N$, y deduzca que $1/b_n\to1/M$. (b) Sea $(a_n)$ una sucesión de reales no nulos. Para cada una de las condiciones siguientes, demuestre que $(1/a_n)$ converge o dé un contraejemplo: (i) $a_n\to L$ para algún $L\in\R$; (ii) $a_n\to L$ para algún $L\ne0$.
\fuente{adaptado de Cambridge, \emph{Tripos} IA 2019, Paper~1~\cite{tripos2019}, pregunta~3E.}
\end{ejercicio}

\begin{ejercicio}{3}{eje:2-alternada}
(a) Demuestre que $((-1)^{n})$ diverge. (b) Más en general, demuestre que si $a_n\to L$ entonces $a_{n+1}-a_n\to0$, y deduzca que $(r^{n})$ diverge para $r\le-1$. (c) ¿Es cierto el recíproco de (b)? Considere $a_n=\sqrt n$.
\fuente{adaptado de Gowers y Chua~\cite{cambA1}, §2.1 y §2.2.}
\end{ejercicio}

\begin{ejercicio}{3}{eje:2-recurrencia}
Sea $a_1=1$ y $a_{n+1}=\sqrt{2+a_n}$. Demuestre por inducción que $0<a_n<2$ y que $(a_n)$ es creciente. Deduzca que converge y calcule su límite.
\end{ejercicio}

\begin{ejercicio}{3}{eje:2-raizn}
(a) Demuestre que $\sqrt[n]{n}\to1$. (Indicación: escriba $\sqrt[n]{n}=1+h_n$ con $h_n\ge0$ y use $(1+h_n)^{n}\ge\binom n2h_n^{2}$.) (b) Deduzca que $\sqrt[n]{a}\to1$ para todo $a>0$.
\end{ejercicio}

\begin{ejercicio}{3}{eje:2-laterales}
Demuestre que $\lim_{x\to a}f(x)=L$ si y solo si $\lim_{x\to a^{+}}f(x)=L$ y $\lim_{x\to a^{-}}f(x)=L$. Aplíquelo a $f(x)=\ReLU(x)/x$ en $a=0$.
\end{ejercicio}

% =====================================================================
\chapter{Continuidad}
\label{tema:3}
% =====================================================================

Una función es continua cuando pequeños cambios en la variable producen pequeños cambios en el valor. Este tema formaliza esa idea con la definición de límite, demuestra que las operaciones usuales conservan la continuidad y prueba los resultados que Spivak agrupa bajo el título «tres teoremas difíciles»: el teorema de Bolzano (con su consecuencia, el de los valores intermedios) y el teorema de Weierstrass, que reúne la acotación de una función continua en un intervalo cerrado y la existencia de máximo y mínimo. Todos dependen de la completitud de los números reales. Se siguen los capítulos~6 a~8 de Spivak~\cite{spivak}, el capítulo~3 de Apostol~\cite{apostol}, el capítulo~4 de los apuntes de \emph{Analysis~I}~\cite{cambA1} y la §2.4 de OpenStax~\cite{osc1}.

\section{Funciones continuas}

\begin{definicion}[Continuidad]\label{def:continuidad}
Sea $f\colon D\to\R$ y $a\in D$. Se dice que $f$ es \emph{continua en $a$} si
\[
\forall\varepsilon>0\ \ \exists\delta>0\ \ \forall x\in D:\quad \abs{x-a}<\delta\ \Rightarrow\ \abs{f(x)-f(a)}<\varepsilon.
\]
Si $a$ es interior a $D$, esto equivale a $\lim_{x\to a}f(x)=f(a)$. Se dice que $f$ es \emph{continua} (en $D$) si lo es en cada punto de $D$.
\end{definicion}

La definición exige tres cosas a la vez: que $f(a)$ exista, que exista el límite y que coincidan. Por el \cref{teo:heine}, $f$ es continua en $a$ si y solo si $x_n\to a$ (con $x_n\in D$) implica $f(x_n)\to f(a)$: una función continua «conmuta con los límites», $\lim f(x_n)=f(\lim x_n)$.

\begin{proposicion}[Operaciones con funciones continuas]\label{prop:operaciones-continuas}
Si $f$ y $g$ son continuas en $a$, también lo son $f+g$, $cf$, $fg$ y, si $g(a)\ne0$, $f/g$. Si $f$ es continua en $a$ y $g$ es continua en $f(a)$, entonces $g\circ f$ es continua en $a$.
\end{proposicion}

\begin{proof}
Las cuatro primeras afirmaciones son el álgebra de límites (\cref{teo:algebra-sucesiones} con el \cref{teo:heine}). Para la composición, sea $x_n\to a$; por la continuidad de $f$, $f(x_n)\to f(a)$, y por la de $g$ en $f(a)$, $g(f(x_n))\to g(f(a))$.
\end{proof}

\begin{proposicion}[Continuidad de las funciones elementales]\label{prop:elementales-continuas}
Son continuas en todo su dominio: los polinomios y las funciones racionales; $\abs{x}$; $\sin$ y $\cos$; la exponencial y el logaritmo; y, en consecuencia, $\tanh$, $\sigma$, $\ReLU$, $\softplus$ y todas las funciones que se obtienen de ellas por sumas, productos, cocientes y composiciones.
\end{proposicion}

\begin{proof}
La función identidad y las constantes son continuas (con $\delta=\varepsilon$ y $\delta$ arbitrario, respectivamente), luego los polinomios y las funciones racionales lo son por la \cref{prop:operaciones-continuas}. Para el valor absoluto, $\bigl|\abs{x}-\abs{a}\bigr|\le\abs{x-a}$ (desigualdad triangular inversa), y basta $\delta=\varepsilon$. Para el seno, por la fórmula de la diferencia de senos y por $\abs{\sin t}\le\abs{t}$ (\cref{prop:senx}),
\[
\abs{\sin x-\sin a}=2\Bigl|\cos\frac{x+a}{2}\Bigr|\,\Bigl|\sin\frac{x-a}{2}\Bigr|\le2\cdot1\cdot\frac{\abs{x-a}}{2}=\abs{x-a},
\]
y basta $\delta=\varepsilon$; el coseno es análogo. La continuidad de la exponencial se admitió en el \cref{tema:2}; la del logaritmo se sigue del \cref{teo:inversa-continua}. Finalmente, $\ReLU(x)=(x+\abs{x})/2$ y las demás son composiciones de las anteriores.
\end{proof}

\begin{definicion}[Tipos de discontinuidad]\label{def:discontinuidades}
Sea $f$ definida en un entorno de $a$ y no continua en $a$. La discontinuidad es
\begin{enumerate}
  \item \emph{evitable} si existe $\lim_{x\to a}f(x)$ y es finito (pero no coincide con $f(a)$, o $f(a)$ no está definido);
  \item \emph{de salto} si existen los límites laterales, finitos y distintos;
  \item \emph{esencial} en los demás casos (algún lateral infinito o inexistente).
\end{enumerate}
\end{definicion}

\begin{ejemplo}[Discontinuidades]\label{ej:discontinuidades}
(a) $f(x)=\dfrac{\sin x}{x}$ para $x\ne0$ tiene en~0 una discontinuidad evitable: definiendo $f(0)=1$ se obtiene una función continua. (b) La función escalón $H(x)=1$ si $x\ge0$ y $H(x)=0$ si $x<0$, que es la activación del perceptrón de Rosenblatt (módulo~9), tiene un salto en~0. (c) $1/x$ y $\sin(1/x)$ tienen discontinuidades esenciales en~0 (\cref{ej:sin1x}). (d) La función de Dirichlet ($1$ en $\Q$, $0$ fuera) es discontinua en todos los puntos (\cref{eje:2-no-existe}).
\end{ejemplo}

\begin{proposicion}[Conservación del signo]\label{prop:signo}
Si $f$ es continua en $a$ y $f(a)>0$, existe $\delta>0$ tal que $f(x)>0$ para todo $x\in D$ con $\abs{x-a}<\delta$.
\end{proposicion}

\begin{proof}
Se aplica la definición con $\varepsilon=f(a)>0$: si $\abs{x-a}<\delta$, entonces $f(x)>f(a)-\varepsilon=0$.
\end{proof}

\section{El teorema de Bolzano y los valores intermedios}

\begin{teorema}[Bolzano]\label{teo:bolzano}
Sea $f\colon[a,b]\to\R$ continua con $f(a)<0<f(b)$. Entonces existe $c\in(a,b)$ con $f(c)=0$.
\end{teorema}

\begin{proof}
Se usa el método de bisección. Sean $a_0=a$, $b_0=b$. Dado $[a_n,b_n]$ con $f(a_n)<0\le f(b_n)$, sea $m=(a_n+b_n)/2$: si $f(m)<0$, se toma $[a_{n+1},b_{n+1}]=[m,b_n]$, y si $f(m)\ge0$, se toma $[a_n,m]$. Se obtienen así $(a_n)$ creciente, $(b_n)$ decreciente, ambas en $[a,b]$, con $b_n-a_n=(b-a)/2^{n}$ y $f(a_n)<0\le f(b_n)$ para todo $n$. Por el \cref{teo:monotona}, $a_n\to c$ y $b_n\to c'$, y $c'-c=\lim(b_n-a_n)=0$, luego $c=c'\in[a,b]$. Por la continuidad, $f(c)=\lim f(a_n)\le0$ y $f(c)=\lim f(b_n)\ge0$, de modo que $f(c)=0$. Como $f(a)\ne0\ne f(b)$, $c\in(a,b)$.
\end{proof}

La demostración es constructiva: proporciona un algoritmo, la \emph{bisección}, que tras $n$ pasos localiza un cero con error menor que $(b-a)/2^{n}$.

\begin{corolario}[Teorema de los valores intermedios]\label{cor:tvi}
Sea $f\colon[a,b]\to\R$ continua y sea $y$ un valor estrictamente comprendido entre $f(a)$ y $f(b)$. Entonces existe $c\in(a,b)$ con $f(c)=y$. En consecuencia, la imagen de un intervalo por una función continua es un intervalo.
\end{corolario}

\begin{proof}
Si $f(a)<y<f(b)$, se aplica el \cref{teo:bolzano} a $g(x)=f(x)-y$; si $f(a)>y>f(b)$, a $g(x)=y-f(x)$. Para la segunda afirmación: si $u<v$ son valores de $f$ en un intervalo $I$, $u=f(s)$, $v=f(t)$, y $u<y<v$, la primera parte aplicada en el intervalo de extremos $s$ y $t$ (contenido en $I$) da $y\in f(I)$; un subconjunto de $\R$ con esta propiedad es un intervalo.
\end{proof}

\begin{ejemplo}[Raíces de polinomios de grado impar]\label{ej:grado-impar}
Todo polinomio $p(x)=x^{n}+a_{n-1}x^{n-1}+\dots+a_0$ de grado $n$ impar tiene al menos una raíz real. En efecto, para $x\ne0$, $p(x)=x^{n}\bigl(1+a_{n-1}/x+\dots+a_0/x^{n}\bigr)$, y el paréntesis tiende a~1 cuando $\abs{x}\to\infty$; como $n$ es impar, $p(x)\to+\infty$ cuando $x\to+\infty$ y $p(x)\to-\infty$ cuando $x\to-\infty$. Existen, pues, $a<b$ con $p(a)<0<p(b)$, y el \cref{teo:bolzano} da una raíz en $(a,b)$.
\end{ejemplo}

\begin{ejemplo}[Bisección]\label{ej:biseccion}
La ecuación $x^{3}+x-1=0$ tiene una raíz en $(0,1)$, porque $p(0)=-1<0<1=p(1)$. Tres pasos de bisección: $p(0{,}5)=-0{,}375<0$, luego la raíz está en $(0{,}5;\,1)$; $p(0{,}75)\approx0{,}172>0$, luego en $(0{,}5;\,0{,}75)$; $p(0{,}625)\approx-0{,}131<0$, luego en $(0{,}625;\,0{,}75)$. La raíz es $0{,}6823\dots$
\end{ejemplo}

\begin{teorema}[Continuidad de la inversa]\label{teo:inversa-continua}
Sea $f\colon I\to\R$ continua y estrictamente monótona en un intervalo $I$. Entonces $J=f(I)$ es un intervalo y $f^{-1}\colon J\to I$ es continua.
\end{teorema}

\begin{proof}
$J$ es un intervalo por el \cref{cor:tvi}. Supóngase $f$ creciente (si no, se cambia $f$ por $-f$). Sea $y_0=f(x_0)\in J$ y $\varepsilon>0$; se trata el caso en que $x_0$ es interior a $I$ (en un extremo el argumento es el mismo por un solo lado). Reduciendo $\varepsilon$ si hace falta, $x_0\pm\varepsilon\in I$. Por la monotonía estricta, $f(x_0-\varepsilon)<y_0<f(x_0+\varepsilon)$. Sea $\delta=\min\bigl(y_0-f(x_0-\varepsilon),\,f(x_0+\varepsilon)-y_0\bigr)>0$. Si $y\in J$ y $\abs{y-y_0}<\delta$, entonces $f(x_0-\varepsilon)<y<f(x_0+\varepsilon)$ y, como $f^{-1}$ es creciente (\cref{prop:inversa-monotona}), $x_0-\varepsilon<f^{-1}(y)<x_0+\varepsilon$.
\end{proof}

Así, $\ln=\exp^{-1}$, $\sqrt[n]{\phantom{x}}$, $\arcsin$, $\arctan$ y $\logit=\sigma^{-1}$ son continuas.

\section{El teorema de Weierstrass}

\begin{lema}[Bolzano--Weierstrass]\label{lem:bw}
Toda sucesión acotada de números reales tiene una subsucesión convergente.
\end{lema}

\begin{proof}
Se dice que $m$ es un \emph{pico} de $(a_n)$ si $a_m\ge a_n$ para todo $n>m$. Si hay infinitos picos $m_1<m_2<\cdots$, la subsucesión $(a_{m_k})$ es decreciente. Si hay un número finito, sea $n_1$ mayor que todos ellos; como $n_1$ no es pico, existe $n_2>n_1$ con $a_{n_2}>a_{n_1}$; como $n_2$ tampoco lo es, existe $n_3>n_2$ con $a_{n_3}>a_{n_2}$, y así sucesivamente: $(a_{n_k})$ es creciente. En ambos casos la subsucesión es monótona y acotada, y converge por el \cref{teo:monotona}.
\end{proof}

\begin{teorema}[Weierstrass]\label{teo:weierstrass}
Sea $f\colon[a,b]\to\R$ continua. Entonces $f$ está acotada y alcanza sus valores máximo y mínimo: existen $x_m,x_M\in[a,b]$ con $f(x_m)\le f(x)\le f(x_M)$ para todo $x\in[a,b]$.
\end{teorema}

\begin{proof}
\emph{Acotación.} Si $f$ no estuviera acotada superiormente, para cada $n$ habría $x_n\in[a,b]$ con $f(x_n)>n$. Por el \cref{lem:bw}, una subsucesión $x_{n_k}$ converge a algún $c$, que está en $[a,b]$ porque $a\le x_{n_k}\le b$. Por la continuidad, $f(x_{n_k})\to f(c)$; pero $f(x_{n_k})>n_k\to\infty$, contradicción.
\emph{Se alcanza el máximo.} Sea $M=\sup f([a,b])$, que existe por lo anterior. Para cada $n$, $M-1/n$ no es cota superior, luego existe $x_n$ con $M-1/n<f(x_n)\le M$. Una subsucesión $x_{n_k}$ converge a algún $x_M\in[a,b]$, y $f(x_M)=\lim f(x_{n_k})=M$ por el sándwich. El mínimo se obtiene aplicando lo anterior a $-f$.
\end{proof}

Las hipótesis son necesarias: en el intervalo abierto $(0,1)$, $f(x)=1/x$ es continua y no acotada, y $f(x)=x$ es acotada pero no alcanza ni su supremo ni su ínfimo; en $[0,1]$, la función que vale $1/x$ para $x>0$ y $0$ en~$0$ no es continua y no está acotada.

\begin{observacion}[Errores frecuentes]\label{obs:3-errores}
(a) Aplicar Bolzano sin comprobar la continuidad en todo el intervalo cerrado: $1/x$ cambia de signo en $[-1,1]$ y no se anula.
(b) Creer que Bolzano da \emph{una única} raíz: solo garantiza al menos una. La unicidad requiere otro argumento, normalmente la monotonía estricta.
(c) Aplicar Weierstrass en intervalos abiertos o no acotados, como $(0,1)$ o $[0,\infty)$.
(d) Confundir «$f$ es continua» con «la gráfica de $f$ se dibuja sin levantar el lápiz»: la idea es útil en intervalos, pero $1/x$ es continua en su dominio $\R\setminus\{0\}$ y su gráfica tiene dos ramas.
\end{observacion}

\begin{observacion}[Conexión con el curso]
El \cref{teo:weierstrass} es la razón por la que, en los módulos~4 y~7, una función de pérdida continua sobre un conjunto cerrado y acotado de parámetros tiene un mínimo. La continuidad de la ReLU (que no es derivable en~0) frente al salto de la función escalón explica por qué el perceptrón clásico no puede entrenarse por descenso de gradiente y las redes con ReLU sí (módulo~9). El argumento de bisección es también el de la búsqueda binaria, y en el módulo~9 el teorema de los valores intermedios aparece en la demostración en una dimensión del teorema de aproximación universal.
\end{observacion}

\begin{lecturas}
\item Spivak~\cite{spivak}, cap.~6 «Continuous Functions», cap.~7 «Three Hard Theorems» y cap.~8 «Least Upper Bounds».
\item Apostol~\cite{apostol}, cap.~3 «Continuous Functions».
\item Gowers y Chua~\cite{cambA1}, cap.~4 «Continuous functions», §4.1.
\item OpenStax~\cite{osc1}, §2.4 «Continuity».
\end{lecturas}

\section*{Ejercicios del tema 3}
\addcontentsline{toc}{section}{Ejercicios del tema 3}

\begin{ejercicio}{1}{eje:3-verdadero-falso}
Diga si es verdadero o falso y justifique con una demostración o un contraejemplo: (a) si $f$ es continua en $[a,b]$ y $f(a)$, $f(b)$ tienen signos opuestos, $f$ tiene exactamente un cero en $[a,b]$; (b) si $f$ es continua en $[a,b]$ y $f(a)>0$, $f(b)>0$, $f$ no tiene ceros en $[a,b]$; (c) si $\abs{f}$ es continua, $f$ es continua; (d) si $f$ es continua, $\abs f$ es continua.
\fuente{adaptado de OpenStax~\cite{osc1}, §2.4, ejercicios de verdadero o falso de la sección.}
\end{ejercicio}

\begin{ejercicio}{1}{eje:3-hipotesis-weierstrass}
Para cada caso, dé un ejemplo de función que no alcance su máximo e indique qué hipótesis del \cref{teo:weierstrass} falla: (a) $f$ continua en $(0,1]$; (b) $f$ continua en $[0,\infty)$; (c) $f$ definida en $[0,1]$, acotada y discontinua en un solo punto.
\end{ejercicio}

\begin{ejercicio}{2}{eje:3-clasificar}
Estudie la continuidad y clasifique las discontinuidades de: (a) $f(x)=\dfrac{x-2}{x^{2}-4}$; (b) $g(x)=\dfrac{1}{\sin x}$; (c) $h(x)=\floor{x}$; (d) $k(x)=e^{-1/x^{2}}$ para $x\ne0$, $k(0)=0$; (e) $m(x)=\dfrac{\abs{x}}{x}$.
\fuente{adaptado de MIT 18.01SC, Jerison y Mattuck~\cite{mit1801ex}, ejercicio 1D-3.}
\end{ejercicio}

\begin{ejercicio}{2}{eje:3-constantes}
Halle los valores de las constantes que hacen continua la función en todo $\R$: (a) $f(x)=\dfrac{x^{2}+3x+2}{x+2}$ si $x\ne-2$, $f(-2)=k$; (b) $g(x)=e^{kx}$ si $x<4$, $g(x)=x+3$ si $x\ge4$; (c) $h(x)=\sqrt{kx}$ si $0\le x\le3$, $h(x)=x+1$ si $x>3$ (en $[0,\infty)$); (d) $u(x)=ax+b$ si $x<0$, $u(x)=\dfrac{\sin 2x}{x}$ si $x>0$, $u(0)=a$.
\fuente{adaptado de OpenStax~\cite{osc1}, §2.4, ejercicios de la sección.}
\end{ejercicio}

\begin{ejercicio}{2}{eje:3-biseccion}
(a) Demuestre que la ecuación $\cos t=t^{3}$ tiene al menos una solución real. (b) Aplique cuatro pasos de bisección para localizarla en un intervalo de longitud $1/16$. (c) ¿Cuántos pasos de bisección garantizan un error menor que $10^{-6}$ partiendo de $[0,1]$?
\fuente{adaptado de OpenStax~\cite{osc1}, §2.4, ejercicio de la sección sobre $\cos t=t^{3}$.}
\end{ejercicio}

\begin{ejercicio}{2}{eje:3-raices}
Demuestre que las ecuaciones siguientes tienen al menos el número indicado de soluciones reales y localícelas entre enteros consecutivos: (a) $x^{3}-3x+1=0$, tres; (b) $e^{x}=3x$, dos; (c) $\sigma(x)=x$, una (la sigmoide tiene un punto fijo).
\end{ejercicio}

\begin{ejercicio}{3}{eje:3-ed-continua}
Demuestre con la definición $\varepsilon$--$\delta$ que: (a) $f(x)=3x-7$ es continua en todo $a\in\R$; (b) $f(x)=x^{2}$ es continua en todo $a\in\R$ (el $\delta$ dependerá de $a$); (c) $f(x)=\sqrt{x}$ es continua en todo $a>0$ y en $a=0$.
\end{ejercicio}

\begin{ejercicio}{3}{eje:3-punto-fijo}
(a) Sea $f\colon[0,1]\to[0,1]$ continua. Demuestre que existe $c\in[0,1]$ con $f(c)=c$. (b) Muestre que el resultado es falso para funciones continuas de $(0,1)$ en $(0,1)$, y para funciones de $[0,1]$ en $[0,1]$ no continuas.
\end{ejercicio}

\begin{ejercicio}{3}{eje:3-cuerda}
Un ciclista recorre un puerto de montaña: sale de la base a las 8:00 y llega a la cima a las 14:00. Al día siguiente baja por el mismo camino, de 8:00 a 14:00. Demuestre que hay un punto del camino por el que pasa a la misma hora los dos días. (Indicación: considere la diferencia de las posiciones, como funciones continuas del tiempo.)
\fuente{adaptado de OpenStax~\cite{osc1}, §2.4, ejercicio de la sección sobre $s(t)-t$.}
\end{ejercicio}

\begin{ejercicio}{3}{eje:3-sucesiones}
Sea $f\colon[a,b]\to\R$ y $y\in[a,b]$. Demuestre que si $f(x_n)\to f(y)$ para toda sucesión $(x_n)$ de $[a,b]$ que tiende a $y$, entonces $f$ es continua en $y$. (Indicación: contrarrecíproco, como en el \cref{teo:heine}.)
\fuente{adaptado de Cambridge, \emph{Tripos} IA 2013, Paper~1~\cite{tripos2013}, pregunta~10E(a).}
\end{ejercicio}

\begin{ejercicio}{3}{eje:3-inyectiva-monotona}
Sea $f\colon[a,b]\to\R$ continua e inyectiva con $f(a)<f(b)$. Demuestre: (a) $f(a)<f(x)<f(b)$ para todo $x\in(a,b)$; (b) $f$ es estrictamente creciente. (Indicación para (a): si fuera $f(x)>f(b)$, aplique el \cref{cor:tvi} en $[a,x]$ y en $[x,b]$ a un valor intermedio adecuado.)
\fuente{adaptado de Cambridge, \emph{Tripos} IA 2013, Paper~1~\cite{tripos2013}, pregunta~10E(b)(ii).}
\end{ejercicio}

\begin{ejercicio}{3}{eje:3-cauchy}
Sea $g\colon\R\to\R$ con $g(x+y)=g(x)+g(y)$ para todos $x,y$. (a) Demuestre que $g(0)=0$, $g(-x)=-g(x)$ y $g(nx)=ng(x)$ para $n\in\N$; deduzca que $g(q)=cq$ para todo $q\in\Q$, con $c=g(1)$. (b) Demuestre que si $g$ es continua, $g(x)=cx$ para todo $x\in\R$. (c) Demuestre que si $g$ es continua en un solo punto $x_0$, es continua en todo $\R$.
\fuente{adaptado de Cambridge, \emph{Tripos} IA 2019, Paper~1~\cite{tripos2019}, pregunta~9D.}
\end{ejercicio}

\begin{ejercicio}{3}{eje:3-exponencial-funcional}
Sea $h\colon\R\to\R$ continua con $h(x+y)=h(x)h(y)$ para todos $x,y$. (a) Demuestre que $h$ es idénticamente nula o positiva en todo punto. (b) En el segundo caso, aplique el \cref{eje:3-cauchy} a $g=\ln h$ y deduzca que $h(x)=a^{x}$ para algún $a>0$.
\fuente{adaptado de Cambridge, \emph{Tripos} IA 2019, Paper~1~\cite{tripos2019}, pregunta~9D.}
\end{ejercicio}

\begin{ejercicio}{3}{eje:3-periodica}
Sea $f\colon\R\to\R$ continua y periódica de periodo $T>0$. (a) Demuestre que $f$ está acotada y alcanza su máximo y su mínimo en $\R$. (b) Demuestre que existe $x$ con $f(x+T/2)=f(x)$. (Indicación: estudie $g(x)=f(x+T/2)-f(x)$ en $[0,T/2]$.)
\end{ejercicio}

% =====================================================================
\chapter{La derivada y las reglas de derivación}
\label{tema:4}
% =====================================================================

La derivada mide la tasa de cambio instantánea de una función: geométricamente, la pendiente de la recta tangente; físicamente, la velocidad. Este tema define la derivada como un límite, demuestra que la derivabilidad implica la continuidad, demuestra las reglas de derivación (suma, producto, cociente, cadena e inversa) y las aplica a las funciones elementales y a las funciones de activación. Termina con la derivación implícita y logarítmica y con las derivadas sucesivas. Se siguen los capítulos~9, 10 y~12 de Spivak~\cite{spivak}, el capítulo~4 de Apostol~\cite{apostol}, las §2.7--2.8 y el capítulo~3 de Stewart~\cite{stewart} y la unidad~1 del curso 18.01SC del MIT~\cite{mit1801}.

\section{Definición de derivada}

\begin{definicion}[Derivada]\label{def:derivada}
Sea $f$ definida en un intervalo abierto que contiene a $a$. Se dice que $f$ es \emph{derivable} en $a$ si existe y es finito el límite
\[
f'(a)=\lim_{h\to0}\frac{f(a+h)-f(a)}{h}=\lim_{x\to a}\frac{f(x)-f(a)}{x-a}.
\]
El número $f'(a)$ es la \emph{derivada} de $f$ en $a$, y la recta $y=f(a)+f'(a)(x-a)$ es la \emph{recta tangente} a la gráfica en $(a,f(a))$. La función $f'\colon x\mapsto f'(x)$, definida donde $f$ es derivable, es la \emph{función derivada}; también se escribe $\dfrac{\dd f}{\dd x}$.
\end{definicion}

El cociente $\frac{f(a+h)-f(a)}{h}$ es la pendiente de la recta secante que une $(a,f(a))$ con $(a+h,f(a+h))$; la derivada es el límite de esas pendientes cuando el segundo punto se acerca al primero. Las \emph{derivadas laterales} $f'_{+}(a)$ y $f'_{-}(a)$ se definen con los límites laterales $h\to0^{+}$ y $h\to0^{-}$; $f$ es derivable en $a$ si y solo si ambas existen y coinciden.

\begin{ejemplo}[Derivadas a partir de la definición]\label{ej:def-derivada}
(a) Para $f(x)=x^{2}$: $\dfrac{(a+h)^{2}-a^{2}}{h}=2a+h\to2a$, luego $f'(a)=2a$.
(b) Para $f(x)=1/x$ con $a\ne0$: $\dfrac{1/(a+h)-1/a}{h}=\dfrac{-1}{a(a+h)}\to-\dfrac{1}{a^{2}}$.
(c) Para $f(x)=\abs{x}$ en $a=0$: el cociente es $\abs h/h$, que vale $1$ si $h>0$ y $-1$ si $h<0$. Las derivadas laterales son $1$ y $-1$, y $\abs{x}$ no es derivable en~0.
(d) Para $\ReLU$ en~0: $f'_{+}(0)=1$ y $f'_{-}(0)=0$; no es derivable en~0. Para $x\ne0$, $\ReLU'(x)=1$ si $x>0$ y $0$ si $x<0$.
\end{ejemplo}

\begin{teorema}[Derivable implica continua]\label{teo:derivable-continua}
Si $f$ es derivable en $a$, entonces $f$ es continua en $a$.
\end{teorema}

\begin{proof}
Para $x\ne a$, $f(x)-f(a)=\dfrac{f(x)-f(a)}{x-a}\,(x-a)$. Cuando $x\to a$, el primer factor tiende a $f'(a)$ y el segundo a~0, luego $f(x)-f(a)\to f'(a)\cdot0=0$.
\end{proof}

El recíproco es falso: $\abs{x}$ y $\ReLU$ son continuas en~0 y no derivables. La siguiente reformulación de la derivada es la que mejor se generaliza a varias variables (módulo~4) y simplifica las demostraciones.

\begin{proposicion}[Derivada como aproximación lineal]\label{prop:caratheodory}
Son equivalentes: (i) $f$ es derivable en $a$ con derivada $L$; (ii) existe una función $\varphi$, continua en $a$ con $\varphi(a)=L$, tal que $f(x)=f(a)+\varphi(x)(x-a)$ para todo $x$ del dominio; (iii) $f(a+h)=f(a)+Lh+r(h)$ con $r(h)/h\to0$ cuando $h\to0$.
\end{proposicion}

\begin{proof}
(i)$\Rightarrow$(ii): se define $\varphi(x)=\frac{f(x)-f(a)}{x-a}$ si $x\ne a$ y $\varphi(a)=L$; la continuidad de $\varphi$ en $a$ es exactamente (i). (ii)$\Rightarrow$(iii): $r(h)=(\varphi(a+h)-L)h$ y $r(h)/h=\varphi(a+h)-\varphi(a)\to0$. (iii)$\Rightarrow$(i): $\frac{f(a+h)-f(a)}{h}=L+\frac{r(h)}{h}\to L$.
\end{proof}

La forma (iii) dice que, cerca de $a$, $f$ se parece a la función afín $f(a)+f'(a)(x-a)$ con un error que es pequeño \emph{incluso comparado con} $h$; esta \emph{aproximación lineal} es el primer escalón del polinomio de Taylor (\cref{tema:6}).

\section{Reglas de derivación}

\begin{teorema}[Suma, producto y cociente]\label{teo:reglas}
Si $f$ y $g$ son derivables en $a$, entonces:
\begin{enumerate}
  \item $(f+g)'(a)=f'(a)+g'(a)$ y $(cf)'(a)=cf'(a)$;
  \item (regla del producto) $(fg)'(a)=f'(a)g(a)+f(a)g'(a)$;
  \item (regla del cociente) si $g(a)\ne0$, $\Bigl(\dfrac fg\Bigr)'(a)=\dfrac{f'(a)g(a)-f(a)g'(a)}{g(a)^{2}}$.
\end{enumerate}
\end{teorema}

\begin{proof}
(i) Es el álgebra de límites aplicada a los cocientes incrementales. (ii) Sumando y restando $f(a+h)g(a)$,
\[
\frac{f(a+h)g(a+h)-f(a)g(a)}{h}=f(a+h)\,\frac{g(a+h)-g(a)}{h}+g(a)\,\frac{f(a+h)-f(a)}{h}.
\]
Cuando $h\to0$, $f(a+h)\to f(a)$ por el \cref{teo:derivable-continua}, y los cocientes tienden a $g'(a)$ y $f'(a)$. (iii) Primero, para $1/g$: como $g$ es continua en $a$ y $g(a)\ne0$, $g(a+h)\ne0$ para $h$ pequeño (\cref{prop:signo}), y
\[
\frac{1/g(a+h)-1/g(a)}{h}=-\frac{g(a+h)-g(a)}{h}\cdot\frac{1}{g(a+h)g(a)}\longrightarrow-\frac{g'(a)}{g(a)^{2}}.
\]
Después, $f/g=f\cdot(1/g)$ y la regla del producto da $\frac{f'(a)}{g(a)}-\frac{f(a)g'(a)}{g(a)^{2}}$, que es la fórmula del enunciado.
\end{proof}

\begin{proposicion}[Potencias enteras]\label{prop:potencias}
Para todo $n\in\Z$, $(x^{n})'=nx^{n-1}$ (con $x\ne0$ si $n<0$).
\end{proposicion}

\begin{proof}
Para $n=0$ y $n=1$ es inmediato. Si vale para $n\ge1$, la regla del producto da $(x^{n+1})'=(x\cdot x^{n})'=x^{n}+x\cdot nx^{n-1}=(n+1)x^{n}$; por inducción, vale para todo $n\ge0$. Para $n=-m<0$, la regla del cociente da $(1/x^{m})'=-mx^{m-1}/x^{2m}=-mx^{-m-1}$.
\end{proof}

\begin{teorema}[Regla de la cadena]\label{teo:cadena}
Si $g$ es derivable en $a$ y $f$ es derivable en $g(a)$, entonces $f\circ g$ es derivable en $a$ y
\[
(f\circ g)'(a)=f'\bigl(g(a)\bigr)\,g'(a).
\]
\end{teorema}

\begin{proof}
Por la \cref{prop:caratheodory}, existe $\psi$ continua en $b=g(a)$ con $\psi(b)=f'(b)$ y $f(y)-f(b)=\psi(y)(y-b)$ para todo $y$. Tomando $y=g(x)$,
\[
\frac{f(g(x))-f(g(a))}{x-a}=\psi\bigl(g(x)\bigr)\,\frac{g(x)-g(a)}{x-a}\qquad(x\ne a).
\]
Cuando $x\to a$, $g(x)\to g(a)$ (\cref{teo:derivable-continua}) y, por la continuidad de $\psi$ en $b$, $\psi(g(x))\to\psi(b)=f'(g(a))$; el segundo factor tiende a $g'(a)$.
\end{proof}

La demostración «ingenua», que multiplica y divide por $g(x)-g(a)$, falla si $g(x)=g(a)$ para valores de $x$ arbitrariamente próximos a $a$ (por ejemplo, $g(x)=x^{2}\sin(1/x)$ en $a=0$); la función $\psi$ evita esa división. En notación de Leibniz, si $y=f(u)$ y $u=g(x)$, la regla se escribe $\dfrac{\dd y}{\dd x}=\dfrac{\dd y}{\dd u}\dfrac{\dd u}{\dd x}$.

\begin{teorema}[Derivada de la función inversa]\label{teo:inversa}
Sea $f$ continua y estrictamente monótona en un intervalo abierto $I$, derivable en $a\in I$ con $f'(a)\ne0$. Entonces $f^{-1}$ es derivable en $b=f(a)$ y
\[
(f^{-1})'(b)=\frac{1}{f'(a)}=\frac{1}{f'\bigl(f^{-1}(b)\bigr)}.
\]
\end{teorema}

\begin{proof}
Sea $y\ne b$ en $f(I)$ y $x=f^{-1}(y)$; entonces $x\ne a$ por la inyectividad, y
\[
\frac{f^{-1}(y)-f^{-1}(b)}{y-b}=\frac{x-a}{f(x)-f(a)}=\Bigl(\frac{f(x)-f(a)}{x-a}\Bigr)^{-1}.
\]
Cuando $y\to b$, $x=f^{-1}(y)\to a$ por la continuidad de $f^{-1}$ (\cref{teo:inversa-continua}), y el cociente tiende a $1/f'(a)$.
\end{proof}

Geométricamente, la gráfica de $f^{-1}$ es la simétrica de la de $f$ respecto de $y=x$, y la simetría transforma una recta de pendiente $m$ en otra de pendiente $1/m$. Si $f'(a)=0$, la tangente a $f^{-1}$ es vertical: $f(x)=x^{3}$ tiene $f'(0)=0$ y $\sqrt[3]{y}$ no es derivable en~0.

\section{Derivadas de las funciones elementales}

\begin{proposicion}[Tabla de derivadas]\label{prop:tabla}
En los dominios indicados:
\begin{enumerate}
  \item $(e^{x})'=e^{x}$ y $(\ln x)'=1/x$ para $x>0$; $(a^{x})'=a^{x}\ln a$ para $a>0$;
  \item $(x^{r})'=rx^{r-1}$ para $x>0$ y $r\in\R$;
  \item $(\sin x)'=\cos x$, $(\cos x)'=-\sin x$, $(\tan x)'=1+\tan^{2}x=1/\cos^{2}x$;
  \item $(\arcsin x)'=1/\sqrt{1-x^{2}}$ en $(-1,1)$ y $(\arctan x)'=1/(1+x^{2})$;
  \item $(\sinh x)'=\cosh x$, $(\cosh x)'=\sinh x$ y $(\tanh x)'=1-\tanh^{2}x$.
\end{enumerate}
\end{proposicion}

\begin{proof}
(i) $\dfrac{e^{x+h}-e^{x}}{h}=e^{x}\,\dfrac{e^{h}-1}{h}\to e^{x}$ por~\eqref{eq:lim-exp}. Por el \cref{teo:inversa}, con $y=e^{x}$, $(\ln)'(y)=1/e^{x}=1/y$. Como $a^{x}=e^{x\ln a}$, la regla de la cadena da $a^{x}\ln a$.
(ii) $x^{r}=e^{r\ln x}$, luego $(x^{r})'=e^{r\ln x}\cdot r/x=rx^{r-1}$.
(iii) Por la fórmula del seno de una suma,
\[
\frac{\sin(x+h)-\sin x}{h}=\sin x\,\frac{\cos h-1}{h}+\cos x\,\frac{\sin h}{h}\longrightarrow\sin x\cdot0+\cos x\cdot1,
\]
usando la \cref{prop:senx} (de $\frac{1-\cos h}{h^{2}}\to\frac12$ se sigue $\frac{\cos h-1}{h}\to0$). El coseno es análogo, y la tangente sale de la regla del cociente.
(iv) Con $y=\sin x$, $x\in(-\pi/2,\pi/2)$: $(\arcsin)'(y)=1/\cos x=1/\sqrt{1-y^{2}}$, porque $\cos x>0$ en ese intervalo. Con $y=\tan x$: $(\arctan)'(y)=1/(1+\tan^{2}x)=1/(1+y^{2})$.
(v) Se deriva la definición: $(\sinh x)'=\frac{e^{x}+e^{-x}}{2}=\cosh x$, y análogamente $\cosh$. Por la regla del cociente y la \cref{prop:hiperbolicas}(ii), $(\tanh x)'=\frac{\cosh^{2}x-\sinh^{2}x}{\cosh^{2}x}=\frac{1}{\cosh^{2}x}=1-\tanh^{2}x$.
\end{proof}

\begin{proposicion}[Derivadas de las funciones de activación]\label{prop:derivadas-activaciones}
Para todo $x\in\R$:
\[
\sigma'(x)=\sigma(x)\bigl(1-\sigma(x)\bigr),\qquad \softplus'(x)=\sigma(x),\qquad \tanh'(x)=1-\tanh^{2}x.
\]
En particular, $\sigma'(0)=\tfrac14$, $\tanh'(0)=1$, $0<\sigma'(x)\le\tfrac14$ y $0<\tanh'(x)\le1$.
\end{proposicion}

\begin{proof}
$\sigma(x)=(1+e^{-x})^{-1}$; por la regla de la cadena, $\sigma'(x)=-(1+e^{-x})^{-2}\cdot(-e^{-x})=\dfrac{1}{1+e^{-x}}\cdot\dfrac{e^{-x}}{1+e^{-x}}$, y el segundo factor es $1-\sigma(x)$. Para la softplus, $\softplus'(x)=\dfrac{e^{x}}{1+e^{x}}=\dfrac{1}{e^{-x}+1}=\sigma(x)$. La de $\tanh$ es la \cref{prop:tabla}(v). Por último, con $s=\sigma(x)\in(0,1)$, $s(1-s)=\tfrac14-(s-\tfrac12)^{2}\in(0,\tfrac14]$; y $\tanh^{2}x\in[0,1)$.
\end{proof}

Que la derivada se exprese mediante el propio valor de la función es muy conveniente al programar la retropropagación: el valor $\sigma(x)$ calculado en la pasada hacia delante basta para obtener $\sigma'(x)$ en la pasada hacia atrás.

\section{Derivación implícita y logarítmica}

Una ecuación $F(x,y)=0$ puede definir $y$ como función de $x$ cerca de un punto sin que sea posible (o cómodo) despejarla. Si se supone que $y=y(x)$ es derivable, basta derivar la ecuación respecto de $x$ con la regla de la cadena y despejar $y'$. La justificación de que tal función $y(x)$ existe y es derivable es el teorema de la función implícita, que se estudia en varias variables.

\begin{ejemplo}[Circunferencia]\label{ej:implicita}
En $x^{2}+y^{2}=25$, derivando: $2x+2yy'=0$, luego $y'=-x/y$ si $y\ne0$. En $(3,4)$, la pendiente es $-3/4$ y la tangente $y=4-\tfrac34(x-3)$, perpendicular al radio, de pendiente $4/3$.
\end{ejemplo}

\begin{ejemplo}[Derivación logarítmica]\label{ej:logaritmica}
Para derivar $f(x)=x^{x}$ ($x>0$), que no es ni potencia ni exponencial, se toma logaritmo: $\ln f(x)=x\ln x$. Derivando, $f'(x)/f(x)=\ln x+1$, luego $f'(x)=x^{x}(1+\ln x)$. El mismo método simplifica los productos y cocientes largos: si $f=\dfrac{u_1u_2\cdots u_k}{v_1\cdots v_m}$ con factores positivos, entonces
\[
\frac{f'}{f}=\sum_{i}\frac{u_i'}{u_i}-\sum_{j}\frac{v_j'}{v_j}.
\]
\end{ejemplo}

\section{Derivadas sucesivas}

Si $f'$ es a su vez derivable, su derivada es la \emph{derivada segunda} $f''=(f')'$; en general, $f^{(n)}=(f^{(n-1)})'$. Se dice que $f$ es \emph{de clase $C^{n}$} en un intervalo si $f^{(n)}$ existe y es continua en él, y de clase $C^{\infty}$ si lo es para todo $n$. Los polinomios, $e^{x}$, $\sin$, $\cos$, $\sigma$ y $\tanh$ son de clase $C^{\infty}$; la $\ReLU$ es continua pero no de clase $C^{1}$ en~$\R$.

\begin{ejemplo}[Derivadas sucesivas]\label{ej:sucesivas}
(a) $(e^{ax})^{(n)}=a^{n}e^{ax}$. (b) $(\sin x)^{(n)}=\sin(x+n\pi/2)$, porque $\sin'x=\cos x=\sin(x+\pi/2)$. (c) $\bigl(\tfrac{1}{1-x}\bigr)^{(n)}=\tfrac{n!}{(1-x)^{n+1}}$ (por inducción). (d) $\sigma''=\bigl(\sigma(1-\sigma)\bigr)'=\sigma'(1-2\sigma)=\sigma(1-\sigma)(1-2\sigma)$, que se anula en $x=0$.
\end{ejemplo}

\begin{observacion}[Errores frecuentes]\label{obs:4-errores}
(a) $(fg)'=f'g'$ y $(f/g)'=f'/g'$ son falsas; por ejemplo, $(x\cdot x)'=2x\ne1$.
(b) Olvidar la derivada interior en la regla de la cadena: $(\sin x^{2})'=2x\cos x^{2}$, no $\cos x^{2}$.
(c) Derivar $a^{x}$ como una potencia, $xa^{x-1}$, o $x^{a}$ como una exponencial; para $x^{x}$ no sirve ninguna de las dos reglas.
(d) Concluir que una función no es derivable porque la fórmula de su derivada no está definida en un punto: $f(x)=x^{2}\sin(1/x)$, $f(0)=0$, es derivable en~0 con $f'(0)=0$ (por la definición), aunque $f'(x)$ no tenga límite en~0.
(e) En la derivación implícita, derivar $y^{2}$ como $2y$ en lugar de $2yy'$.
\end{observacion}

\begin{observacion}[Conexión con el curso]
La regla de la cadena es el fundamento de la retropropagación: una red neuronal es una composición de muchas funciones, y su derivada es el producto de las derivadas de cada capa (módulos~4, 9 y~10). Las fórmulas $\sigma'=\sigma(1-\sigma)\le\tfrac14$ y $\tanh'=1-\tanh^{2}\le1$ explican el desvanecimiento del gradiente en redes profundas con estas activaciones: un producto de muchos factores menores que~1 tiende a~0 (módulo~11). La ReLU, con derivada~1 para $x>0$, evita el problema; su falta de derivabilidad en~0 se resuelve en la práctica asignando a $\ReLU'(0)$ el valor $0$ o~$1$.
\end{observacion}

\begin{lecturas}
\item Spivak~\cite{spivak}, cap.~9 «Derivatives», cap.~10 «Differentiation» y cap.~12 «Inverse Functions».
\item Apostol~\cite{apostol}, cap.~4 «Differential Calculus».
\item Stewart~\cite{stewart}, §2.7, §2.8 y cap.~3, especialmente §3.4 «The Chain Rule», §3.5 «Implicit Differentiation», §3.6 «Derivatives of Logarithmic Functions» y §3.11 «Hyperbolic Functions».
\item MIT 18.01SC~\cite{mit1801}, unidad~1: sesiones 1, 6 y 9--20.
\end{lecturas}

\section*{Ejercicios del tema 4}
\addcontentsline{toc}{section}{Ejercicios del tema 4}

\begin{ejercicio}{1}{eje:4-grafica-derivada}
Sin calcular, describa la gráfica de la derivada de: (a) una función constante; (b) $\ReLU$; (c) $\abs{x}$; (d) $\sigma$ (¿dónde es máxima $\sigma'$?). Razone por qué la derivada de una función par es impar, con un dibujo.
\end{ejercicio}

\begin{ejercicio}{2}{eje:4-definicion}
Calcule con la definición de derivada: (a) $f'(a)$ para $f(x)=x^{3}$; (b) $f'(a)$ para $f(x)=\sqrt{x}$, $a>0$; (c) $f'(1)$ para $f(x)=\dfrac{1}{2x+1}$, y escriba la recta tangente en $x=1$.
\fuente{adaptado de MIT 18.01SC, Jerison y Mattuck~\cite{mit1801ex}, ejercicio 1C-4.}
\end{ejercicio}

\begin{ejercicio}{2}{eje:4-mecanico-1}
Derive: (a) $x^{10}+3x^{5}-\pi^{3}$; (b) $(x^{3}+x)(x^{5}+x^{2})$; (c) $\dfrac{x+2}{x^{2}-1}$; (d) $\dfrac{x^{4}+1}{x}$; (e) $(x^{2}+2)^{100}$; (f) $x^{10}(x^{2}+1)^{10}$.
\fuente{adaptado de MIT 18.01SC, Jerison y Mattuck~\cite{mit1801ex}, ejercicios 1E-1, 1E-5, 1F-1 y 1F-2.}
\end{ejercicio}

\begin{ejercicio}{2}{eje:4-mecanico-2}
Derive: (a) $e^{-x^{2}/2}$; (b) $\ln(\cos x)$; (c) $\sin^{2}(3x)$; (d) $\arctan\sqrt{x}$; (e) $\ln\bigl(x+\sqrt{x^{2}+1}\bigr)$; (f) $\sigma(wx+b)$ respecto de $x$ y respecto de $w$; (g) $\tanh(x^{2})$; (h) $\ln\sigma(x)$.
\end{ejercicio}

\begin{ejercicio}{2}{eje:4-logaritmica}
Derive mediante derivación logarítmica: (a) $x^{\sin x}$ ($x>0$); (b) $(1+1/x)^{x}$ ($x>0$); (c) $\dfrac{(x+1)^{3}\sqrt{x^{2}+4}}{(2x+3)^{5}}$ ($x>0$).
\end{ejercicio}

\begin{ejercicio}{2}{eje:4-implicita}
(a) Halle la recta tangente a la astroide $x^{2/3}+y^{2/3}=4$ en el punto $(-\sqrt{27},1)$. (b) Halle los puntos de la curva $\sin x+\sin y=\tfrac12$ en los que la tangente es horizontal.
\fuente{adaptado de MIT 18.01SC, examen~1~\cite{mit1801ex}, problema~2, y Jerison y Mattuck, ejercicio 1F-5.}
\end{ejercicio}

\begin{ejercicio}{2}{eje:4-tangentes-origen}
Halle todas las rectas tangentes a la parábola $y=1+(x-1)^{2}$ que pasan por el origen.
\fuente{adaptado de MIT 18.01SC, Jerison y Mattuck~\cite{mit1801ex}, ejercicio 1C-5.}
\end{ejercicio}

\begin{ejercicio}{2}{eje:4-pegado}
Halle las constantes $a$ y $b$ que hacen derivable en todo $\R$ la función: (a) $f(x)=x^{2}+4x+1$ si $x\ge0$ y $f(x)=ax+b$ si $x<0$; (b) la misma, con el cambio de definición en $x=1$ en lugar de en $x=0$.
\fuente{adaptado de MIT 18.01SC, Jerison y Mattuck~\cite{mit1801ex}, ejercicio 1D-6.}
\end{ejercicio}

\begin{ejercicio}{2}{eje:4-sucesivas}
(a) Calcule $f^{(5)}(x)$ para $f(x)=xe^{x}$ y conjeture y demuestre por inducción una fórmula para $f^{(n)}$. (b) Calcule $\sigma''(x)$ y $\sigma'''(0)$. (c) Calcule $\tanh''(x)$ en función de $t=\tanh x$.
\fuente{apartado (a) adaptado de MIT 18.01SC, examen~1~\cite{mit1801ex}, problema~1(c).}
\end{ejercicio}

\begin{ejercicio}{3}{eje:4-producto}
Enuncie la regla del producto y demuéstrela a partir de la definición de derivada, justificando cada paso (en particular, dónde se usa la continuidad). Use después la regla del producto y la inducción para demostrar la fórmula de Leibniz $(fg)''=f''g+2f'g'+fg''$ y, en general, $(fg)^{(n)}=\sum_{k=0}^{n}\binom nkf^{(k)}g^{(n-k)}$.
\fuente{primera parte adaptada de MIT 18.01SC, examen~1~\cite{mit1801ex}, problema~4.}
\end{ejercicio}

\begin{ejercicio}{3}{eje:4-par-impar}
Sea $f\colon\R\to\R$ derivable. Demuestre que si $f$ es par, $f'$ es impar, y que si $f$ es impar, $f'$ es par. Compruébelo con $\cosh$, $\sinh$ y $\sigma(x)-\tfrac12$.
\fuente{adaptado de MIT 18.01SC, Jerison y Mattuck~\cite{mit1801ex}, ejercicio 1F-6.}
\end{ejercicio}

\begin{ejercicio}{3}{eje:4-funcional}
Sea $f\colon\R\to\R$ tal que $f(x+y)=f(x)+f(y)+x^{2}y+xy^{2}$ para todos $x,y$, y $\lim_{x\to0}f(x)/x=1$. Halle $f(0)$, demuestre que $f$ es derivable en todo punto y calcule $f'(x)$.
\fuente{adaptado de MIT 18.01SC, examen~1~\cite{mit1801ex}, problema~6.}
\end{ejercicio}

\begin{ejercicio}{3}{eje:4-raiz-doble}
Sea $p$ un polinomio. Se dice que $a$ es una \emph{raíz doble} de $p$ si $(x-a)^{2}$ divide a $p(x)$. Demuestre que $a$ es raíz doble de $p$ si y solo si $p(a)=p'(a)=0$. Aplíquelo para hallar $c$ de modo que $x^{3}-3x+c$ tenga una raíz doble.
\fuente{adaptado de Spivak~\cite{spivak}, cap.~10, problema~24.}
\end{ejercicio}

\begin{ejercicio}{3}{eje:4-x2sen}
Sea $f(x)=x^{2}\sin(1/x)$ para $x\ne0$ y $f(0)=0$. (a) Demuestre que $f$ es derivable en todo $\R$ y calcule $f'$. (b) Demuestre que $f'$ no es continua en~0. (c) Sea $g(x)=x\sin(1/x)$, $g(0)=0$: demuestre que $g$ es continua pero no derivable en~0.
\end{ejercicio}

\begin{ejercicio}{3}{eje:4-cadena-n}
Sean $f_1,\dots,f_n$ funciones derivables de $\R$ en $\R$ y $F=f_n\circ\cdots\circ f_1$. Demuestre por inducción que, con $x_0=x$ y $x_k=f_k(x_{k-1})$,
\[
F'(x)=\prod_{k=1}^{n}f_k'(x_{k-1}).
\]
Deduzca que si $F=\sigma\circ\sigma\circ\cdots\circ\sigma$ ($n$ veces), entonces $0<F'(x)\le4^{-n}$, e interprete el resultado en términos del desvanecimiento del gradiente.
\end{ejercicio}

% =====================================================================
\chapter{Teoremas del valor medio y aplicaciones}
\label{tema:5}
% =====================================================================

La derivada es una información local: describe $f$ en un entorno infinitesimal de cada punto. Los teoremas de este tema convierten esa información local en información global sobre un intervalo: si la derivada es positiva en todo un intervalo, la función crece en él; si está acotada, la función no puede variar más que una cantidad proporcional a la longitud del intervalo. El instrumento es el teorema del valor medio, del que se derivan la regla de L'Hôpital, los criterios de extremos y la teoría de la convexidad. Se siguen el capítulo~11 de Spivak~\cite{spivak}, el capítulo~4 de Apostol~\cite{apostol}, la §5.3 de los apuntes de \emph{Analysis~I}~\cite{cambA1}, las §4.1--4.5 y~4.7 de Stewart~\cite{stewart} y la unidad~2 del curso 18.01SC del MIT~\cite{mit1801}.

\section{Extremos locales y el teorema de Rolle}

\begin{definicion}[Extremos]\label{def:extremos}
Sea $f\colon D\to\R$ y $c\in D$. Se dice que $f$ tiene en $c$ un \emph{máximo absoluto} (o global) si $f(x)\le f(c)$ para todo $x\in D$, y un \emph{máximo local} (o relativo) si existe $\delta>0$ con $f(x)\le f(c)$ para todo $x\in D$ con $\abs{x-c}<\delta$. Análogamente se definen los mínimos. Un punto $c$ interior a $D$ con $f'(c)=0$ se llama \emph{punto crítico} o estacionario.
\end{definicion}

\begin{teorema}[Fermat: condición necesaria de extremo]\label{teo:fermat}
Si $f$ tiene un extremo local en un punto $c$ interior a su dominio y es derivable en $c$, entonces $f'(c)=0$.
\end{teorema}

\begin{proof}
Supóngase que es un máximo local. Para $0<h<\delta$, $f(c+h)-f(c)\le0$, luego $\dfrac{f(c+h)-f(c)}{h}\le0$ y, en el límite $h\to0^{+}$, $f'(c)\le0$. Para $-\delta<h<0$, el cociente es $\ge0$ (numerador $\le0$, denominador $<0$) y, en el límite $h\to0^{-}$, $f'(c)\ge0$. Luego $f'(c)=0$. El caso del mínimo se obtiene aplicando lo anterior a $-f$.
\end{proof}

El recíproco es falso: $x^{3}$ tiene $f'(0)=0$ y no tiene extremo en~0. Tampoco se puede omitir la hipótesis de derivabilidad ($\abs x$ tiene un mínimo en~0) ni la de punto interior ($x$ en $[0,1]$ tiene su máximo en~1, con $f'(1)=1$).

\begin{teorema}[Rolle]\label{teo:rolle}
Sea $f$ continua en $[a,b]$ y derivable en $(a,b)$, con $f(a)=f(b)$. Entonces existe $c\in(a,b)$ con $f'(c)=0$.
\end{teorema}

\begin{proof}
Por el \cref{teo:weierstrass}, $f$ alcanza en $[a,b]$ un máximo $M$ y un mínimo $m$. Si ambos se alcanzan en los extremos, entonces $M=m=f(a)$, $f$ es constante y $f'(c)=0$ para todo $c\in(a,b)$. Si no, uno de ellos se alcanza en un punto $c\in(a,b)$, que es interior, y el \cref{teo:fermat} da $f'(c)=0$.
\end{proof}

\section{El teorema del valor medio}

\begin{teorema}[Teorema del valor medio de Lagrange]\label{teo:tvm}
Sea $f$ continua en $[a,b]$ y derivable en $(a,b)$. Entonces existe $c\in(a,b)$ con
\[
f'(c)=\frac{f(b)-f(a)}{b-a},\qquad\text{es decir,}\qquad f(b)-f(a)=f'(c)(b-a).
\]
\end{teorema}

\begin{proof}
Sea $m=\frac{f(b)-f(a)}{b-a}$, la pendiente de la cuerda, y $g(x)=f(x)-m(x-a)$, que es $f$ menos la recta de pendiente $m$. Entonces $g$ es continua en $[a,b]$, derivable en $(a,b)$, y $g(a)=f(a)$, $g(b)=f(b)-(f(b)-f(a))=f(a)$. Por el \cref{teo:rolle}, existe $c\in(a,b)$ con $0=g'(c)=f'(c)-m$.
\end{proof}

Geométricamente: hay un punto en el que la tangente es paralela a la cuerda. Cinemáticamente: si un coche recorre 121 kilómetros en una hora y cincuenta minutos, su velocidad media es de 66~km/h, y en algún instante el velocímetro marca exactamente esa velocidad.

\begin{corolario}[Consecuencias del teorema del valor medio]\label{cor:tvm}
Sea $f$ continua en un intervalo $I$ y derivable en su interior.
\begin{enumerate}
  \item Si $f'=0$ en el interior de $I$, $f$ es constante en $I$.
  \item Si $f'=g'$ en el interior de $I$, entonces $f-g$ es constante.
  \item Si $f'>0$ (respectivamente $f'\ge0$) en el interior de $I$, $f$ es estrictamente creciente (respectivamente creciente) en $I$; análogamente con $<0$ y $\le0$.
  \item Si $\abs{f'(x)}\le M$ en el interior de $I$, entonces $\abs{f(x)-f(y)}\le M\abs{x-y}$ para todos $x,y\in I$; se dice que $f$ es \emph{lipschitziana} de constante~$M$.
\end{enumerate}
\end{corolario}

\begin{proof}
Sean $x<y$ en $I$. Por el \cref{teo:tvm} en $[x,y]$, $f(y)-f(x)=f'(c)(y-x)$ para algún $c\in(x,y)$. (i) $f'(c)=0$ da $f(y)=f(x)$. (ii) Es (i) aplicado a $f-g$. (iii) $f'(c)>0$ y $y-x>0$ dan $f(y)>f(x)$. (iv) $\abs{f(y)-f(x)}=\abs{f'(c)}\abs{y-x}\le M\abs{y-x}$.
\end{proof}

\begin{ejemplo}[Desigualdades con el teorema del valor medio]\label{ej:desigualdades}
(a) $\abs{\sin x-\sin y}\le\abs{x-y}$ para todos $x,y$, porque $\abs{\sin'}=\abs{\cos}\le1$.
(b) $e^{x}\ge1+x$ para todo $x\in\R$, con igualdad solo en $x=0$. Si $x>0$, el \cref{teo:tvm} en $[0,x]$ da $e^{x}-1=e^{c}x$ con $c>0$, luego $e^{x}-1>x$. Si $x<0$, en $[x,0]$: $1-e^{x}=e^{c}(-x)$ con $c<0$, luego $1-e^{x}<-x$.
(c) $\ln(1+x)\le x$ para $x>-1$: es (b) con $x$ sustituido por $\ln(1+x)$.
(d) $\sigma$ es lipschitziana de constante $\tfrac14$ y $\tanh$ de constante~1 (\cref{prop:derivadas-activaciones}).
\end{ejemplo}

\begin{teorema}[Teorema del valor medio de Cauchy]\label{teo:cauchy}
Sean $f$ y $g$ continuas en $[a,b]$ y derivables en $(a,b)$, con $g'(x)\ne0$ para todo $x\in(a,b)$. Entonces $g(b)\ne g(a)$ y existe $c\in(a,b)$ con
\[
\frac{f'(c)}{g'(c)}=\frac{f(b)-f(a)}{g(b)-g(a)}.
\]
\end{teorema}

\begin{proof}
Si fuera $g(b)=g(a)$, Rolle daría un punto con $g'=0$. Sea $k=\frac{f(b)-f(a)}{g(b)-g(a)}$ y $h(x)=f(x)-kg(x)$. Entonces $h(a)=h(b)$ (\cref{eje:5-cauchy}) y, por Rolle, existe $c\in(a,b)$ con $0=h'(c)=f'(c)-kg'(c)$.
\end{proof}

\section{La regla de L'Hôpital}

\begin{teorema}[Regla de L'Hôpital]\label{teo:lhopital}
Sean $f$ y $g$ derivables en un intervalo abierto que contiene a $a$, salvo quizá en $a$, con $g'(x)\ne0$ para $x\ne a$, y supóngase que $\lim_{x\to a}f(x)=\lim_{x\to a}g(x)=0$. Si existe $\lim_{x\to a}\dfrac{f'(x)}{g'(x)}=\ell$ (finito o infinito), entonces
\[
\lim_{x\to a}\frac{f(x)}{g(x)}=\ell.
\]
El resultado vale también para límites laterales, para $a=\pm\infty$, y cuando $\lim f=\lim g=\pm\infty$ (forma $\infty/\infty$).
\end{teorema}

\begin{proof}
Se prueba el caso $0/0$ con $a$ finito y $x\to a^{+}$; el caso $x\to a^{-}$ es simétrico. Se redefine (o define) $f(a)=g(a)=0$, con lo que $f$ y $g$ son continuas en $[a,x]$ para $x>a$ próximo. Por el \cref{teo:cauchy} en $[a,x]$, existe $c_x\in(a,x)$ con
\[
\frac{f(x)}{g(x)}=\frac{f(x)-f(a)}{g(x)-g(a)}=\frac{f'(c_x)}{g'(c_x)}.
\]
Cuando $x\to a^{+}$, también $c_x\to a^{+}$ (pues $a<c_x<x$), luego el último cociente tiende a $\ell$. La demostración de los demás casos puede verse en~\cite[cap.~11]{spivak}.
\end{proof}

\begin{ejemplo}[Aplicaciones de L'Hôpital]\label{ej:lhopital}
(a) $\lim_{x\to0}\dfrac{x-\sin x}{x^{3}}\overset{0/0}{=}\lim\dfrac{1-\cos x}{3x^{2}}=\dfrac16$, por la \cref{prop:senx}.
(b) $\lim_{x\to+\infty}\dfrac{\ln x}{x^{p}}\overset{\infty/\infty}{=}\lim\dfrac{1/x}{px^{p-1}}=\lim\dfrac{1}{px^{p}}=0$ para $p>0$.
(c) Forma $0\cdot\infty$: $\lim_{x\to0^{+}}x\ln x=\lim\dfrac{\ln x}{1/x}\overset{\infty/\infty}{=}\lim\dfrac{1/x}{-1/x^{2}}=\lim(-x)=0$.
(d) Forma $0^{0}$: $x^{x}=e^{x\ln x}\to e^{0}=1$ cuando $x\to0^{+}$, por (c).
(e) Forma $1^{\infty}$: $\lim_{x\to0}(1+\sin x)^{1/x}=\exp\Bigl(\lim\dfrac{\ln(1+\sin x)}{x}\Bigr)$, y por L'Hôpital el exponente tiende a $\lim\frac{\cos x}{1+\sin x}=1$; el límite es $e$.
\end{ejemplo}

\begin{proposicion}[Jerarquía de crecimientos]\label{prop:crecimientos}
Para todos $p>0$ y $b>1$: $\displaystyle\lim_{x\to+\infty}\frac{\ln x}{x^{p}}=0$ y $\displaystyle\lim_{x\to+\infty}\frac{x^{p}}{b^{x}}=0$.
\end{proposicion}

\begin{proof}
La primera es el \cref{ej:lhopital}(b). Para la segunda, sea $n\ge p$ entero. Aplicando L'Hôpital $n$ veces, $\lim\frac{x^{n}}{b^{x}}=\lim\frac{n!}{b^{x}(\ln b)^{n}}=0$, y $0<\frac{x^{p}}{b^{x}}\le\frac{x^{n}}{b^{x}}$ para $x\ge1$.
\end{proof}

\section{Monotonía, extremos y optimización}

\begin{proposicion}[Criterios de extremo local]\label{prop:criterios}
Sea $f$ derivable en un entorno de un punto crítico $c$.
\begin{enumerate}
  \item (\emph{Criterio de la primera derivada}) Si $f'\ge0$ a la izquierda de $c$ y $f'\le0$ a la derecha (en un entorno), $f$ tiene un máximo local en $c$; con los signos opuestos, un mínimo local.
  \item (\emph{Criterio de la segunda derivada}) Si $f''(c)$ existe y $f''(c)<0$, $f$ tiene un máximo local estricto en $c$; si $f''(c)>0$, un mínimo local estricto. Si $f''(c)=0$, el criterio no decide.
\end{enumerate}
\end{proposicion}

\begin{proof}
(i) Por el \cref{cor:tvm}(iii), $f$ es creciente en $(c-\delta,c]$ y decreciente en $[c,c+\delta)$, luego $f(x)\le f(c)$ en $(c-\delta,c+\delta)$. (ii) Si $f''(c)<0$, entonces $\lim_{h\to0}\frac{f'(c+h)-f'(c)}{h}=\lim\frac{f'(c+h)}{h}<0$, y para $h$ pequeño $f'(c+h)/h<0$: $f'>0$ a la izquierda de $c$ y $f'<0$ a la derecha. Por la demostración de (i), con monotonía estricta, $c$ es un máximo local estricto. Los ejemplos $x^{4}$, $-x^{4}$ y $x^{3}$ en $c=0$ muestran que con $f''(c)=0$ puede haber mínimo, máximo o ninguno de los dos.
\end{proof}

\begin{proposicion}[Extremos absolutos en un intervalo cerrado]\label{prop:intervalo-cerrado}
Sea $f$ continua en $[a,b]$. Su máximo y su mínimo absolutos (que existen por el \cref{teo:weierstrass}) se alcanzan en los extremos $a$, $b$, en puntos críticos de $(a,b)$ o en puntos de $(a,b)$ donde $f$ no es derivable. Basta, pues, evaluar $f$ en esos candidatos y comparar.
\end{proposicion}

\begin{proof}
Si el máximo se alcanza en un punto $c\in(a,b)$ en el que $f$ es derivable, el \cref{teo:fermat} da $f'(c)=0$.
\end{proof}

\begin{ejemplo}[Un problema de optimización]\label{ej:optimizacion}
Se quiere construir una caja abierta recortando cuadrados de lado $x$ en las esquinas de una lámina cuadrada de lado~12 y doblando los bordes. El volumen es $V(x)=x(12-2x)^{2}$ para $x\in[0,6]$. Entonces $V'(x)=(12-2x)^{2}-4x(12-2x)=(12-2x)(12-6x)$, que se anula en $x=2$ y $x=6$. Los candidatos dan $V(0)=V(6)=0$ y $V(2)=2\cdot64=128$: el volumen máximo es $128$, con $x=2$.
\end{ejemplo}

\section{Convexidad}

\begin{definicion}[Función convexa]\label{def:convexa}
Una función $f\colon I\to\R$ en un intervalo $I$ es \emph{convexa} si para todos $x,y\in I$ y todo $t\in[0,1]$
\[
f\bigl(tx+(1-t)y\bigr)\le tf(x)+(1-t)f(y),
\]
es decir, si el segmento que une dos puntos de la gráfica queda por encima de ella. Es \emph{estrictamente convexa} si la desigualdad es estricta para $x\ne y$ y $t\in(0,1)$, y \emph{cóncava} si $-f$ es convexa. Un \emph{punto de inflexión} es un punto donde $f$ es continua y pasa de convexa a cóncava o al revés.
\end{definicion}

\begin{teorema}[Caracterizaciones de la convexidad]\label{teo:convexidad}
Sea $f$ derivable en un intervalo abierto $I$.
\begin{enumerate}
  \item Si $f'$ es creciente en $I$, entonces $f$ está por encima de sus tangentes: $f(y)\ge f(x)+f'(x)(y-x)$ para todos $x,y\in I$.
  \item Si $f$ está por encima de sus tangentes, $f$ es convexa.
  \item En particular, si $f''\ge0$ en $I$, $f$ es convexa; si $f''>0$, es estrictamente convexa.
\end{enumerate}
\end{teorema}

\begin{proof}
(i) Sea $y>x$. Por el \cref{teo:tvm}, $f(y)-f(x)=f'(c)(y-x)$ con $c\in(x,y)$, y $f'(c)\ge f'(x)$ por el crecimiento de $f'$; como $y-x>0$, $f(y)-f(x)\ge f'(x)(y-x)$. Si $y<x$, se aplica el TVM en $[y,x]$: $f(x)-f(y)=f'(c)(x-y)\le f'(x)(x-y)$, que es la misma desigualdad.
(ii) Sean $x,y\in I$, $t\in[0,1]$ y $z=tx+(1-t)y$. Las desigualdades de la tangente en $z$ dan $f(x)\ge f(z)+f'(z)(x-z)$ y $f(y)\ge f(z)+f'(z)(y-z)$. Multiplicando por $t$ y por $1-t$ y sumando, $tf(x)+(1-t)f(y)\ge f(z)+f'(z)\bigl(tx+(1-t)y-z\bigr)=f(z)$.
(iii) Si $f''\ge0$, $f'$ es creciente por el \cref{cor:tvm}(iii), y se aplican (i) y (ii). La versión estricta se obtiene igual con desigualdades estrictas.
\end{proof}

\begin{ejemplo}[Convexidad de funciones elementales y de activación]\label{ej:convexas}
$x^{2}$, $e^{x}$ y $-\ln x$ son estrictamente convexas ($f''=2$, $e^{x}$ y $1/x^{2}$). La softplus es estrictamente convexa: $\softplus''=\sigma'>0$. La sigmoide es convexa en $(-\infty,0]$ y cóncava en $[0,\infty)$, porque $\sigma''=\sigma(1-\sigma)(1-2\sigma)$ tiene el signo de $1-2\sigma$; su único punto de inflexión es $x=0$. La $\ReLU$ es convexa (es el máximo de dos funciones afines, \cref{eje:5-max-convexas}), aunque no sea derivable.
\end{ejemplo}

\begin{observacion}[Errores frecuentes]\label{obs:5-errores}
(a) Aplicar L'Hôpital sin comprobar que la forma es $0/0$ o $\infty/\infty$: $\lim_{x\to0}\frac{\cos x}{x+1}=1$, mientras que el cociente de derivadas tiende a~$0$.
(b) Derivar el cociente con la regla del cociente en vez de derivar numerador y denominador por separado.
(c) Concluir que si $\lim f'/g'$ no existe, tampoco existe $\lim f/g$: $\frac{x+\sin x}{x}\to1$ cuando $x\to\infty$, aunque $\frac{1+\cos x}{1}$ no tenga límite.
(d) Identificar «punto crítico» con «extremo», o buscar extremos absolutos solo entre los puntos críticos, olvidando los extremos del intervalo y los puntos sin derivada.
(e) Pensar que $f''(c)=0$ implica punto de inflexión: $x^{4}$ tiene $f''(0)=0$ y es convexa en todo $\R$.
\end{observacion}

\begin{observacion}[Conexión con el curso]
El \cref{cor:tvm}(iv) es la idea del módulo~7: si el gradiente está acotado (o es lipschitziano), un paso pequeño del descenso de gradiente cambia poco la pérdida, y el \emph{lema de descenso} cuantifica cuánto disminuye. La convexidad garantiza que todo mínimo local es global (módulos~4 y~7), y la desigualdad de la tangente $e^{x}\ge1+x$, o su forma $\ln t\le t-1$, es el paso clave de la desigualdad de Gibbs del módulo~6. La regla de L'Hôpital justifica el convenio $0\ln0=0$ en la entropía ($x\ln x\to0$, \cref{ej:lhopital}(c)), y la \cref{prop:crecimientos} es el lenguaje con el que se comparan los costes cuadrático y lineal de la atención o las leyes de potencia del módulo~16.
\end{observacion}

\begin{lecturas}
\item Spivak~\cite{spivak}, cap.~11 «Significance of the Derivative».
\item Apostol~\cite{apostol}, cap.~4 «Differential Calculus».
\item Gowers y Chua~\cite{cambA1}, §5.3 «Differentiation theorems».
\item Stewart~\cite{stewart}, §4.1--4.5 y §4.7; OpenStax~\cite{osc1}, §4.3--4.5, §4.7 y §4.8.
\item MIT 18.01SC~\cite{mit1801}, unidad~2: sesiones 27--30 (gráficas y optimización) y 34--35 (teorema del valor medio).
\end{lecturas}

\section*{Ejercicios del tema 5}
\addcontentsline{toc}{section}{Ejercicios del tema 5}

\begin{ejercicio}{1}{eje:5-hipotesis}
Para cada función, diga si se cumplen las hipótesis del teorema del valor medio en el intervalo dado y, si no, si aun así existe algún $c$ con $f'(c)=\frac{f(b)-f(a)}{b-a}$: (a) $\abs{x-\tfrac12}$ en $[-1,1]$; (b) $1/x^{2}$ en $[-1,1]$; (c) $e^{x}$ en $[0,1]$; (d) $\sqrt[3]{x}$ en $[-1,1]$.
\fuente{adaptado de OpenStax~\cite{osc1}, §4.4, ejercicios de la sección.}
\end{ejercicio}

\begin{ejercicio}{2}{eje:5-valores-c}
Halle todos los puntos $c$ que proporciona el teorema del valor medio para: (a) $x^{2}$ en $[0,1]$; (b) $\ln x$ en $[1,2]$; (c) $x^{3}-x$ en $[-2,2]$; (d) $x^{3}$ en $[0,2]$.
\fuente{adaptado de MIT 18.01SC, Jerison y Mattuck~\cite{mit1801ex}, ejercicio 2G-1, y OpenStax~\cite{osc1}, §4.4.}
\end{ejercicio}

\begin{ejercicio}{2}{eje:5-lhopital}
Calcule: (a) $\lim_{x\to0}\dfrac{e^{x}-x-1}{x^{2}}$; (b) $\lim_{x\to0}\dfrac{\sin x-\tan x}{x^{3}}$; (c) $\lim_{x\to1}\dfrac{x-1}{\ln x}$; (d) $\lim_{x\to\infty}x^{2}e^{-x}$; (e) $\lim_{x\to0^{+}}x^{2x}$; (f) $\lim_{x\to\infty}\Bigl(1-\dfrac1x\Bigr)^{x}$; (g) $\lim_{x\to0}\Bigl(\dfrac{1}{\sin x}-\dfrac1x\Bigr)$; (h) $\lim_{x\to0}\dfrac{3^{x}-2^{x}}{x}$.
\fuente{adaptado de OpenStax~\cite{osc1}, §4.8, ejercicios de la sección.}
\end{ejercicio}

\begin{ejercicio}{2}{eje:5-grafica}
Estudie dominio, asíntotas, intervalos de crecimiento, extremos locales, convexidad y puntos de inflexión, y esboce la gráfica de: (a) $f(x)=\dfrac{x^{2}}{x-1}$; (b) $g(x)=x-3x^{1/3}$; (c) $h(x)=e^{-x^{2}}$.
\fuente{adaptado de MIT 18.01SC, Jerison y Mattuck~\cite{mit1801ex}, ejercicios 2B-1(d) y 2B-1(h), y examen~2, problema~3.}
\end{ejercicio}

\begin{ejercicio}{2}{eje:5-optimizacion}
(a) Maximice $xy$ sujeto a $x+2y=a$, con $a>0$. (b) Halle el máximo de $x^{2}\ln(1/x)$ en $(0,\infty)$. (c) Estudie si $(x+1)e^{-x}$ alcanza un máximo y un mínimo en $[0,\infty)$. (d) Halle el rectángulo de área máxima inscrito en una semicircunferencia de radio $a$, con un lado sobre el diámetro.
\fuente{adaptado de MIT 18.01SC, Jerison y Mattuck~\cite{mit1801ex}, ejercicios 2C-3, 2C-7, 2C-14(a) y 2C-15.}
\end{ejercicio}

\begin{ejercicio}{2}{eje:5-cubica}
Sea $y=x^{3}+ax^{2}+bx+c$. Halle una condición necesaria y suficiente sobre $a$, $b$, $c$ para que $y$ tenga un máximo local y un mínimo local. Compruébela con $x^{3}-3x+1$ y con $x^{3}+x$.
\fuente{adaptado de MIT 18.01SC, Jerison y Mattuck~\cite{mit1801ex}, ejercicio 2B-3.}
\end{ejercicio}

\begin{ejercicio}{2}{eje:5-una-raiz}
Demuestre que la ecuación $x^{3}+3x^{2}+16=0$ tiene exactamente una raíz real, y hállela. Demuestre también que $e^{x}=2-x$ tiene exactamente una solución real.
\fuente{adaptado de OpenStax~\cite{osc1}, §4.4, ejercicio de la sección sobre $x^{3}+3x^{2}+16$.}
\end{ejercicio}

\begin{ejercicio}{3}{eje:5-desigualdades}
Demuestre con el teorema del valor medio: (a) $\sin x<x$ para $x>0$; (b) $\sqrt{1+x}<1+\dfrac x2$ para $x>0$; (c) $\dfrac{x}{1+x}<\ln(1+x)<x$ para $x>0$; (d) $\abs{\arctan x-\arctan y}\le\abs{x-y}$.
\fuente{adaptado de MIT 18.01SC, Jerison y Mattuck~\cite{mit1801ex}, ejercicio 2G-2, y examen~2, problema~6.}
\end{ejercicio}

\begin{ejercicio}{3}{eje:5-comparacion}
(a) Demuestre que si $f$ y $g$ son derivables en $[a,\infty)$, $f(a)\ge g(a)$ y $f'(x)>g'(x)$ para $x>a$, entonces $f(x)>g(x)$ para $x>a$. (b) Aplíquelo sucesivamente para demostrar que, para $x>0$, $e^{x}>1+x$ y $e^{x}>1+x+\dfrac{x^{2}}{2}$. (c) Demuestre por inducción que $e^{x}>\sum_{k=0}^{n}x^{k}/k!$ para $x>0$ y todo $n$.
\fuente{adaptado de MIT 18.01SC, Jerison y Mattuck~\cite{mit1801ex}, ejercicio 2G-7.}
\end{ejercicio}

\begin{ejercicio}{3}{eje:5-raices-polinomio}
(a) Demuestre que si un polinomio de grado $n$ tiene $n$ raíces reales distintas, su derivada tiene $n-1$ raíces reales distintas. (b) Sea $f$ dos veces derivable con ceros en $a<b<c$; demuestre que existe $p\in(a,c)$ con $f''(p)=0$. (c) Demuestre que $x^{5}+x-1$ tiene exactamente una raíz real.
\fuente{adaptado de MIT 18.01SC, Jerison y Mattuck~\cite{mit1801ex}, ejercicios 2G-4 y 2G-5.}
\end{ejercicio}

\begin{ejercicio}{3}{eje:5-tvm-demo}
(a) Enuncie el teorema de Rolle y demuestre a partir de él el teorema del valor medio. (b) Sea $f$ dos veces derivable en $\R$. Demuestre que
\[
f''(a)=\lim_{h\to0}\frac{f(a+h)+f(a-h)-2f(a)}{h^{2}}.
\]
(Indicación: aplique L'Hôpital una vez y use la definición de $f''(a)$.)
\fuente{adaptado de Cambridge, \emph{Tripos} IA 2013, Paper~1~\cite{tripos2013}, pregunta~11E(i), (ii) y (iv).}
\end{ejercicio}

\begin{ejercicio}{3}{eje:5-cauchy}
(a) Complete la demostración del \cref{teo:cauchy} comprobando que $h(a)=h(b)$. (b) Interprete el teorema para la curva paramétrica $t\mapsto(g(t),f(t))$. (c) Deduzca la regla de L'Hôpital del caso $0/0$ con $x\to a^{+}$ y explique dónde se usa la hipótesis $g'\ne0$.
\fuente{adaptado de Cambridge, \emph{Tripos} IA 2013, Paper~1~\cite{tripos2013}, pregunta~11E(iii).}
\end{ejercicio}

\begin{ejercicio}{3}{eje:5-darboux}
(Teorema de Darboux.) Sea $g$ derivable en un intervalo abierto que contiene a $[a,b]$, con $g'(a)<k<g'(b)$. Demuestre que existe $c\in(a,b)$ con $g'(c)=k$. (Indicación: la función $\varphi(x)=g(x)-kx$ es continua en $[a,b]$; demuestre que su mínimo no se alcanza en $a$ ni en $b$.) Deduzca que una derivada no puede tener discontinuidades de salto y que la función escalón no es la derivada de ninguna función.
\fuente{adaptado de Cambridge, \emph{Tripos} IA 2019, Paper~1~\cite{tripos2019}, pregunta~10D.}
\end{ejercicio}

\begin{ejercicio}{3}{eje:5-creciente-derivada}
(a) Demuestre que si $f$ es creciente en un intervalo abierto y derivable en un punto $a$ de él, entonces $f'(a)\ge0$. (b) Explique por qué el argumento no da $f'(a)>0$ aunque $f$ sea estrictamente creciente, y dé un contraejemplo. (c) Demuestre que $\sigma$ y $\tanh$ son estrictamente crecientes y que $\ReLU$ es creciente pero no estrictamente.
\fuente{adaptado de MIT 18.01SC, Jerison y Mattuck~\cite{mit1801ex}, ejercicio 2B-7.}
\end{ejercicio}

\begin{ejercicio}{3}{eje:5-max-convexas}
(a) Demuestre con la \cref{def:convexa} que el máximo de dos funciones convexas es convexo, y deduzca que $\ReLU$ y $\abs{x}$ son convexas. (b) Demuestre que la suma de funciones convexas es convexa y que, si $f$ es convexa y $g$ es convexa y creciente, $g\circ f$ es convexa. (c) Deduzca que $x\mapsto\softplus(x^{2}-1)$ es convexa.
\end{ejercicio}

\begin{ejercicio}{3}{eje:5-jensen}
(a) Demuestre por inducción la desigualdad de Jensen finita: si $f$ es convexa, $t_1,\dots,t_n\ge0$ con $\sum t_i=1$, entonces $f\bigl(\sum t_ix_i\bigr)\le\sum t_if(x_i)$. (b) Aplíquela a $-\ln$ para obtener la desigualdad entre las medias aritmética y geométrica: $\sqrt[n]{x_1\cdots x_n}\le\frac{x_1+\dots+x_n}{n}$ para $x_i>0$.
\end{ejercicio}

% =====================================================================
\chapter{Polinomio de Taylor}
\label{tema:6}
% =====================================================================

La recta tangente es la mejor aproximación de una función por un polinomio de grado~1 cerca de un punto. El polinomio de Taylor generaliza esa idea a cualquier grado: es el polinomio que tiene en el punto las mismas derivadas que la función hasta un orden dado. El teorema de Taylor con resto de Lagrange dice cuánto se equivoca esa aproximación, lo que permite calcular valores de funciones con un error garantizado, calcular límites y clasificar puntos críticos. Se siguen el capítulo~20 de Spivak~\cite{spivak}, el capítulo~7 de Apostol~\cite{apostol}, la §5.3 de los apuntes de \emph{Analysis~I}~\cite{cambA1}, la §6.3 del segundo volumen de OpenStax~\cite{osc2} y las sesiones~23--26 del curso 18.01SC del MIT~\cite{mit1801}.

\section{El polinomio de Taylor}

\begin{definicion}[Polinomio de Taylor]\label{def:taylor}
Sea $f$ derivable $n$ veces en $a$. El \emph{polinomio de Taylor} de orden $n$ de $f$ en $a$ es
\[
P_{n,a}f(x)=\sum_{k=0}^{n}\frac{f^{(k)}(a)}{k!}(x-a)^{k}=f(a)+f'(a)(x-a)+\frac{f''(a)}{2}(x-a)^{2}+\dots+\frac{f^{(n)}(a)}{n!}(x-a)^{n},
\]
y el \emph{resto} es $R_{n,a}f(x)=f(x)-P_{n,a}f(x)$. Si $a=0$ se habla de polinomio de \emph{Maclaurin}.
\end{definicion}

\begin{proposicion}[Caracterización]\label{prop:taylor-unico}
$P=P_{n,a}f$ es el único polinomio de grado a lo sumo $n$ que cumple $P^{(k)}(a)=f^{(k)}(a)$ para $k=0,1,\dots,n$.
\end{proposicion}

\begin{proof}
Todo polinomio de grado $\le n$ se escribe de forma única como $P(x)=\sum_{k=0}^{n}c_k(x-a)^{k}$. Derivando $j$ veces, el término $c_k(x-a)^{k}$ da $\frac{k!}{(k-j)!}c_k(x-a)^{k-j}$ si $k\ge j$ y $0$ si $k<j$; en $x=a$ solo sobrevive $k=j$, luego $P^{(j)}(a)=j!\,c_j$. Así, $P^{(j)}(a)=f^{(j)}(a)$ si y solo si $c_j=f^{(j)}(a)/j!$.
\end{proof}

\begin{ejemplo}[Polinomios de Maclaurin básicos]\label{ej:maclaurin}
(a) $f(x)=e^{x}$: todas las derivadas valen $e^{0}=1$ en~0, luego $P_{n,0}(x)=\sum_{k=0}^{n}\frac{x^{k}}{k!}$.
(b) $\sin$: las derivadas en~0 son, cíclicamente, $0,1,0,-1$, luego $P_{2m+1,0}\sin(x)=x-\frac{x^{3}}{3!}+\frac{x^{5}}{5!}-\dots+(-1)^{m}\frac{x^{2m+1}}{(2m+1)!}$. Análogamente, $P_{2m,0}\cos(x)=1-\frac{x^{2}}{2!}+\dots+(-1)^{m}\frac{x^{2m}}{(2m)!}$.
(c) $\ln(1+x)$: $f^{(k)}(x)=(-1)^{k-1}(k-1)!/(1+x)^{k}$ para $k\ge1$, luego $f^{(k)}(0)/k!=(-1)^{k-1}/k$ y $P_{n,0}(x)=x-\frac{x^{2}}{2}+\frac{x^{3}}{3}-\dots+(-1)^{n-1}\frac{x^{n}}{n}$.
(d) $(1+x)^{\alpha}$, $\alpha\in\R$: $f^{(k)}(0)=\alpha(\alpha-1)\cdots(\alpha-k+1)$, luego $P_{n,0}(x)=\sum_{k=0}^{n}\binom{\alpha}{k}x^{k}$, con $\binom{\alpha}{k}=\frac{\alpha(\alpha-1)\cdots(\alpha-k+1)}{k!}$.
\end{ejemplo}

\section{El teorema de Taylor}

\begin{teorema}[Taylor con resto de Lagrange]\label{teo:taylor}
Sea $f$ derivable $n+1$ veces en un intervalo abierto $I$ que contiene a $a$. Para cada $x\in I$ existe $\xi$ entre $a$ y $x$ tal que
\[
f(x)=\sum_{k=0}^{n}\frac{f^{(k)}(a)}{k!}(x-a)^{k}+\frac{f^{(n+1)}(\xi)}{(n+1)!}(x-a)^{n+1}.
\]
\end{teorema}

\begin{proof}
Si $x=a$ no hay nada que probar. Sea $x\ne a$ fijo y, para $t$ en el intervalo cerrado de extremos $a$ y $x$, sea
\[
g(t)=f(x)-\sum_{k=0}^{n}\frac{f^{(k)}(t)}{k!}(x-t)^{k}-K\frac{(x-t)^{n+1}}{(n+1)!},
\]
donde la constante $K$ se elige para que $g(a)=0$ (es posible porque $(x-a)^{n+1}\ne0$). También $g(x)=f(x)-f(x)=0$, pues en $t=x$ solo sobrevive el término $k=0$ del sumatorio. Al derivar respecto de $t$, cada término del sumatorio da, por la regla del producto, $\frac{f^{(k+1)}(t)}{k!}(x-t)^{k}-\frac{f^{(k)}(t)}{(k-1)!}(x-t)^{k-1}$ (sin el segundo sumando si $k=0$); la suma es telescópica y queda
\[
g'(t)=-\frac{f^{(n+1)}(t)}{n!}(x-t)^{n}+K\frac{(x-t)^{n}}{n!}.
\]
Por el teorema de Rolle aplicado a $g$ entre $a$ y $x$, existe $\xi$ estrictamente entre ellos con $g'(\xi)=0$; como $x-\xi\ne0$, se obtiene $K=f^{(n+1)}(\xi)$. Sustituyendo en $g(a)=0$ se llega a la fórmula del enunciado.
\end{proof}

Para $n=0$ el teorema es exactamente el teorema del valor medio: $f(x)=f(a)+f'(\xi)(x-a)$. El punto $\xi$ es desconocido, pero basta conocer una cota de $f^{(n+1)}$ para acotar el error.

\begin{corolario}[Cota del error y forma de Peano]\label{cor:cota}
Si además $\abs{f^{(n+1)}(t)}\le M$ para todo $t$ entre $a$ y $x$, entonces
\[
\abs{R_{n,a}f(x)}\le\frac{M}{(n+1)!}\abs{x-a}^{n+1}.
\]
En particular, si $f^{(n+1)}$ está acotada cerca de $a$, entonces $R_{n,a}f(x)/(x-a)^{n}\to0$ cuando $x\to a$, lo que se escribe $f(x)=P_{n,a}f(x)+o\bigl((x-a)^{n}\bigr)$.
\end{corolario}

\begin{proof}
La cota es inmediata a partir del \cref{teo:taylor}. Para la segunda parte, $\abs{R_{n,a}f(x)/(x-a)^{n}}\le\frac{M}{(n+1)!}\abs{x-a}\to0$.
\end{proof}

La notación $o(g(x))$ («o pequeña») designa cualquier función $h$ con $h(x)/g(x)\to0$; se cumplen $o(x^{n})+o(x^{n})=o(x^{n})$, $x^{m}o(x^{n})=o(x^{n+m})$ y $o(x^{m})=o(x^{n})$ si $m\ge n$, cuando $x\to0$.

\begin{table}[h]
\centering\small
\caption{Desarrollos de Maclaurin más usados, cuando $x\to0$.}\label{tab:maclaurin}
\begin{tabular}{@{}ll@{}}
\toprule
$e^{x}=1+x+\dfrac{x^{2}}{2}+\dfrac{x^{3}}{6}+o(x^{3})$ & $\ln(1+x)=x-\dfrac{x^{2}}{2}+\dfrac{x^{3}}{3}+o(x^{3})$ \\[2mm]
$\sin x=x-\dfrac{x^{3}}{6}+o(x^{4})$ & $\cos x=1-\dfrac{x^{2}}{2}+\dfrac{x^{4}}{24}+o(x^{5})$ \\[2mm]
$(1+x)^{\alpha}=1+\alpha x+\dfrac{\alpha(\alpha-1)}{2}x^{2}+o(x^{2})$ & $\dfrac{1}{1-x}=1+x+x^{2}+x^{3}+o(x^{3})$ \\[2mm]
$\tanh x=x-\dfrac{x^{3}}{3}+o(x^{4})$ & $\sigma(x)=\dfrac12+\dfrac{x}{4}-\dfrac{x^{3}}{48}+o(x^{4})$ \\
\bottomrule
\end{tabular}
\end{table}

En la tabla, $\sin x=x-x^{3}/6+o(x^{4})$ (y no solo $o(x^{3})$) porque el término de grado~4 es nulo: $P_{3,0}\sin=P_{4,0}\sin$. Lo mismo ocurre con $\cos$, $\tanh$ y $\sigma-\tfrac12$, que son pares o impares.

\begin{ejemplo}[Desarrollos de $\tanh$ y de $\sigma$]\label{ej:taylor-activaciones}
Con $t=\tanh x$, $t'=1-t^{2}$, $t''=-2t(1-t^{2})$ y $t'''=-2(1-t^{2})(1-3t^{2})$ (\cref{eje:4-sucesivas}). En $x=0$, $t=0$, luego $t'(0)=1$, $t''(0)=0$ y $t'''(0)=-2$: $P_{3,0}\tanh(x)=x-\frac{2}{6}x^{3}=x-\frac{x^{3}}{3}$. Como $\sigma(x)=\tfrac12+\tfrac12\tanh(x/2)$ (\cref{eje:1-sigmoide-simetria}), $\sigma(x)=\tfrac12+\tfrac12\bigl(\tfrac x2-\tfrac{x^{3}}{24}\bigr)+o(x^{4})=\tfrac12+\tfrac x4-\tfrac{x^{3}}{48}+o(x^{4})$. Cerca del origen, $\sigma$ y $\tanh$ son prácticamente lineales; la no linealidad aparece en el término cúbico.
\end{ejemplo}

\section{Aproximación con error controlado}

\begin{ejemplo}[Cálculo de $e$]\label{ej:calculo-e}
Con $f=\exp$, $a=0$, $x=1$: $e=\sum_{k=0}^{n}\frac{1}{k!}+\frac{e^{\xi}}{(n+1)!}$ con $0<\xi<1$. Como $e^{\xi}<e<3$, el error es menor que $3/(n+1)!$. Para que sea menor que $10^{-6}$ basta $(n+1)!>3\cdot10^{6}$; como $9!=362\,880$ y $10!=3\,628\,800$, basta $n=9$: $\sum_{k=0}^{9}1/k!=2{,}7182815\dots$, y $e=2{,}7182818\dots$
\end{ejemplo}

\begin{ejemplo}[Aproximación de una raíz]\label{ej:raiz}
Se aproxima $\sqrt{1{,}1}$ con $f(x)=\sqrt{1+x}$ en $a=0$ y $n=1$: $P_{1,0}f(x)=1+x/2$, y $f''(t)=-\tfrac14(1+t)^{-3/2}$. Para $x=0{,}1$, $\sqrt{1{,}1}\approx1{,}05$ con $\abs{R}\le\frac{1}{2}\cdot\frac14\cdot(0{,}1)^{2}=0{,}00125$, porque $(1+t)^{-3/2}\le1$ para $t\ge0$. Además el resto es negativo ($f''<0$), luego $1{,}05$ es una aproximación por exceso. El valor exacto es $1{,}048808\dots$, con error $0{,}00119$.
\end{ejemplo}

\section{Aplicaciones a límites y a extremos}

\begin{ejemplo}[Límites mediante desarrollos]\label{ej:limites-taylor}
(a) $\dfrac{x-\sin x}{x^{3}}=\dfrac{x-(x-x^{3}/6+o(x^{4}))}{x^{3}}=\dfrac16+o(x)\to\dfrac16$, como en el \cref{ej:lhopital}(a), pero sin derivar tres veces.
(b) $\dfrac{e^{x}-1-x}{x^{2}}=\dfrac{x^{2}/2+o(x^{2})}{x^{2}}\to\dfrac12$.
(c) $\dfrac{\ln(1+x)-x}{x^{2}}=\dfrac{-x^{2}/2+o(x^{2})}{x^{2}}\to-\dfrac12$.
\end{ejemplo}

\begin{proposicion}[Criterio de la derivada $n$-ésima]\label{prop:criterio-n}
Sea $f$ de clase $C^{n}$ en un entorno de $a$, con $f'(a)=\dots=f^{(n-1)}(a)=0$ y $f^{(n)}(a)\ne0$, $n\ge2$. Si $n$ es par, $f$ tiene en $a$ un mínimo local estricto si $f^{(n)}(a)>0$ y un máximo local estricto si $f^{(n)}(a)<0$. Si $n$ es impar, $f$ no tiene extremo local en $a$.
\end{proposicion}

\begin{proof}
Por el \cref{teo:taylor} de orden $n-1$, $f(x)-f(a)=\frac{f^{(n)}(\xi_x)}{n!}(x-a)^{n}$ con $\xi_x$ entre $a$ y $x$. Por la continuidad de $f^{(n)}$ y la \cref{prop:signo}, $f^{(n)}(\xi_x)$ tiene el signo de $f^{(n)}(a)$ si $x$ está suficientemente próximo a $a$. Si $n$ es par, $(x-a)^{n}>0$ para $x\ne a$ y $f(x)-f(a)$ tiene el signo de $f^{(n)}(a)$. Si $n$ es impar, $(x-a)^{n}$ cambia de signo en $a$ y $f(x)-f(a)$ también.
\end{proof}

\begin{observacion}[Errores frecuentes]\label{obs:6-errores}
(a) Evaluar las derivadas en $x$ en lugar de en $a$: los coeficientes del polinomio de Taylor son números, $f^{(k)}(a)/k!$.
(b) Olvidar el factorial, o escribir $(x-a)^{k}$ como $x^{k}-a^{k}$.
(c) Tratar $\xi$ como un número conocido o como si fuera igual a $a$; solo se sabe que está entre $a$ y $x$, y por eso se usa una cota de $f^{(n+1)}$ en todo ese intervalo.
(d) Usar el polinomio de Taylor lejos de $a$: $P_{n,0}$ de $\ln(1+x)$ no aproxima $\ln(1+x)$ para $x>1$, por grande que sea $n$, y en general el error crece como $\abs{x-a}^{n+1}$.
(e) Truncar los desarrollos a distinto orden en el numerador y el denominador de un límite: hay que llegar, en ambos, hasta el primer término no nulo.
\end{observacion}

\begin{observacion}[Conexión con el curso]
El desarrollo de Taylor de segundo orden, en varias variables, es la herramienta del módulo~4 para clasificar puntos críticos con el hessiano y para analizar en el módulo~7 el descenso de gradiente, que se basa en la aproximación $L(\theta+h)\approx L(\theta)+L'(\theta)h+\tfrac12L''(\theta)h^{2}$. El mismo teorema da el error de las diferencias finitas con que se comprueban los gradientes (módulo~4, \cref{eje:6-diferencias}). En el módulo~11, $\ln(1+x)\approx x$ para $x$ pequeño explica por qué las bibliotecas numéricas ofrecen una función específica para $\ln(1+x)$, y en el módulo~9 la aproximación de la GELU mediante $\tanh$ con un término cúbico es un ajuste del mismo tipo que los de este tema.
\end{observacion}

\begin{lecturas}
\item Spivak~\cite{spivak}, cap.~20 «Approximation by Polynomial Functions».
\item Apostol~\cite{apostol}, cap.~7 «Polynomial Approximations to Functions».
\item Gowers y Chua~\cite{cambA1}, §5.3 «Differentiation theorems» (teorema de Taylor con resto de Lagrange).
\item OpenStax~\cite{osc2}, §6.3 «Taylor and Maclaurin Series».
\item MIT 18.01SC~\cite{mit1801}, sesiones 23--26 (aproximaciones lineal y cuadrática).
\end{lecturas}

\section*{Ejercicios del tema 6}
\addcontentsline{toc}{section}{Ejercicios del tema 6}

\begin{ejercicio}{1}{eje:6-coeficientes}
(a) Se sabe que $P_{2,0}f(x)=3-x+5x^{2}$. Halle $f(0)$, $f'(0)$ y $f''(0)$. ¿Tiene $f$ un extremo local en~0? (b) ¿Cuál es el polinomio de Taylor de orden~3 en $a=1$ de $p(x)=x^{3}$? Escríbalo en potencias de $(x-1)$ y compruebe que coincide con $p$.
\end{ejercicio}

\begin{ejercicio}{1}{eje:6-por-que}
Explique por qué el polinomio de Taylor de orden~$n$ de un polinomio de grado $\le n$ coincide con él, por qué el de una función par en $a=0$ solo tiene potencias pares y por qué $P_{3,0}\sin=P_{4,0}\sin$.
\end{ejercicio}

\begin{ejercicio}{2}{eje:6-maclaurin}
Calcule a partir de la definición: (a) $P_{4,0}$ de $\cos x$; (b) $P_{3,0}$ de $e^{-x}$; (c) $P_{3,0}$ de $\dfrac{1}{1-x}$; (d) $P_{2,0}$ de $\sqrt{1+x}$; (e) $P_{3,0}$ de $\arctan x$.
\end{ejercicio}

\begin{ejercicio}{2}{eje:6-otros-puntos}
Calcule: (a) $P_{3,1}$ de $\ln x$; (b) $P_{2,4}$ de $\sqrt{x}$; (c) $P_{3,1}$ de $1/x$. Use (a) para aproximar $\ln 1{,}1$ y compare con el valor real.
\fuente{apartado (a) adaptado de OpenStax~\cite{osc2}, §6.3, ejemplo~6.11.}
\end{ejercicio}

\begin{ejercicio}{2}{eje:6-ln12}
(a) Escriba la aproximación cuadrática general de una función dos veces derivable cerca de $x=a$. (b) Úsela para aproximar $\ln 1{,}2$. (c) Acote el error con el resto de Lagrange y compruebe la cota con el valor real.
\fuente{adaptado de MIT 18.01SC, examen~2~\cite{mit1801ex}, problema~1.}
\end{ejercicio}

\begin{ejercicio}{2}{eje:6-cuadraticas}
Halle las aproximaciones cuadráticas cerca de $x=0$ de: (a) $\tan x$; (b) $\ln\cos x$; (c) $\sin 2x$ y $\cos 2x$; (d) $\dfrac{1}{\cos x}$; (e) $e^{x^{2}}$; (f) $\dfrac{1}{a+bx}$ con $a\ne0$.
\fuente{adaptado de MIT 18.01SC, Jerison y Mattuck~\cite{mit1801ex}, ejercicios 2A-6, 2A-12(d) y 2A-13.}
\end{ejercicio}

\begin{ejercicio}{2}{eje:6-error}
(a) Aproxime $\sin 0{,}1$ con $P_{3,0}\sin$ y dé una cota del error. (b) ¿Qué orden $n$ de $P_{n,0}$ de $e^{x}$ garantiza $e^{0{,}5}$ con error menor que $10^{-5}$? (c) Aproxime $\sqrt[3]{8{,}1}$ con un polinomio de Taylor de orden~1 de $\sqrt[3]{x}$ en $a=8$ y acote el error.
\end{ejercicio}

\begin{ejercicio}{2}{eje:6-limites}
Calcule con desarrollos de Taylor los límites cuando $x\to0$ de:
\begin{enumerate}[label=(\alph*)]
  \item $\dfrac{e^{x}-1-x-x^{2}/2}{x^{3}}$;
  \item $\dfrac{\cos x-1+x^{2}/2}{x^{4}}$;
  \item $\dfrac{\sin x-x\cos x}{x^{3}}$;
  \item $\dfrac{\tanh x-x}{x^{3}}$;
  \item $\dfrac{\sigma(x)-\tfrac12-\tfrac x4}{x^{3}}$.
\end{enumerate}
\end{ejercicio}

\begin{ejercicio}{2}{eje:6-poligono}
(a) Con el teorema del coseno, demuestre que el lado de un polígono regular de $n$ lados inscrito en la circunferencia unidad mide $\sqrt{2-2\cos(2\pi/n)}$. (b) Con la aproximación cuadrática del coseno, calcule el límite del perímetro cuando $n\to\infty$. (c) Calcule el perímetro para $n=6$ y $n=100$ y compare con $2\pi$.
\fuente{adaptado de MIT 18.01SC, Jerison y Mattuck~\cite{mit1801ex}, ejercicio 2A-16.}
\end{ejercicio}

\begin{ejercicio}{3}{eje:6-taylor-orden1}
Sea $f$ dos veces derivable en un intervalo abierto que contiene a $a$ y a $x$. Demuestre, aplicando dos veces el teorema del valor medio de Cauchy a $F(t)=f(t)-f(a)-f'(a)(t-a)$ y $G(t)=(t-a)^{2}$, que existe $\xi$ entre $a$ y $x$ con $f(x)=f(a)+f'(a)(x-a)+\tfrac12f''(\xi)(x-a)^{2}$. Deduzca que si $f''\ge0$, la gráfica está por encima de la tangente.
\end{ejercicio}

\begin{ejercicio}{3}{eje:6-binomica}
(a) Enuncie el teorema de Taylor. (b) Sea $c_n=\binom{1/2}{n}$. Demuestre que $\abs{c_n}\le1$ para todo $n\ge0$. (c) Demuestre que, para $x\in(0,1)$, $\sqrt{1+x}=\lim_{N\to\infty}\sum_{n=0}^{N}c_nx^{n}$, acotando el resto de Lagrange.
\fuente{adaptado de Cambridge, \emph{Tripos} IA 2013, Paper~1~\cite{tripos2013}, pregunta~9D(b).}
\end{ejercicio}

\begin{ejercicio}{3}{eje:6-ln2}
Demuestre, acotando el resto de Lagrange, que para $x\in[0,1]$ se cumple $\abs{\ln(1+x)-P_{n,0}(x)}\le\dfrac{x^{n+1}}{n+1}$. Deduzca que $\ln2=\lim_{n\to\infty}\Bigl(1-\dfrac12+\dfrac13-\dots+\dfrac{(-1)^{n-1}}{n}\Bigr)$ y diga cuántos términos garantizan un error menor que $10^{-3}$.
\end{ejercicio}

\begin{ejercicio}{3}{eje:6-e-irracional}
Demuestre que $e$ es irracional. (Indicación: si $e=p/q$ con $p,q\in\N$, multiplique la igualdad del \cref{ej:calculo-e} con $n=q$ por $q!$ y observe que el resto multiplicado por $q!$ tendría que ser un entero estrictamente entre $0$ y $1$.)
\end{ejercicio}

\begin{ejercicio}{3}{eje:6-criterio-n}
(a) Aplique la \cref{prop:criterio-n} para clasificar el punto crítico $x=0$ de: $x^{4}$, $x^{5}$, $\cos x-1+\dfrac{x^{2}}{2}$ y $x\sin x$. (b) Demuestre que $\sin x\ge x-\dfrac{x^{3}}{6}$ para todo $x\ge0$ y que $\cos x\ge1-\dfrac{x^{2}}{2}$ para todo $x\in\R$, usando el signo del resto de Lagrange.
\end{ejercicio}

\begin{ejercicio}{3}{eje:6-diferencias}
Sea $f$ de clase $C^{3}$ con $\abs{f''}\le M_2$ y $\abs{f'''}\le M_3$. Demuestre que
\[
\Bigl|\frac{f(a+h)-f(a)}{h}-f'(a)\Bigr|\le\frac{M_2}{2}\abs h,\qquad
\Bigl|\frac{f(a+h)-f(a-h)}{2h}-f'(a)\Bigr|\le\frac{M_3}{6}h^{2}.
\]
Compruebe numéricamente las dos cotas con $f=\sin$, $a=1$ y $h=0{,}01$, e indique cuál de las dos fórmulas es preferible para comprobar gradientes.
\end{ejercicio}


\appendix
% =====================================================================
\chapter{Soluciones de los ejercicios}
\label{ap:soluciones}
% =====================================================================

Cada solución se presenta en cinco partes: datos, planteamiento, desarrollo, resultado y comprobación. En los ejercicios de demostración ($\star\star\star$), los pasos del desarrollo, leídos en orden, forman la demostración redactada tal como debería entregarse. La comprobación indica una forma independiente de verificar el resultado (un valor numérico, un caso particular, una gráfica) o señala dónde se ha usado cada hipótesis. Los valores numéricos se han comprobado con un ordenador y se redondean al número de cifras indicado.

% ---------------------------------------------------------------------
\section{Tema 1}

\begin{solucion}{eje:1-dominios}
\datos cuatro fórmulas.
\planteamiento dominio natural (\cref{def:funcion}): raíz cuadrada de argumento $\ge0$, logaritmo de argumento $>0$, denominador no nulo.
\pasos
\begin{enumerate}[label=(\alph*)]
  \item $x^{2}-5x+6=(x-2)(x-3)\ge0$ si y solo si $x\le2$ o $x\ge3$: $\dom f=(-\infty,2]\cup[3,\infty)$.
  \item Se necesita $\frac{1+x}{1-x}>0$ (y $x\ne1$): numerador y denominador del mismo signo. Ambos positivos: $-1<x<1$; ambos negativos: $x<-1$ y $x>1$, imposible. $\dom g=(-1,1)$.
  \item Se necesita $x>0$ y $\ln x\ne0$, es decir, $x\ne1$: $\dom h=(0,1)\cup(1,\infty)$.
  \item Se necesita $x>0$ y $\ln x\ge0$, es decir, $x\ge1$: $\dom k=[1,\infty)$.
\end{enumerate}
\resultado (a) $(-\infty,2]\cup[3,\infty)$; (b) $(-1,1)$; (c) $(0,1)\cup(1,\infty)$; (d) $[1,\infty)$.
\comprobacion valores de prueba: $f(2{,}5)=\sqrt{-0{,}25}$ no existe; $g(0)=\ln1=0$ existe; $h(1)=1/0$ no existe; $k(0{,}5)=\sqrt{-0{,}69}$ no existe.
\end{solucion}

\begin{solucion}{eje:1-transformaciones}
\datos $(0,0)$, $(1,2)$, $(3,-1)$ están en la gráfica de $f$.
\planteamiento si $(x_0,y_0)$ está en la gráfica de $f$, se busca el punto correspondiente en cada transformación (\cref{prop:transformaciones}).
\pasos
\begin{enumerate}[label=(\alph*)]
  \item $y=f(x-2)+1$: traslación de 2 a la derecha y 1 hacia arriba; $(x_0,y_0)\mapsto(x_0+2,y_0+1)$: $(2,1)$, $(3,3)$, $(5,0)$.
  \item $y=3f(x)$: dilatación vertical de factor~3; $(x_0,3y_0)$: $(0,0)$, $(1,6)$, $(3,-3)$.
  \item $y=f(2x)$: compresión horizontal de factor~2; $(x_0/2,y_0)$: $(0,0)$, $(\tfrac12,2)$, $(\tfrac32,-1)$.
  \item $y=-f(-x)$: simetría respecto del origen; $(-x_0,-y_0)$: $(0,0)$, $(-1,-2)$, $(-3,1)$.
\end{enumerate}
\resultado los puntos indicados en cada apartado.
\comprobacion en (c), $x=\tfrac12$ da $y=f(2\cdot\tfrac12)=f(1)=2$; en (d), $x=-1$ da $y=-f(1)=-2$.
\end{solucion}

\begin{solucion}{eje:1-paridad}
\datos seis funciones.
\planteamiento calcular $f(-x)$ y compararlo con $f(x)$ y $-f(x)$ (\cref{def:propiedades}); para descartar ambas propiedades basta un valor.
\pasos
\begin{enumerate}[label=(\alph*)]
  \item $(-x)^{3}-(-x)=-(x^{3}-x)$: \textbf{impar}.
  \item $\abs{-x}+\cos(-x)=\abs{x}+\cos x$: \textbf{par}.
  \item $e^{-1}\ne e^{1}$ y $e^{-1}\ne-e^{1}$: \textbf{ninguna}.
  \item $\ln\frac{1+x}{1-x}=-\ln\frac{1-x}{1+x}$: \textbf{impar}.
  \item $(-x)\sigma(-x)=-x(1-\sigma(x))=x\sigma(x)-x$, que no es $\pm x\sigma(x)$; con $x=1$: $f(1)=\sigma(1)\approx0{,}731$ y $f(-1)=-\sigma(-1)\approx-0{,}269$. \textbf{Ninguna}.
  \item Por la \cref{prop:relu}(ii), $\softplus(-x)=\softplus(x)-x$, luego $\softplus(-x)-\frac{-x}{2}=\softplus(x)-x+\frac x2=\softplus(x)-\frac x2$: \textbf{par}.
\end{enumerate}
\resultado impares: (a), (d); pares: (b), (f); ninguna: (c), (e).
\comprobacion (f) en $x=\pm2$: $\softplus(2)-1=2{,}1269-1=1{,}1269$ y $\softplus(-2)+1=0{,}1269+1=1{,}1269$.
\end{solucion}

\begin{solucion}{eje:1-composicion}
\datos $f(x)=\sqrt x$ ($\dom f=[0,\infty)$), $g(x)=1-x^{2}$ ($\dom g=\R$).
\planteamiento $\dom(u\circ v)=\{x\in\dom v: v(x)\in\dom u\}$.
\pasos
\begin{enumerate}[label=\arabic*.]
  \item $(f\circ g)(x)=\sqrt{1-x^{2}}$; se necesita $1-x^{2}\ge0$: dominio $[-1,1]$.
  \item $(g\circ f)(x)=1-(\sqrt x)^{2}=1-x$, pero solo para $x\in\dom f=[0,\infty)$.
  \item $(f\circ f)(x)=\sqrt{\sqrt x}=x^{1/4}$, dominio $[0,\infty)$.
  \item $(g\circ g)(x)=1-(1-x^{2})^{2}=2x^{2}-x^{4}$, dominio $\R$.
\end{enumerate}
\resultado $f\circ g=\sqrt{1-x^{2}}$ en $[-1,1]$; $g\circ f=1-x$ en $[0,\infty)$; $f\circ f=x^{1/4}$ en $[0,\infty)$; $g\circ g=2x^{2}-x^{4}$ en $\R$.
\comprobacion $g\circ f$ no es la función $1-x$ en $\R$: en $x=-1$ la composición no está definida, aunque la fórmula simplificada sí lo esté. Valor de prueba: $(g\circ g)(1)=g(0)=1=2-1$.
\end{solucion}

\begin{solucion}{eje:1-ecuaciones}
\datos cuatro ecuaciones.
\planteamiento aislar la exponencial y tomar logaritmos; en las logarítmicas, usar $\ln a+\ln b=\ln(ab)$ y comprobar el dominio.
\pasos
\begin{enumerate}[label=(\alph*)]
  \item $e^{3x}=15$, $3x=\ln15$, $x=\frac{\ln15}{3}\approx0{,}9027$.
  \item $(3x-2)\ln7=\ln11$, $x=\frac13\bigl(2+\frac{\ln11}{\ln7}\bigr)\approx1{,}0774$.
  \item Se necesita $x>2$. $\ln(x(x-2))=\ln4$, y como $\ln$ es inyectiva, $x^{2}-2x-4=0$, $x=1\pm\sqrt5$. Solo $1+\sqrt5\approx3{,}236$ cumple $x>2$.
  \item $4^{x+1}=32$, $2^{2x+2}=2^{5}$, $2x+2=5$, $x=\tfrac32$.
\end{enumerate}
\resultado (a) $\frac{\ln15}{3}$; (b) $\frac{2}{3}+\frac{\ln11}{3\ln7}$; (c) $1+\sqrt5$; (d) $\tfrac32$.
\comprobacion (c): $(1+\sqrt5)(\sqrt5-1)=5-1=4$. (d): $4^{2{,}5}=2^{5}=32$.
\end{solucion}

\begin{solucion}{eje:1-hiperbolicas-calculo}
\datos definiciones de $\sinh$ y $\cosh$.
\planteamiento sustituir la \cref{def:hiperbolicas} y usar la \cref{prop:hiperbolicas}(iii).
\pasos
\begin{enumerate}[label=(\alph*)]
  \item $2\cosh(\ln x)=e^{\ln x}+e^{-\ln x}=x+\frac1x$.
  \item $\cosh4x+\sinh4x=e^{4x}$.
  \item $\cosh2x-\sinh2x=e^{-2x}$.
  \item $\ln(e^{x})+\ln(e^{-x})=x-x=0$ (o bien $\ln(\cosh^{2}x-\sinh^{2}x)=\ln1=0$).
\end{enumerate}
\resultado (a) $x+1/x$; (b) $e^{4x}$; (c) $e^{-2x}$; (d) $0$.
\comprobacion (a) con $x=2$: $2\cosh(\ln2)=2\cdot1{,}25=2{,}5=2+\tfrac12$.
\end{solucion}

\begin{solucion}{eje:1-sigmoide-valores}
\datos $\sigma(x)=1/(1+e^{-x})$, $\logit(p)=\ln\frac{p}{1-p}$.
\planteamiento evaluar directamente y usar la inversa (\cref{prop:sigmoide}(iii)).
\pasos
\begin{enumerate}[label=(\alph*)]
  \item $\sigma(0)=\frac12$; $\sigma(\ln3)=\frac{1}{1+1/3}=\frac34$; $\sigma(-\ln3)=\frac{1}{1+3}=\frac14$.
  \item $x=\logit(0{,}9)=\ln9\approx2{,}197$; $x=\logit(0{,}01)=\ln\frac{1}{99}=-\ln99\approx-4{,}595$.
  \item $\sigma$ es estrictamente creciente, luego $\sigma(x)>0{,}99\iff x>\logit(0{,}99)=\ln99\approx4{,}595$.
  \item $\tanh(\ln2)=\frac{2-1/2}{2+1/2}=\frac{3/2}{5/2}=\frac35$.
\end{enumerate}
\resultado (a) $\tfrac12$, $\tfrac34$, $\tfrac14$; (b) $\ln9$ y $-\ln99$; (c) $x>\ln99$; (d) $\tfrac35$.
\comprobacion $\sigma(\ln3)+\sigma(-\ln3)=1$, como exige la simetría $\sigma(-x)=1-\sigma(x)$. Con (iv) de la \cref{prop:sigmoide}: $2\sigma(2\ln2)-1=2\sigma(\ln4)-1=2\cdot\frac45-1=\frac35$.
\end{solucion}

\begin{solucion}{eje:1-inversa-racional}
\datos $f(x)=\frac{x+1}{x-1}$, $x\ne1$.
\planteamiento componer y simplificar; una función con $f\circ f=\mathrm{id}$ es biyectiva y su propia inversa.
\pasos
\begin{enumerate}[label=\arabic*.]
  \item Para $x\ne1$, $f(x)\ne1$ (si $\frac{x+1}{x-1}=1$, entonces $x+1=x-1$, absurdo), luego $f\circ f$ está definida en $\R\setminus\{1\}$.
  \item $f(f(x))=\dfrac{\frac{x+1}{x-1}+1}{\frac{x+1}{x-1}-1}=\dfrac{(x+1)+(x-1)}{(x+1)-(x-1)}=\dfrac{2x}{2}=x$.
  \item Como $f\circ f=\mathrm{id}$, $f$ es inyectiva (si $f(x)=f(y)$, aplicando $f$, $x=y$) y sobreyectiva sobre $\R\setminus\{1\}$ (todo $y\ne1$ es $f(f(y))$). Luego $f^{-1}=f$ y la imagen es $\R\setminus\{1\}$.
  \item La gráfica de $f^{-1}$ es la simétrica de la de $f$ respecto de $y=x$ (\cref{prop:inversa-monotona}; el argumento de la simetría no usa la monotonía); como $f^{-1}=f$, la gráfica es simétrica de sí misma.
\end{enumerate}
\resultado $f\circ f=\mathrm{id}$, $f^{-1}=f$, imagen $\R\setminus\{1\}$; gráfica simétrica respecto de $y=x$.
\comprobacion $f(3)=2$ y $f(2)=3$: el punto $(3,2)$ y su simétrico $(2,3)$ están en la gráfica.
\end{solucion}

\begin{solucion}{eje:1-sigmoide-simetria}
\datos $\sigma(x)=\frac{1}{1+e^{-x}}$, $\tanh x=\frac{e^{x}-e^{-x}}{e^{x}+e^{-x}}$.
\planteamiento cálculo algebraico directo con exponenciales.
\pasos
\begin{enumerate}[label=\arabic*.]
  \item Sea $x\in\R$. Multiplicando numerador y denominador por $e^{-x}$, $\sigma(-x)=\frac{1}{1+e^{x}}=\frac{e^{-x}}{e^{-x}+1}$.
  \item Por otra parte, $1-\sigma(x)=\frac{1+e^{-x}-1}{1+e^{-x}}=\frac{e^{-x}}{1+e^{-x}}$, que coincide con lo anterior. Luego $\sigma(-x)=1-\sigma(x)$.
  \item Para (iv): $2\sigma(2x)-1=\frac{2}{1+e^{-2x}}-1=\frac{1-e^{-2x}}{1+e^{-2x}}$. Multiplicando numerador y denominador por $e^{x}$ se obtiene $\frac{e^{x}-e^{-x}}{e^{x}+e^{-x}}=\tanh x$.
  \item Sustituyendo $x$ por $x/2$ en la identidad anterior, $\tanh\frac x2=2\sigma(x)-1$, y despejando, $\sigma(x)=\frac12\bigl(1+\tanh\frac x2\bigr)$. \hfill$\blacksquare$
\end{enumerate}
\resultado $\sigma(-x)=1-\sigma(x)$, $\tanh x=2\sigma(2x)-1$ y $\sigma(x)=\frac12(1+\tanh\frac x2)$.
\comprobacion con $x=\ln3$: $\frac12(1+\tanh\frac{\ln3}{2})=\frac12\bigl(1+\frac{\sqrt3-1/\sqrt3}{\sqrt3+1/\sqrt3}\bigr)=\frac12(1+\frac{2}{4})=\frac34=\sigma(\ln3)$.
\end{solucion}

\begin{solucion}{eje:1-hiperbolicas-suma}
\datos $\cosh u\pm\sinh u=e^{\pm u}$ (\cref{prop:hiperbolicas}(iii)).
\planteamiento desarrollar el miembro derecho con exponenciales; para (c), dividir (a) entre (b).
\pasos
\begin{enumerate}[label=\arabic*.]
  \item Sean $x,y\in\R$. $\sinh x\cosh y+\cosh x\sinh y=\frac14\bigl[(e^{x}-e^{-x})(e^{y}+e^{-y})+(e^{x}+e^{-x})(e^{y}-e^{-y})\bigr]$.
  \item Desarrollando: el primer producto es $e^{x+y}+e^{x-y}-e^{-x+y}-e^{-x-y}$ y el segundo $e^{x+y}-e^{x-y}+e^{-x+y}-e^{-x-y}$. La suma es $2e^{x+y}-2e^{-(x+y)}$, y al dividir por~4 se obtiene $\sinh(x+y)$. Esto prueba (a).
  \item Análogamente, $\cosh x\cosh y+\sinh x\sinh y=\frac14\bigl[(e^{x+y}+e^{x-y}+e^{-x+y}+e^{-x-y})+(e^{x+y}-e^{x-y}-e^{-x+y}+e^{-x-y})\bigr]=\frac{e^{x+y}+e^{-(x+y)}}{2}=\cosh(x+y)$. Esto prueba (b).
  \item Por (a) y (b), $\tanh(x+y)=\frac{\sinh x\cosh y+\cosh x\sinh y}{\cosh x\cosh y+\sinh x\sinh y}$. Dividiendo numerador y denominador por $\cosh x\cosh y>0$ se obtiene $\frac{\tanh x+\tanh y}{1+\tanh x\tanh y}$, que prueba (c). El denominador no se anula porque $\abs{\tanh x\tanh y}<1$. \hfill$\blacksquare$
\end{enumerate}
\resultado las tres identidades. Frente a las trigonométricas, $\sinh$ se comporta como $\sin$, pero en $\cosh(x+y)$ y en el denominador de $\tanh(x+y)$ aparece un signo $+$ donde las trigonométricas tienen $-$.
\comprobacion con $y=x$: (b) da $\cosh2x=\cosh^{2}x+\sinh^{2}x$; en $x=0$, $1=1+0$. (c) con $x=y=\ln2$: $\tanh(\ln4)=\frac{4-1/4}{4+1/4}=\frac{15}{17}$ y $\frac{2\cdot\frac35}{1+\frac{9}{25}}=\frac{6/5}{34/25}=\frac{15}{17}$.
\end{solucion}

\begin{solucion}{eje:1-arsinh}
\datos $\sinh x=\frac{e^{x}-e^{-x}}{2}$, $\tanh x=\frac{e^{2x}-1}{e^{2x}+1}$.
\planteamiento resolver $f(x)=y$ con el cambio $u=e^{x}>0$ (o $u=e^{2x}$); una solución única para cada $y$ prueba la biyectividad.
\pasos
\begin{enumerate}[label=\arabic*.]
  \item (a) Sea $y\in\R$. Con $u=e^{x}>0$, $\sinh x=y$ equivale a $u-\frac1u=2y$, es decir, $u^{2}-2yu-1=0$, cuyas raíces son $u=y\pm\sqrt{y^{2}+1}$.
  \item Como $\sqrt{y^{2}+1}>\abs y$, la raíz $y-\sqrt{y^{2}+1}$ es negativa y $y+\sqrt{y^{2}+1}$ es positiva. Hay, pues, exactamente un $u>0$ y exactamente un $x=\ln u$: $\sinh$ es biyectiva y $\arsinh y=\ln\bigl(y+\sqrt{y^{2}+1}\bigr)$.
  \item (b) Sea $y\in\R$. Con $u=e^{2x}>0$, $\tanh x=y$ equivale a $u-1=y(u+1)$, es decir, $u(1-y)=1+y$. Si $y=1$ no hay solución; si $y\ne1$, $u=\frac{1+y}{1-y}$, que es positivo si y solo si $-1<y<1$.
  \item Por tanto, cada $y\in(-1,1)$ tiene exactamente una preimagen $x=\frac12\ln\frac{1+y}{1-y}$, y ningún $y\notin(-1,1)$ la tiene: $\tanh\colon\R\to(-1,1)$ es biyectiva con la inversa del enunciado.
  \item (c) Con $p=\frac{1+y}{2}$, $1-p=\frac{1-y}{2}$ y $\frac{p}{1-p}=\frac{1+y}{1-y}$, luego $\frac12\logit(p)=\frac12\ln\frac{1+y}{1-y}=\artanh y$. \hfill$\blacksquare$
\end{enumerate}
\resultado $\arsinh y=\ln(y+\sqrt{y^{2}+1})$ en $\R$; $\artanh y=\frac12\ln\frac{1+y}{1-y}$ en $(-1,1)$; $\artanh y=\frac12\logit\frac{1+y}{2}$.
\comprobacion $\arsinh(\frac34)=\ln(\frac34+\frac54)=\ln2$ y $\sinh(\ln2)=\frac{2-1/2}{2}=\frac34$. $\artanh(\frac35)=\frac12\ln\frac{8/5}{2/5}=\frac12\ln4=\ln2$, coherente con el \cref{eje:1-sigmoide-valores}(d).
\end{solucion}

\begin{solucion}{eje:1-estricta}
\datos $f\colon[a,b]\to\R$ estrictamente creciente.
\planteamiento la inyectividad, por casos; la segunda afirmación, por contrarrecíproco.
\pasos
\begin{enumerate}[label=\arabic*.]
  \item Sean $x\ne y$ en $[a,b]$. Si $x<y$, $f(x)<f(y)$; si $y<x$, $f(y)<f(x)$. En ambos casos $f(x)\ne f(y)$: $f$ es inyectiva.
  \item Supóngase $f(x)<f(y)$ y, por reducción al absurdo, $x\ge y$. Si $x=y$, $f(x)=f(y)$; si $x>y$, $f(x)>f(y)$. En ambos casos $f(x)\ge f(y)$, contradicción. Luego $x<y$.
  \item Ejemplo: $g\colon[0,2]\to\R$ con $g(x)=x$ en $[0,1]$ y $g(x)=4-x$ en $(1,2]$. Cada tramo es inyectivo, y los valores del primero, $[0,1]$, y del segundo, $[2,3)$, son disjuntos; luego $g$ es inyectiva. Pero $0<1{,}5<2$ y $g(0)=0<g(1{,}5)=2{,}5>g(2)=2$, de modo que $g$ no es creciente ni decreciente.
\end{enumerate}
\resultado $f$ es inyectiva y $f(x)<f(y)\Rightarrow x<y$. La función $g(x)=x$ en $[0,1]$, $g(x)=4-x$ en $(1,2]$ es inyectiva y no monótona.
\comprobacion la $g$ del ejemplo es discontinua en $x=1$; el \cref{eje:3-inyectiva-monotona} prueba que una función \emph{continua} e inyectiva en un intervalo sí es estrictamente monótona, lo que explica que el contraejemplo necesite un salto.
\end{solucion}

\begin{solucion}{eje:1-monotonia-composicion}
\datos $f,g\colon\R\to\R$.
\planteamiento aplicar las definiciones a $x<y$ arbitrarios.
\pasos
\begin{enumerate}[label=\arabic*.]
  \item (a) Sean $x<y$. Como $f$ es creciente, $f(x)\le f(y)$; como $g$ es creciente, $g(f(x))\le g(f(y))$. Luego $g\circ f$ es creciente.
  \item (b) Sean $x<y$. Entonces $f(x)\le f(y)$ y, como $g$ es decreciente, $g(f(x))\ge g(f(y))$: $g\circ f$ es decreciente.
  \item (c) Si $x<y$, $f(x)\le f(y)$ y $g(x)\le g(y)$; sumando, $(f+g)(x)\le(f+g)(y)$.
  \item (d) Contraejemplo: $f(x)=g(x)=x$ son crecientes, pero $f(x)g(x)=x^{2}$ no lo es ($(-1)^{2}>0^{2}$).
  \item Aplicación: $x\mapsto x^{3}$ es creciente y $x\mapsto\sigma(-x)$ es decreciente (es $\sigma$, creciente, compuesta con $-x$, decreciente, que por un argumento como (b) da una función decreciente). Por (b), $x\mapsto\sigma(-x^{3})$ es decreciente. \hfill$\blacksquare$
\end{enumerate}
\resultado (a)--(c) demostradas; (d) falsa en general; $\sigma(-x^{3})$ es decreciente.
\comprobacion $\sigma(-0^{3})=0{,}5>\sigma(-1)=0{,}269$. En (d), el producto de crecientes \emph{positivas} sí es creciente: si $0<f(x)\le f(y)$ y $0<g(x)\le g(y)$, entonces $f(x)g(x)\le f(y)g(y)$.
\end{solucion}

\begin{solucion}{eje:1-par-impar}
\datos $D$ simétrico, $f\colon D\to\R$.
\planteamiento unicidad: suponer $f=p+q$ y obtener $p$ y $q$; existencia: comprobar que las fórmulas obtenidas funcionan.
\pasos
\begin{enumerate}[label=\arabic*.]
  \item \emph{Unicidad.} Supóngase $f=p+q$ con $p$ par y $q$ impar. Para $x\in D$, $f(x)=p(x)+q(x)$ y $f(-x)=p(x)-q(x)$. Sumando y restando, $p(x)=\frac{f(x)+f(-x)}{2}$ y $q(x)=\frac{f(x)-f(-x)}{2}$: $p$ y $q$ quedan determinadas por $f$.
  \item \emph{Existencia.} Se definen $p$ y $q$ por esas fórmulas. Entonces $p(-x)=\frac{f(-x)+f(x)}{2}=p(x)$, $q(-x)=\frac{f(-x)-f(x)}{2}=-q(x)$ y $p+q=f$.
  \item Para $f(x)=e^{x}$: $p=\frac{e^{x}+e^{-x}}{2}=\cosh x$ y $q=\sinh x$.
  \item Para $f=\softplus$: por la \cref{prop:relu}(ii), $f(x)-f(-x)=x$, luego $q(x)=\frac x2$, y $p(x)=\softplus(x)-\frac x2$, que es la función par del \cref{eje:1-paridad}(f). \hfill$\blacksquare$
\end{enumerate}
\resultado $p(x)=\frac{f(x)+f(-x)}{2}$, $q(x)=\frac{f(x)-f(-x)}{2}$; para $e^{x}$, $\cosh$ y $\sinh$; para la softplus, $\softplus(x)-\frac x2$ y $\frac x2$.
\comprobacion $\cosh1+\sinh1=1{,}5431+1{,}1752=2{,}7183=e$.
\end{solucion}

\begin{solucion}{eje:1-softplus}
\datos $\softplus(x)=\ln(1+e^{x})$, $\sigma(x)=1/(1+e^{-x})$.
\planteamiento manipular el argumento del logaritmo; para la monotonía, componer funciones crecientes.
\pasos
\begin{enumerate}[label=\arabic*.]
  \item $\softplus(x)=\ln\bigl(e^{x}(e^{-x}+1)\bigr)=x+\ln(1+e^{-x})=x+\softplus(-x)$.
  \item $-\ln\sigma(-x)=-\ln\frac{1}{1+e^{x}}=\ln(1+e^{x})=\softplus(x)$.
  \item Monotonía: $x\mapsto e^{x}$ es estrictamente creciente, luego también $1+e^{x}$, y $\ln$ es estrictamente creciente; la composición de estrictamente crecientes es estrictamente creciente (como en el \cref{eje:1-monotonia-composicion}(a), con desigualdades estrictas).
  \item Cambiando $x$ por $-x$ en el paso~2: $\softplus(-x)=-\ln\sigma(x)$, es decir, $\ln\sigma(x)=-\softplus(-x)$.
  \item Utilidad numérica: para $x=-1000$, $\sigma(x)=1/(1+e^{1000})$ es menor que el menor número representable en coma flotante y se redondea a~0, y $\ln0=-\infty$. En cambio, $-\softplus(1000)=-(1000+\softplus(-1000))=-(1000+\ln(1+e^{-1000}))$ se calcula sin problemas y vale $-1000$ con precisión de máquina. \hfill$\blacksquare$
\end{enumerate}
\resultado las identidades del enunciado; $\ln\sigma(x)=-\softplus(-x)$ evita el desbordamiento por defecto.
\comprobacion con $x=0$: $\ln\sigma(0)=\ln\frac12=-\ln2=-\softplus(0)$.
\end{solucion}
% ---------------------------------------------------------------------
\section{Tema 2}

\begin{solucion}{eje:2-negacion}
\datos las definiciones~\ref{def:limite-sucesion} y~\ref{def:limite-funcion}.
\planteamiento negar cuantificadores: $\neg\forall=\exists\neg$, $\neg\exists=\forall\neg$, y $\neg(P\Rightarrow Q)\equiv P\wedge\neg Q$ (módulo~0.1).
\pasos
\begin{enumerate}[label=\arabic*.]
  \item $a_n\not\to L$: $\exists\varepsilon>0\ \forall N\in\N\ \exists n\ge N:\ \abs{a_n-L}\ge\varepsilon$.
  \item $\lim_{x\to a}f(x)\ne L$: $\exists\varepsilon>0\ \forall\delta>0\ \exists x:\ 0<\abs{x-a}<\delta\ \wedge\ \abs{f(x)-L}\ge\varepsilon$.
  \item Para probar $a_n\not\to L$ hay que exhibir un $\varepsilon>0$ concreto y mostrar que, por grande que sea $N$, hay términos con índice $\ge N$ a distancia al menos $\varepsilon$ de $L$; equivalentemente, que hay infinitos términos fuera del intervalo $(L-\varepsilon,L+\varepsilon)$.
\end{enumerate}
\resultado las dos negaciones anteriores.
\comprobacion para $a_n=(-1)^{n}$ y $L=1$, $\varepsilon=1$ sirve: todos los términos de índice impar valen $-1$ y distan~2 de~1.
\end{solucion}

\begin{solucion}{eje:2-N}
\datos $a_n=\frac{2n}{n+3}$, $L=2$.
\planteamiento calcular $\abs{a_n-2}$ exactamente y despejar $n$.
\pasos
\begin{enumerate}[label=\arabic*.]
  \item $\abs{a_n-2}=\abs{\frac{2n-2n-6}{n+3}}=\frac{6}{n+3}$, que es decreciente en $n$.
  \item $\frac{6}{n+3}<0{,}01\iff n+3>600\iff n>597$. El menor $N$ es $598$.
  \item $\frac{6}{n+3}<10^{-6}\iff n+3>6\cdot10^{6}\iff n>5\,999\,997$. El menor $N$ es $5\,999\,998$.
\end{enumerate}
\resultado $N=598$ para $\varepsilon=0{,}01$ y $N=5\,999\,998$ para $\varepsilon=10^{-6}$.
\comprobacion $\abs{a_{597}-2}=6/600=0{,}01$, que no es $<0{,}01$; $\abs{a_{598}-2}=6/601<0{,}01$. Como la cota es decreciente, todos los $n\ge N$ sirven.
\end{solucion}

\begin{solucion}{eje:2-sucesiones-calculo}
\datos seis sucesiones.
\planteamiento dividir por la potencia dominante, multiplicar por el conjugado, usar el sándwich o la \cref{prop:notables-exp}(iii).
\pasos
\begin{enumerate}[label=(\alph*)]
  \item Dividiendo por $n^{2}$: $\frac{3-1/n+4/n^{2}}{5+2/n^{2}}\to\frac35$ (\cref{teo:algebra-sucesiones} y $1/n\to0$).
  \item $\sqrt{n^{2}+n}-n=\frac{(n^{2}+n)-n^{2}}{\sqrt{n^{2}+n}+n}=\frac{n}{\sqrt{n^{2}+n}+n}=\frac{1}{\sqrt{1+1/n}+1}\to\frac12$, usando la continuidad de la raíz (\cref{tema:3}).
  \item $\frac{2^{n}+3^{n}}{3^{n+1}}=\frac13\bigl(\frac23\bigr)^{n}+\frac13\to\frac13$ por la \cref{prop:geometrica}.
  \item $-\frac1n\le\frac{\sin n}{n}\le\frac1n$; por el sándwich, el límite es~0.
  \item Por la \cref{prop:notables-exp}(iii) con $t=-2$: $e^{-2}\approx0{,}1353$.
  \item $3^{n}\le2^{n}+3^{n}\le2\cdot3^{n}$, luego $3\le\sqrt[n]{2^{n}+3^{n}}\le3\cdot2^{1/n}$. Como $2^{1/n}\to1$ (\cref{eje:2-raizn}(b)), el sándwich da~3.
\end{enumerate}
\resultado (a) $\tfrac35$; (b) $\tfrac12$; (c) $\tfrac13$; (d) $0$; (e) $e^{-2}$; (f) $3$.
\comprobacion numérica: para $n=10^{6}$, (b) vale $0{,}49999988$ y (e) vale $0{,}1353350$, frente a $e^{-2}=0{,}1353353$.
\end{solucion}

\begin{solucion}{eje:2-funciones-calculo}
\datos cinco límites de funciones racionales o con raíces.
\planteamiento factorizar y simplificar el factor que se anula; en $\infty$, dividir por la potencia dominante; en las raíces, multiplicar por el conjugado.
\pasos
\begin{enumerate}[label=(\alph*)]
  \item $\frac{(x-2)(x+2)}{(x-1)(x-2)}=\frac{x+2}{x-1}$ para $x\ne2$; el límite es $\frac41=4$.
  \item $\frac{x+1}{(x-1)(x+1)}=\frac{1}{x-1}$ para $x\ne\pm1$. Cuando $x\to1^{+}$, $x-1\to0^{+}$ y el cociente tiende a $+\infty$; cuando $x\to1^{-}$, a $-\infty$. El límite (bilateral) no existe.
  \item $\frac{2-1/x^{2}}{1+4/x}\to2$.
  \item $\frac{\sqrt{1+x}-1}{x}=\frac{(1+x)-1}{x(\sqrt{1+x}+1)}=\frac{1}{\sqrt{1+x}+1}\to\frac12$.
  \item $x^{3}+8=(x+2)(x^{2}-2x+4)$, luego el cociente es $x^{2}-2x+4\to4+4+4=12$.
\end{enumerate}
\resultado (a) $4$; (b) $+\infty$ y $-\infty$; (c) $2$; (d) $\tfrac12$; (e) $12$.
\comprobacion numérica: en $x=2{,}000001$, (a) vale $3{,}999997$; en $x=10^{-6}$, (d) vale $0{,}49999988$.
\end{solucion}

\begin{solucion}{eje:2-notables}
\datos la \cref{tab:notables}.
\planteamiento reducir cada límite a uno conocido mediante un cambio de variable o un producto de factores con límite conocido; demostrar primero las entradas de la tabla que faltan.
\pasos
\begin{enumerate}[label=\arabic*.]
  \item \emph{Entradas pendientes.} $\frac{\tan x}{x}=\frac{\sin x}{x}\cdot\frac{1}{\cos x}\to1$. $\frac{a^{x}-1}{x}=\ln a\cdot\frac{e^{u}-1}{u}$ con $u=x\ln a\to0$ (si $a\ne1$), que tiende a $\ln a$ por~\eqref{eq:lim-exp}. Para $p\ne0$, $\frac{(1+x)^{p}-1}{x}=\frac{e^{u}-1}{u}\cdot\frac{p\ln(1+x)}{x}$ con $u=p\ln(1+x)\to0$, que tiende a $1\cdot p$ por la \cref{prop:notables-exp}(i). $\frac{\tanh x}{x}=\frac{e^{2x}-1}{2x}\cdot\frac{2}{e^{2x}+1}\to1\cdot\frac22=1$.
  \item (a) $\frac{\sin5x}{3x}=\frac53\cdot\frac{\sin5x}{5x}\to\frac53$.
  \item (b) $1$, por el paso~1.
  \item (c) $\frac{e^{2x}-1}{\sin x}=2\cdot\frac{e^{2x}-1}{2x}\cdot\frac{x}{\sin x}\to2$.
  \item (d) $\ln2$, por el paso~1 con $a=2$.
  \item (e) $1$, por el paso~1.
  \item (f) Con $t=1/x\to0^{+}$: $x\sin(1/x)=\frac{\sin t}{t}\to1$.
\end{enumerate}
\resultado (a) $\tfrac53$; (b) $1$; (c) $2$; (d) $\ln2\approx0{,}6931$; (e) $1$; (f) $1$.
\comprobacion numérica con $x=10^{-6}$: (a) $1{,}6666667$; (c) $2{,}000002$; (d) $0{,}6931474$.
\end{solucion}

\begin{solucion}{eje:2-asintotas}
\datos cuatro funciones.
\planteamiento verticales: puntos donde algún lateral es infinito; horizontales y oblicuas: \cref{prop:oblicua}.
\pasos
\begin{enumerate}[label=(\alph*)]
  \item División: $2x^{2}-3x=(x+1)(2x-5)+5$, luego $f(x)=2x-5+\frac{5}{x+1}$. Vertical $x=-1$ ($f\to\pm\infty$ por la derecha y la izquierda, respectivamente); oblicua $y=2x-5$ en $\pm\infty$.
  \item $g$ está definida en $\R$ y es continua: no hay verticales. Para $x>0$, $g(x)=\frac{1}{\sqrt{1+1/x^{2}}}\to1$; para $x<0$, $\sqrt{x^{2}}=-x$ y $g(x)=\frac{-1}{\sqrt{1+1/x^{2}}}\to-1$. Horizontales $y=1$ en $+\infty$ e $y=-1$ en $-\infty$.
  \item $h(x)=xe^{1/x}$, $x\ne0$. Si $x\to0^{+}$, con $t=1/x\to+\infty$, $h=e^{t}/t\to+\infty$ (\cref{prop:crecimientos}): vertical $x=0$ por la derecha. Si $x\to0^{-}$, $e^{1/x}\to0$ y $h\to0$. En $\pm\infty$: $m=\lim e^{1/x}=1$ y $b=\lim x(e^{1/x}-1)=\lim_{t\to0}\frac{e^{t}-1}{t}=1$. Oblicua $y=x+1$.
  \item $k(x)=\softplus(2x)$: en $-\infty$, $k\to\ln1=0$ (horizontal $y=0$); en $+\infty$, $k(x)-2x=\softplus(-2x)\to0$ (oblicua $y=2x$). No hay verticales.
\end{enumerate}
\resultado (a) $x=-1$, $y=2x-5$; (b) $y=1$ y $y=-1$; (c) $x=0$ (por la derecha), $y=x+1$; (d) $y=0$ en $-\infty$, $y=2x$ en $+\infty$.
\comprobacion (c) en $x=100$: $100e^{0{,}01}=101{,}005$, frente a $x+1=101$.
\end{solucion}

\begin{solucion}{eje:2-eN}
\datos cuatro sucesiones y sus límites.
\planteamiento trabajo previo: acotar $\abs{a_n-L}$ por $C/n$ o $C/\sqrt n$; después, redactar con los cuantificadores en orden.
\pasos
\begin{enumerate}[label=(\alph*)]
  \item $\abs{\frac{n^{2}}{n^{2}+1}-1}=\frac{1}{n^{2}+1}<\frac1n$. Sea $\varepsilon>0$ y $N\in\N$ con $N>1/\varepsilon$. Para $n\ge N$, $\abs{a_n-1}<\frac1n\le\frac1N<\varepsilon$.
  \item Sea $\varepsilon>0$ y $N\in\N$ con $N>1/\varepsilon^{2}$. Para $n\ge N$, $\sqrt n\ge\sqrt N>1/\varepsilon$, luego $\frac{1}{\sqrt n}<\varepsilon$.
  \item $\abs{\frac{(-1)^{n}}{n^{2}}}=\frac{1}{n^{2}}\le\frac1n$. Sea $\varepsilon>0$ y $N>1/\varepsilon$; para $n\ge N$, $\abs{a_n}\le\frac1n<\varepsilon$.
  \item $\abs{\frac{n+5}{2n-1}-\frac12}=\frac{\abs{2n+10-2n+1}}{2(2n-1)}=\frac{11}{2(2n-1)}\le\frac{11}{2n}$, porque $2n-1\ge n$ para $n\ge1$. Sea $\varepsilon>0$ y $N>\frac{11}{2\varepsilon}$; para $n\ge N$, $\abs{a_n-\frac12}\le\frac{11}{2n}<\varepsilon$. \hfill$\blacksquare$
\end{enumerate}
\resultado las cuatro convergencias, con $N>1/\varepsilon$, $N>1/\varepsilon^{2}$, $N>1/\varepsilon$ y $N>11/(2\varepsilon)$.
\comprobacion (d) con $\varepsilon=0{,}1$: $N=56$ y $\abs{a_{56}-\frac12}=\frac{11}{222}\approx0{,}0495<0{,}1$.
\end{solucion}

\begin{solucion}{eje:2-ed}
\datos cuatro límites.
\planteamiento factorizar $\abs{f(x)-L}$ como $\abs{x-a}$ por un factor y acotar ese factor (\cref{ej:cuadratico}).
\pasos
\begin{enumerate}[label=(\alph*)]
  \item $\abs{(5x+8)-18}=5\abs{x-2}$. Dado $\varepsilon>0$, sea $\delta=\varepsilon/5$. Si $0<\abs{x-2}<\delta$, $\abs{(5x+8)-18}<5\delta=\varepsilon$.
  \item $\abs{x^{2}+2x-8}=\abs{x-2}\abs{x+4}$. Si $\abs{x-2}<1$, entonces $1<x<3$ y $\abs{x+4}<7$. Dado $\varepsilon>0$, sea $\delta=\min(1,\varepsilon/7)$. Si $0<\abs{x-2}<\delta$, $\abs{x^{2}+2x-8}<7\abs{x-2}<7\cdot\frac{\varepsilon}{7}=\varepsilon$.
  \item Para $x\ne3$, $\frac{x^{2}-9}{x-3}=x+3$ y $\abs{(x+3)-6}=\abs{x-3}$. Dado $\varepsilon>0$, $\delta=\varepsilon$ sirve; la condición $0<\abs{x-3}$ es esencial, pues la función no está definida en~3.
  \item Dado $\varepsilon>0$, sea $\delta=\varepsilon^{1/4}$. Si $0<\abs{x}<\delta$, $\abs{x^{4}}=\abs{x}^{4}<\delta^{4}=\varepsilon$. \hfill$\blacksquare$
\end{enumerate}
\resultado $\delta=\varepsilon/5$; $\delta=\min(1,\varepsilon/7)$; $\delta=\varepsilon$; $\delta=\varepsilon^{1/4}$.
\comprobacion (b) con $\varepsilon=0{,}07$: $\delta=0{,}01$ y, en el peor caso $x=2{,}01$, $x^{2}+2x=8{,}0601$, a distancia $0{,}0601<0{,}07$.
\end{solucion}

\begin{solucion}{eje:2-ed-dificiles}
\datos tres límites.
\planteamiento como en el ejercicio anterior; en (b), acotar $1/\abs x$ alejando $x$ de~0; en (c), usar la \cref{def:limites-variantes}(ii).
\pasos
\begin{enumerate}[label=(\alph*)]
  \item Para $x\ge0$, $\abs{\sqrt x-2}=\frac{\abs{x-4}}{\sqrt x+2}\le\frac{\abs{x-4}}{2}$. Dado $\varepsilon>0$, sea $\delta=\min(4,2\varepsilon)$; la condición $\delta\le4$ garantiza $x>0$, de modo que $\sqrt x$ está definida. Si $0<\abs{x-4}<\delta$, $\abs{\sqrt x-2}\le\frac{\abs{x-4}}{2}<\varepsilon$.
  \item $\abs{\frac1x-1}=\frac{\abs{x-1}}{\abs x}$. Si $\abs{x-1}<\frac12$, entonces $x>\frac12$ y $\frac{1}{\abs x}<2$. Dado $\varepsilon>0$, sea $\delta=\min(\frac12,\frac\varepsilon2)$; si $0<\abs{x-1}<\delta$, $\abs{\frac1x-1}<2\abs{x-1}<\varepsilon$.
  \item Sea $K\in\R$. Si $K\le0$, $\frac{1}{x^{2}}>0\ge K$ para todo $x\ne0$ y cualquier $\delta$ sirve. Si $K>0$, sea $\delta=1/\sqrt K$; si $0<\abs x<\delta$, $x^{2}<1/K$ y $\frac{1}{x^{2}}>K$. \hfill$\blacksquare$
\end{enumerate}
\resultado $\delta=\min(4,2\varepsilon)$; $\delta=\min(\frac12,\frac\varepsilon2)$; $\delta=1/\sqrt K$.
\comprobacion (a) con $\varepsilon=0{,}1$: $\delta=0{,}2$; $\sqrt{4{,}2}=2{,}0494$ y $\sqrt{3{,}8}=1{,}9494$, ambos a distancia menor que $0{,}1$ de~2.
\end{solucion}

\begin{solucion}{eje:2-no-existe}
\datos $f(x)=\frac{\abs{x-1}}{x-1}$; la función de Dirichlet $D$.
\planteamiento (a) reducción al absurdo con la definición; (b) dos sucesiones con imágenes de límites distintos (\cref{teo:heine}).
\pasos
\begin{enumerate}[label=\arabic*.]
  \item (a) Supóngase $\lim_{x\to1}f(x)=L$. Con $\varepsilon=1$ existe $\delta>0$ tal que $\abs{f(x)-L}<1$ si $0<\abs{x-1}<\delta$.
  \item Para $x=1+\delta/2$, $f(x)=1$, luego $\abs{1-L}<1$ y $L>0$. Para $x=1-\delta/2$, $f(x)=-1$, luego $\abs{-1-L}<1$ y $L<0$. Contradicción: el límite no existe.
  \item (b) Sean $x_n=1/n\in\Q$ e $y_n=\sqrt2/n\notin\Q$ (si fuera racional, $\sqrt2=ny_n$ también lo sería). Ambas tienden a~0 sin anularse, y $D(x_n)=1\to1$, $D(y_n)=0\to0$. Por el \cref{teo:heine}, $\lim_{x\to0}D(x)$ no existe. \hfill$\blacksquare$
\end{enumerate}
\resultado ninguno de los dos límites existe.
\comprobacion en (a) los límites laterales existen y valen $1$ y $-1$; en (b) ni siquiera existen los laterales, porque el mismo argumento vale con sucesiones positivas o negativas.
\end{solucion}

\begin{solucion}{eje:2-producto}
\datos $f\to L$ y $g\to M$ cuando $x\to a$.
\planteamiento $f(x)g(x)-LM=f(x)\bigl(g(x)-M\bigr)+M\bigl(f(x)-L\bigr)$; acotar $\abs{f(x)}$ cerca de $a$.
\pasos
\begin{enumerate}[label=\arabic*.]
  \item Con $\varepsilon=1$ en la definición para $f$, existe $\delta_1>0$ tal que $\abs{f(x)-L}<1$, y por tanto $\abs{f(x)}<\abs L+1=:K$, si $0<\abs{x-a}<\delta_1$.
  \item Sea $\varepsilon>0$. Existe $\delta_2>0$ con $\abs{g(x)-M}<\frac{\varepsilon}{2K}$ si $0<\abs{x-a}<\delta_2$, y existe $\delta_3>0$ con $\abs{f(x)-L}<\frac{\varepsilon}{2(\abs M+1)}$ si $0<\abs{x-a}<\delta_3$.
  \item Sea $\delta=\min(\delta_1,\delta_2,\delta_3)$. Si $0<\abs{x-a}<\delta$,
  \begin{align*}
  \abs{f(x)g(x)-LM}&\le\abs{f(x)}\abs{g(x)-M}+\abs M\abs{f(x)-L}\\
  &<K\cdot\frac{\varepsilon}{2K}+\abs M\cdot\frac{\varepsilon}{2(\abs M+1)}<\frac\varepsilon2+\frac\varepsilon2=\varepsilon.
  \end{align*}
  Luego $\lim_{x\to a}f(x)g(x)=LM$. \hfill$\blacksquare$
\end{enumerate}
\resultado la ley del producto.
\comprobacion se divide por $\abs M+1$ y no por $\abs M$ para cubrir el caso $M=0$; con $f(x)=g(x)=x$ y $a=2$, la ley da $\lim x^{2}=4$, coherente con el \cref{ej:cuadratico}.
\end{solucion}

\begin{solucion}{eje:2-cociente}
\datos $b_n\to M\ne0$, $b_n\ne0$; la sucesión $(a_n)$ de reales no nulos.
\planteamiento (a) alejar $b_n$ de~0 y acotar $\abs{1/b_n-1/M}$; (b) contraejemplo en (i) y (a) en (ii).
\pasos
\begin{enumerate}[label=\arabic*.]
  \item Con $\varepsilon=\abs M/2$ existe $N_1$ tal que $\abs{b_n-M}<\abs M/2$ para $n\ge N_1$, y entonces $\abs{b_n}\ge\abs M-\abs{M-b_n}>\abs M/2$.
  \item Para $n\ge N_1$, $\abs{\frac1{b_n}-\frac1M}=\frac{\abs{M-b_n}}{\abs{b_n}\abs M}<\frac{2}{M^{2}}\abs{b_n-M}$.
  \item Sea $\varepsilon>0$. Existe $N_2$ con $\abs{b_n-M}<\frac{\varepsilon M^{2}}{2}$ para $n\ge N_2$. Para $n\ge\max(N_1,N_2)$, $\abs{\frac1{b_n}-\frac1M}<\varepsilon$. Esto prueba que $1/b_n\to1/M$, y con el \cref{teo:algebra-sucesiones}(iii), $a_n/b_n\to L/M$.
  \item (b)(i) No basta: $a_n=1/n\to0$, y $1/a_n=n$ no converge.
  \item (b)(ii) Basta: es exactamente el apartado (a). \hfill$\blacksquare$
\end{enumerate}
\resultado $1/b_n\to1/M$; la condición (i) no garantiza la convergencia de $(1/a_n)$ y la (ii) sí.
\comprobacion $b_n=2+1/n\to2$: $1/b_n=\frac{n}{2n+1}\to\frac12$, y la cota del paso~2 da $\abs{\frac{n}{2n+1}-\frac12}=\frac{1}{2(2n+1)}<\frac{2}{4}\cdot\frac1n$.
\end{solucion}

\begin{solucion}{eje:2-alternada}
\datos $a_n=(-1)^{n}$; una sucesión convergente cualquiera.
\planteamiento probar primero (b) y deducir (a) y el caso $r\le-1$ por contrarrecíproco.
\pasos
\begin{enumerate}[label=\arabic*.]
  \item (b) Sea $a_n\to L$ y $\varepsilon>0$. Existe $N$ con $\abs{a_n-L}<\varepsilon/2$ para $n\ge N$; entonces también $\abs{a_{n+1}-L}<\varepsilon/2$ y $\abs{a_{n+1}-a_n}\le\abs{a_{n+1}-L}+\abs{L-a_n}<\varepsilon$. Luego $a_{n+1}-a_n\to0$.
  \item (a) Para $a_n=(-1)^{n}$, $\abs{a_{n+1}-a_n}=2$ para todo $n$, que no tiende a~0. Por el contrarrecíproco de (b), $(a_n)$ diverge.
  \item Si $r\le-1$, $\abs{r^{n+1}-r^{n}}=\abs r^{n}\abs{r-1}\ge1\cdot2=2$, y $(r^{n})$ diverge por la misma razón.
  \item (c) El recíproco es falso: con $a_n=\sqrt n$, $a_{n+1}-a_n=\frac{1}{\sqrt{n+1}+\sqrt n}\to0$, pero $\sqrt n\to+\infty$. \hfill$\blacksquare$
\end{enumerate}
\resultado (a) y (b) demostrados; (c) el recíproco es falso.
\comprobacion $\sqrt{10^{6}+1}-\sqrt{10^{6}}\approx0{,}0005$, mientras que $\sqrt{10^{6}}=1000$.
\end{solucion}

\begin{solucion}{eje:2-recurrencia}
\datos $a_1=1$, $a_{n+1}=\sqrt{2+a_n}$.
\planteamiento inducción para acotación y monotonía; \cref{teo:monotona}; paso al límite en la recurrencia.
\pasos
\begin{enumerate}[label=\arabic*.]
  \item \emph{Acotación.} $0<a_1=1<2$. Si $0<a_n<2$, entonces $2<2+a_n<4$ y $\sqrt2<a_{n+1}<2$. Por inducción, $0<a_n<2$ para todo $n$.
  \item \emph{Monotonía.} Como $a_n>0$ y $a_{n+1}>0$, $a_{n+1}>a_n\iff2+a_n>a_n^{2}\iff(2-a_n)(1+a_n)>0$, que es cierto porque $0<a_n<2$. Luego $(a_n)$ es estrictamente creciente.
  \item Por el \cref{teo:monotona}, $a_n\to L$ con $1\le L\le2$.
  \item Elevando al cuadrado, $a_{n+1}^{2}=2+a_n$. Como $(a_{n+1})$ también tiende a $L$, el álgebra de límites da $L^{2}=2+L$, es decir, $(L-2)(L+1)=0$, y como $L\ge1$, $L=2$. \hfill$\blacksquare$
\end{enumerate}
\resultado $(a_n)$ es creciente, está en $(0,2)$ y converge a~2.
\comprobacion $a_2=\sqrt3\approx1{,}732$, $a_3\approx1{,}9319$, $a_4\approx1{,}9829$, $a_5\approx1{,}9957$.
\end{solucion}

\begin{solucion}{eje:2-raizn}
\datos $\sqrt[n]{n}=1+h_n$ con $h_n\ge0$ (porque $n\ge1$).
\planteamiento desarrollar $(1+h_n)^{n}$ por el binomio quedándose con un término; sándwich.
\pasos
\begin{enumerate}[label=\arabic*.]
  \item (a) Sea $n\ge2$. Todos los términos del binomio son $\ge0$, luego $n=(1+h_n)^{n}\ge\binom n2h_n^{2}=\frac{n(n-1)}{2}h_n^{2}$.
  \item Despejando, $0\le h_n\le\sqrt{\frac{2}{n-1}}$, y el miembro derecho tiende a~0. Por el sándwich, $h_n\to0$ y $\sqrt[n]n\to1$.
  \item (b) Si $a\ge1$: para $n\ge a$, $1\le\sqrt[n]a\le\sqrt[n]n$, y el sándwich da $\sqrt[n]a\to1$.
  \item Si $0<a<1$: $\sqrt[n]a=1/\sqrt[n]{1/a}$ con $1/a>1$; por el caso anterior y el \cref{teo:algebra-sucesiones}(iv), $\sqrt[n]a\to1$. \hfill$\blacksquare$
\end{enumerate}
\resultado $\sqrt[n]n\to1$ y $\sqrt[n]a\to1$ para $a>0$.
\comprobacion $\sqrt[100]{100}\approx1{,}0471$ y la cota da $h_{100}\le\sqrt{2/99}\approx0{,}142$.
\end{solucion}

\begin{solucion}{eje:2-laterales}
\datos la \cref{def:limite-funcion} y la \cref{def:limites-variantes}(i).
\planteamiento ($\Rightarrow$) el mismo $\delta$ sirve; ($\Leftarrow$) tomar el mínimo de los dos $\delta$.
\pasos
\begin{enumerate}[label=\arabic*.]
  \item ($\Rightarrow$) Sea $\varepsilon>0$ y $\delta$ el del límite. Si $a<x<a+\delta$, entonces $0<\abs{x-a}<\delta$ y $\abs{f(x)-L}<\varepsilon$; igual si $a-\delta<x<a$. Luego ambos laterales valen $L$.
  \item ($\Leftarrow$) Sea $\varepsilon>0$, y sean $\delta_{+}$ y $\delta_{-}$ los de los límites laterales. Con $\delta=\min(\delta_{+},\delta_{-})$, si $0<\abs{x-a}<\delta$, o bien $a<x<a+\delta_{+}$ o bien $a-\delta_{-}<x<a$; en ambos casos $\abs{f(x)-L}<\varepsilon$.
  \item Aplicación: $\frac{\ReLU(x)}{x}=1$ si $x>0$ y $0$ si $x<0$. Los laterales en~0 son $1$ y $0$, distintos, luego el límite no existe. \hfill$\blacksquare$
\end{enumerate}
\resultado la equivalencia; $\lim_{x\to0}\ReLU(x)/x$ no existe.
\comprobacion los laterales $1$ y $0$ son las derivadas laterales de la $\ReLU$ en~0 (\cref{ej:def-derivada}(d)).
\end{solucion}
% ---------------------------------------------------------------------
\section{Tema 3}

\begin{solucion}{eje:3-verdadero-falso}
\datos cuatro afirmaciones sobre funciones continuas.
\planteamiento buscar contraejemplos para las falsas y demostrar la verdadera.
\pasos
\begin{enumerate}[label=(\alph*)]
  \item \textbf{Falsa.} $\cos$ en $[-\pi,2\pi]$: $\cos(-\pi)=-1<0<1=\cos(2\pi)$, y $\cos$ se anula en $-\frac\pi2$, $\frac\pi2$ y $\frac{3\pi}{2}$: tres ceros. Bolzano garantiza al menos uno, no uno solo.
  \item \textbf{Falsa.} $(x-1)^{2}$ en $[-2,2]$: vale $9$ y $1$ en los extremos y se anula en $x=1$.
  \item \textbf{Falsa.} Sea $f(x)=1$ si $x\in\Q$ y $f(x)=-1$ si no. $\abs f=1$ es continua, pero $f$ no es continua en ningún punto (como la función de Dirichlet, \cref{eje:2-no-existe}).
  \item \textbf{Verdadera.} Por la desigualdad triangular inversa, $\bigl|\abs{f(x)}-\abs{f(a)}\bigr|\le\abs{f(x)-f(a)}$; el mismo $\delta$ que sirve para $f$ sirve para $\abs f$. (También: $\abs f$ es la composición de $f$ con $\abs{\cdot}$, \cref{prop:operaciones-continuas}.)
\end{enumerate}
\resultado (a) F; (b) F; (c) F; (d) V.
\comprobacion (c) y (d) muestran que la implicación «$f$ continua $\Rightarrow\abs f$ continua» no tiene recíproco.
\end{solucion}

\begin{solucion}{eje:3-hipotesis-weierstrass}
\datos el \cref{teo:weierstrass} exige $f$ continua en un intervalo cerrado y acotado.
\planteamiento en cada caso se elimina exactamente una hipótesis.
\pasos
\begin{enumerate}[label=(\alph*)]
  \item $f(x)=1/x$ en $(0,1]$: continua, no acotada superiormente; no tiene máximo. Falla que el intervalo sea cerrado.
  \item $f(x)=\arctan x$ en $[0,\infty)$: continua y acotada, con $\sup=\pi/2$ no alcanzado. Falla que el intervalo sea acotado.
  \item $f(x)=x$ en $[0,1)$ y $f(1)=0$: acotada, con $\sup=1$ no alcanzado; solo es discontinua en $x=1$. Falla la continuidad.
\end{enumerate}
\resultado los tres ejemplos anteriores.
\comprobacion en (b), $\arctan x<\pi/2$ para todo $x$, pero $\arctan x\to\pi/2$: el supremo es $\pi/2$ y no es máximo.
\end{solucion}

\begin{solucion}{eje:3-clasificar}
\datos cinco funciones.
\planteamiento por la \cref{prop:elementales-continuas}, cada función es continua en su dominio salvo, quizá, en los puntos de cambio de definición; en los puntos aislados donde no está definida se estudian los límites laterales y se clasifica según la \cref{def:discontinuidades}.
\pasos
\begin{enumerate}[label=(\alph*)]
  \item $f(x)=\frac{x-2}{(x-2)(x+2)}=\frac{1}{x+2}$ para $x\ne\pm2$. En $x=2$, el límite es $\frac14$: discontinuidad \emph{evitable}. En $x=-2$, los laterales son $\pm\infty$: \emph{esencial} (asíntota vertical).
  \item $1/\sin x$ no está definida en $x=k\pi$, $k\in\Z$, y allí $\sin x\to0$ con signos opuestos a cada lado: laterales infinitos, discontinuidad \emph{esencial}.
  \item $\floor x$ es constante en cada $[n,n+1)$, luego continua en los no enteros. En $n\in\Z$: $\lim_{x\to n^{-}}\floor x=n-1$ y $\lim_{x\to n^{+}}\floor x=n$: discontinuidad \emph{de salto}, de altura~1.
  \item Para $x\ne0$, $k$ es continua. Cuando $x\to0$, $-1/x^{2}\to-\infty$ y $e^{-1/x^{2}}\to0=k(0)$. $k$ es \textbf{continua en todo $\R$}.
  \item $\abs x/x$ vale $1$ para $x>0$ y $-1$ para $x<0$: en $x=0$, salto (no evitable, porque los laterales difieren).
\end{enumerate}
\resultado (a) evitable en $2$, esencial en $-2$; (b) esenciales en $k\pi$; (c) saltos en los enteros; (d) continua; (e) salto en~0.
\comprobacion (d) numérica: $k(0{,}1)=e^{-100}\approx3{,}7\cdot10^{-44}$.
\end{solucion}

\begin{solucion}{eje:3-constantes}
\datos cuatro funciones definidas a trozos.
\planteamiento en el punto de unión, igualar el valor de la función con los límites laterales.
\pasos
\begin{enumerate}[label=(\alph*)]
  \item Para $x\ne-2$, $\frac{(x+1)(x+2)}{x+2}=x+1\to-1$: $k=-1$.
  \item En $x=4$: el límite por la izquierda es $e^{4k}$ y por la derecha y el valor, $7$. $e^{4k}=7$, $k=\frac{\ln7}{4}\approx0{,}4865$.
  \item En $x=3$: $\sqrt{3k}=4$, $3k=16$, $k=\frac{16}{3}$.
  \item En $x=0$: $\lim_{x\to0^{+}}\frac{\sin2x}{x}=2\lim\frac{\sin2x}{2x}=2$; $\lim_{x\to0^{-}}(ax+b)=b$; $u(0)=a$. Continuidad: $a=b=2$.
\end{enumerate}
\resultado (a) $k=-1$; (b) $k=\frac{\ln7}{4}$; (c) $k=\frac{16}{3}$; (d) $a=b=2$.
\comprobacion (c): con $k=16/3$, $h(3)=\sqrt{16}=4=3+1$.
\end{solucion}

\begin{solucion}{eje:3-biseccion}
\datos $g(t)=\cos t-t^{3}$, continua en $\R$.
\planteamiento Bolzano en $[0,1]$ y bisección (\cref{teo:bolzano}).
\pasos
\begin{enumerate}[label=\arabic*.]
  \item (a) $g(0)=1>0$ y $g(1)=\cos1-1\approx-0{,}460<0$. Por Bolzano hay una solución en $(0,1)$.
  \item (b) $g(0{,}5)\approx0{,}753>0$: raíz en $[0{,}5;1]$. $g(0{,}75)\approx0{,}310>0$: en $[0{,}75;1]$. $g(0{,}875)\approx-0{,}029<0$: en $[0{,}75;0{,}875]$. $g(0{,}8125)\approx0{,}151>0$: en $[0{,}8125;0{,}875]$, de longitud $1/16$.
  \item (c) Tras $n$ pasos el intervalo mide $2^{-n}$. Para que cualquier punto del intervalo diste menos de $10^{-6}$ de la raíz basta $2^{-n}<10^{-6}$, es decir, $n\ge20$, porque $2^{20}=1\,048\,576$. Si se toma el punto medio, el error es menor que $2^{-(n+1)}$ y bastan $19$ pasos.
\end{enumerate}
\resultado (a) existe solución en $(0,1)$; (b) está en $[0{,}8125;\,0{,}875]$; (c) 20 pasos (19 con el punto medio).
\comprobacion la solución, calculada numéricamente, es $t\approx0{,}86547$, dentro del intervalo de (b).
\end{solucion}

\begin{solucion}{eje:3-raices}
\datos tres ecuaciones.
\planteamiento evaluar en enteros y aplicar Bolzano en cada cambio de signo.
\pasos
\begin{enumerate}[label=(\alph*)]
  \item $p(x)=x^{3}-3x+1$: $p(-2)=-1$, $p(-1)=3$, $p(0)=1$, $p(1)=-1$, $p(2)=3$. Cambios de signo en $(-2,-1)$, $(0,1)$ y $(1,2)$: tres raíces, y no hay más por ser de grado~3.
  \item $h(x)=e^{x}-3x$: $h(0)=1$, $h(1)=e-3\approx-0{,}28$, $h(2)=e^{2}-6\approx1{,}39$. Raíces en $(0,1)$ y $(1,2)$.
  \item $g(x)=\sigma(x)-x$: $g(0)=\frac12>0$ y $g(1)=\sigma(1)-1\approx-0{,}27<0$. Raíz en $(0,1)$.
\end{enumerate}
\resultado (a) en $(-2,-1)$, $(0,1)$, $(1,2)$; (b) en $(0,1)$ y $(1,2)$; (c) en $(0,1)$.
\comprobacion valores numéricos de las raíces: (a) $-1{,}879$; $0{,}347$; $1{,}532$; (b) $0{,}619$ y $1{,}512$; (c) $0{,}659$.
\end{solucion}

\begin{solucion}{eje:3-ed-continua}
\datos tres funciones.
\planteamiento \cref{def:continuidad}: dado $\varepsilon$, encontrar $\delta$ (que puede depender de $a$).
\pasos
\begin{enumerate}[label=(\alph*)]
  \item $\abs{(3x-7)-(3a-7)}=3\abs{x-a}$. Dado $\varepsilon>0$, $\delta=\varepsilon/3$.
  \item $\abs{x^{2}-a^{2}}=\abs{x-a}\abs{x+a}$, y $\abs{x+a}\le\abs{x-a}+2\abs a<1+2\abs a$ si $\abs{x-a}<1$. Dado $\varepsilon>0$, sea $\delta=\min\bigl(1,\frac{\varepsilon}{1+2\abs a}\bigr)$; si $\abs{x-a}<\delta$, $\abs{x^{2}-a^{2}}<(1+2\abs a)\delta\le\varepsilon$.
  \item Si $a>0$: para $x\ge0$, $\abs{\sqrt x-\sqrt a}=\frac{\abs{x-a}}{\sqrt x+\sqrt a}\le\frac{\abs{x-a}}{\sqrt a}$; dado $\varepsilon>0$, $\delta=\varepsilon\sqrt a$ sirve. Si $a=0$: para $x\ge0$, $\sqrt x<\varepsilon\iff x<\varepsilon^{2}$; $\delta=\varepsilon^{2}$ sirve. \hfill$\blacksquare$
\end{enumerate}
\resultado $\delta=\varepsilon/3$; $\delta=\min(1,\varepsilon/(1+2\abs a))$; $\delta=\varepsilon\sqrt a$ o $\delta=\varepsilon^{2}$.
\comprobacion en (b), $\delta$ disminuye al crecer $\abs a$: la parábola es más empinada lejos del origen, y el mismo $\varepsilon$ exige un $\delta$ menor. Ningún $\delta$ sirve a la vez para todos los $a$.
\end{solucion}

\begin{solucion}{eje:3-punto-fijo}
\datos $f\colon[0,1]\to[0,1]$ continua.
\planteamiento aplicar Bolzano a $g(x)=f(x)-x$.
\pasos
\begin{enumerate}[label=\arabic*.]
  \item $g$ es continua en $[0,1]$, $g(0)=f(0)\ge0$ y $g(1)=f(1)-1\le0$, porque $f$ toma valores en $[0,1]$.
  \item Si $g(0)=0$ o $g(1)=0$, el punto fijo es $0$ o $1$. Si no, $g(0)>0>g(1)$ y, por Bolzano (aplicado a $-g$), existe $c\in(0,1)$ con $g(c)=0$, es decir, $f(c)=c$.
  \item (b) $f(x)=x^{2}$ es continua de $(0,1)$ en $(0,1)$ y $x^{2}<x$ en $(0,1)$: no hay punto fijo. $f(x)=1$ si $x<\frac12$ y $f(x)=0$ si $x\ge\frac12$ va de $[0,1]$ en $[0,1]$ y no tiene punto fijo: $f(x)=x$ exigiría $x=1<\frac12$ o $x=0\ge\frac12$. \hfill$\blacksquare$
\end{enumerate}
\resultado existe $c$ con $f(c)=c$; los ejemplos de (b) muestran que hacen falta el intervalo cerrado y la continuidad.
\comprobacion con $f=\cos$, que lleva $[0,1]$ en $[\cos1,1]\subseteq[0,1]$, el punto fijo es $c\approx0{,}739$.
\end{solucion}

\begin{solucion}{eje:3-cuerda}
\datos posiciones $s_1(t)$ (subida) y $s_2(t)$ (bajada), medidas como distancia a la base a lo largo del camino, continuas en $t\in[8,14]$; longitud del camino $D>0$.
\planteamiento Bolzano aplicado a $h=s_1-s_2$.
\pasos
\begin{enumerate}[label=\arabic*.]
  \item $s_1(8)=0$, $s_1(14)=D$, $s_2(8)=D$, $s_2(14)=0$.
  \item $h=s_1-s_2$ es continua, $h(8)=-D<0$ y $h(14)=D>0$.
  \item Por Bolzano, existe $t_0\in(8,14)$ con $h(t_0)=0$: a la hora $t_0$, el ciclista está el primer día y el segundo en el mismo punto del camino. \hfill$\blacksquare$
\end{enumerate}
\resultado existe tal hora y tal punto.
\comprobacion el argumento no dice cuál es la hora ni que sea única; si un día se detiene a descansar, puede haber varias.
\end{solucion}

\begin{solucion}{eje:3-sucesiones}
\datos $f(x_n)\to f(y)$ para toda $(x_n)\subseteq[a,b]$ con $x_n\to y$.
\planteamiento contrarrecíproco: si $f$ no es continua en $y$, construir una sucesión mala.
\pasos
\begin{enumerate}[label=\arabic*.]
  \item Supóngase que $f$ no es continua en $y$. Negando la \cref{def:continuidad}: existe $\varepsilon>0$ tal que para todo $\delta>0$ hay $x\in[a,b]$ con $\abs{x-y}<\delta$ y $\abs{f(x)-f(y)}\ge\varepsilon$.
  \item Para cada $n\in\N$, con $\delta=1/n$, se elige $x_n\in[a,b]$ con $\abs{x_n-y}<1/n$ y $\abs{f(x_n)-f(y)}\ge\varepsilon$.
  \item Entonces $x_n\to y$, pero $f(x_n)\not\to f(y)$, en contra de la hipótesis. Luego $f$ es continua en $y$. \hfill$\blacksquare$
\end{enumerate}
\resultado $f$ es continua en $y$.
\comprobacion el recíproco (continua implica que conserva los límites de sucesiones) es la parte ($\Rightarrow$) del \cref{teo:heine}; juntos dan la caracterización por sucesiones de la continuidad.
\end{solucion}

\begin{solucion}{eje:3-inyectiva-monotona}
\datos $f$ continua e inyectiva en $[a,b]$, $f(a)<f(b)$.
\planteamiento (a) reducción al absurdo con el TVI; (b) aplicar (a) a subintervalos.
\pasos
\begin{enumerate}[label=\arabic*.]
  \item (a) Sea $x\in(a,b)$. Por la inyectividad, $f(x)\ne f(a)$ y $f(x)\ne f(b)$.
  \item Si fuera $f(x)>f(b)$, entonces $f(a)<f(b)<f(x)$ y el \cref{cor:tvi} en $[a,x]$ daría $c\in(a,x)$ con $f(c)=f(b)$; como $c<x<b$, $c\ne b$, en contra de la inyectividad.
  \item Si fuera $f(x)<f(a)$, entonces $f(x)<f(a)<f(b)$ y el TVI en $[x,b]$ daría $c\in(x,b)$ con $f(c)=f(a)$, $c\ne a$, de nuevo contradicción. Luego $f(a)<f(x)<f(b)$.
  \item (b) Sean $a\le x<y\le b$. Por (a) (o por la hipótesis si $x=a$), $f(x)<f(b)$ cuando $x<b$. La restricción de $f$ a $[x,b]$ es continua, inyectiva y cumple $f(x)<f(b)$; si $y<b$, (a) aplicado a ella da $f(x)<f(y)$, y si $y=b$ ya se sabe. Luego $f$ es estrictamente creciente. \hfill$\blacksquare$
\end{enumerate}
\resultado (a) $f(a)<f(x)<f(b)$ en $(a,b)$; (b) $f$ es estrictamente creciente.
\comprobacion la continuidad es imprescindible: la función inyectiva y no monótona del \cref{eje:1-estricta} tiene un salto.
\end{solucion}

\begin{solucion}{eje:3-cauchy}
\datos $g(x+y)=g(x)+g(y)$ para todos $x,y$.
\planteamiento (a) sustituciones particulares e inducción; (b) densidad de $\Q$ y continuidad; (c) trasladar la continuidad con la ecuación.
\pasos
\begin{enumerate}[label=\arabic*.]
  \item Con $x=y=0$: $g(0)=2g(0)$, luego $g(0)=0$. Con $y=-x$: $0=g(0)=g(x)+g(-x)$, luego $g(-x)=-g(x)$.
  \item $g(1\cdot x)=g(x)$, y si $g(nx)=ng(x)$, entonces $g((n+1)x)=g(nx)+g(x)=(n+1)g(x)$. Por inducción, vale para todo $n\in\N$.
  \item Sea $q=m/n$ con $m\in\Z$, $n\in\N$, y $c=g(1)$. Por el paso~2 (y la imparidad si $m<0$), $ng(m/n)=g(m)=mg(1)=mc$, luego $g(q)=cq$.
  \item (b) Sea $x\in\R$ y $(q_k)$ una sucesión de racionales con $q_k\to x$ (por ejemplo, los truncamientos decimales de $x$). Por la continuidad, $g(x)=\lim g(q_k)=\lim cq_k=cx$.
  \item (c) Sea $g$ continua en $x_0$ y sea $x\in\R$. Para todo $h$, $g(x+h)-g(x)=g(h)=g(x_0+h)-g(x_0)$, y el último miembro tiende a~0 cuando $h\to0$. Luego $g$ es continua en $x$. \hfill$\blacksquare$
\end{enumerate}
\resultado $g(q)=g(1)q$ en $\Q$; si $g$ es continua (basta en un punto), $g(x)=cx$ en $\R$.
\comprobacion $g(x)=cx$ cumple la ecuación: $c(x+y)=cx+cy$. Sin ninguna hipótesis de continuidad existen soluciones no lineales, pero su construcción requiere el axioma de elección.
\end{solucion}

\begin{solucion}{eje:3-exponencial-funcional}
\datos $h$ continua con $h(x+y)=h(x)h(y)$.
\planteamiento (a) escribir $h(x)$ como un cuadrado; (b) tomar logaritmos y aplicar el \cref{eje:3-cauchy}.
\pasos
\begin{enumerate}[label=\arabic*.]
  \item (a) $h(x)=h(\frac x2+\frac x2)=h(\frac x2)^{2}\ge0$ para todo $x$.
  \item Si $h(x_0)=0$ para algún $x_0$, entonces $h(x)=h(x_0)h(x-x_0)=0$ para todo $x$. Luego, o $h\equiv0$, o $h(x)\ne0$ para todo $x$, y entonces $h(x)>0$ por el paso~1.
  \item (b) Si $h>0$, $g=\ln h$ es continua (composición) y $g(x+y)=\ln\bigl(h(x)h(y)\bigr)=g(x)+g(y)$.
  \item Por el \cref{eje:3-cauchy}(b), $g(x)=cx$ con $c=g(1)$, luego $h(x)=e^{cx}=a^{x}$ con $a=e^{c}=h(1)>0$. \hfill$\blacksquare$
\end{enumerate}
\resultado $h\equiv0$ o $h(x)=a^{x}$ con $a=h(1)>0$.
\comprobacion $a^{x+y}=a^{x}a^{y}$; y $h(0)=h(0)^{2}$ da $h(0)\in\{0,1\}$, coherente con $a^{0}=1$.
\end{solucion}

\begin{solucion}{eje:3-periodica}
\datos $f$ continua en $\R$, $f(x+T)=f(x)$.
\planteamiento (a) reducir a $[0,T]$ y aplicar Weierstrass; (b) Bolzano a $g$.
\pasos
\begin{enumerate}[label=\arabic*.]
  \item (a) Por inducción, $f(x+kT)=f(x)$ para $k\in\N$, y también para $k$ negativo, porque $f(x)=f((x-T)+T)=f(x-T)$.
  \item Todo $x\in\R$ se escribe $x=kT+r$ con $k=\floor{x/T}\in\Z$ y $r\in[0,T)$, y $f(x)=f(r)$. Luego $f(\R)=f([0,T])$.
  \item Por Weierstrass, $f$ alcanza en $[0,T]$ un máximo y un mínimo, que lo son también en $\R$ por el paso~2.
  \item (b) Sea $g(x)=f(x+T/2)-f(x)$, continua. $g(0)=f(T/2)-f(0)$ y $g(T/2)=f(T)-f(T/2)=f(0)-f(T/2)=-g(0)$.
  \item Si $g(0)=0$, sirve $x=0$. Si no, $g(0)$ y $g(T/2)$ tienen signos opuestos y Bolzano da $x\in(0,T/2)$ con $g(x)=0$. \hfill$\blacksquare$
\end{enumerate}
\resultado $f$ alcanza máximo y mínimo; existe $x$ con $f(x+T/2)=f(x)$.
\comprobacion con $f=\sin$, $T=2\pi$: $\sin(x+\pi)=\sin x$ equivale a $-\sin x=\sin x$, que se cumple en $x=0$. Interpretación clásica: en todo momento hay dos puntos antípodas del ecuador con la misma temperatura.
\end{solucion}
% ---------------------------------------------------------------------
\section{Tema 4}

\begin{solucion}{eje:4-grafica-derivada}
\datos cuatro funciones; la derivada es la pendiente de la tangente.
\planteamiento leer la pendiente en cada tramo de la gráfica.
\pasos
\begin{enumerate}[label=(\alph*)]
  \item Constante: pendiente nula; $f'=0$ (el eje de abscisas).
  \item $\ReLU$: pendiente $0$ para $x<0$ y $1$ para $x>0$; $\ReLU'$ es la función escalón, no definida en~0.
  \item $\abs x$: pendiente $-1$ para $x<0$ y $1$ para $x>0$; la derivada es la función signo, no definida en~0.
  \item $\sigma$: siempre creciente, con la máxima pendiente en el centro de simetría; $\sigma'$ es una «campana» positiva, par, con máximo $\sigma'(0)=\frac14$, que tiende a~0 en $\pm\infty$.
  \item Paridad: si la gráfica es simétrica respecto del eje de ordenadas, la tangente en $-x$ es la reflejada de la tangente en $x$, y una reflexión vertical cambia el signo de la pendiente: $f'(-x)=-f'(x)$.
\end{enumerate}
\resultado (a) $0$; (b) escalón; (c) signo; (d) campana con máximo $\frac14$ en $x=0$; la derivada de una función par es impar.
\comprobacion (c) es la derivada de una función par ($\abs x$) y es impar; (d) es la derivada de $\sigma-\frac12$, impar, y es par (\cref{eje:4-par-impar}).
\end{solucion}

\begin{solucion}{eje:4-definicion}
\datos $f(x)=x^{3}$, $\sqrt x$ y $\frac{1}{2x+1}$.
\planteamiento cociente incremental, simplificación del factor $h$ y límite.
\pasos
\begin{enumerate}[label=(\alph*)]
  \item $\frac{(a+h)^{3}-a^{3}}{h}=\frac{3a^{2}h+3ah^{2}+h^{3}}{h}=3a^{2}+3ah+h^{2}\to3a^{2}$.
  \item $\frac{\sqrt{a+h}-\sqrt a}{h}=\frac{(a+h)-a}{h(\sqrt{a+h}+\sqrt a)}=\frac{1}{\sqrt{a+h}+\sqrt a}\to\frac{1}{2\sqrt a}$, por la continuidad de la raíz.
  \item $f(1)=\frac13$ y $\frac{1}{h}\Bigl(\frac{1}{3+2h}-\frac13\Bigr)=\frac{3-(3+2h)}{3h(3+2h)}=\frac{-2}{3(3+2h)}\to-\frac29$. Tangente: $y=\frac13-\frac29(x-1)$, es decir, $y=\frac{5-2x}{9}$.
\end{enumerate}
\resultado (a) $3a^{2}$; (b) $\frac{1}{2\sqrt a}$; (c) $f'(1)=-\frac29$, tangente $y=\frac{5-2x}{9}$.
\comprobacion (c) con la regla de la cadena: $f'(x)=-\frac{2}{(2x+1)^{2}}$, y $f'(1)=-\frac29$.
\end{solucion}

\begin{solucion}{eje:4-mecanico-1}
\datos seis funciones.
\planteamiento reglas de la suma, del producto, del cociente y de la cadena (\cref{teo:reglas,teo:cadena}).
\pasos
\begin{enumerate}[label=(\alph*)]
  \item $10x^{9}+15x^{4}$ ($\pi^{3}$ es constante).
  \item $(3x^{2}+1)(x^{5}+x^{2})+(x^{3}+x)(5x^{4}+2x)=8x^{7}+6x^{5}+5x^{4}+3x^{2}$.
  \item $\frac{(x^{2}-1)-(x+2)\cdot2x}{(x^{2}-1)^{2}}=-\frac{x^{2}+4x+1}{(x^{2}-1)^{2}}$.
  \item $\frac{x^{4}+1}{x}=x^{3}+x^{-1}$, luego $3x^{2}-\frac{1}{x^{2}}$.
  \item $100(x^{2}+2)^{99}\cdot2x=200x(x^{2}+2)^{99}$.
  \item $10x^{9}(x^{2}+1)^{10}+x^{10}\cdot10(x^{2}+1)^{9}\cdot2x=10x^{9}(x^{2}+1)^{9}(3x^{2}+1)$.
\end{enumerate}
\resultado los indicados.
\comprobacion (b) multiplicando primero: $(x^{3}+x)(x^{5}+x^{2})=x^{8}+x^{6}+x^{5}+x^{3}$, cuya derivada es $8x^{7}+6x^{5}+5x^{4}+3x^{2}$.
\end{solucion}

\begin{solucion}{eje:4-mecanico-2}
\datos ocho funciones.
\planteamiento regla de la cadena, \cref{prop:tabla,prop:derivadas-activaciones}.
\pasos
\begin{enumerate}[label=(\alph*)]
  \item $e^{-x^{2}/2}\cdot(-x)=-xe^{-x^{2}/2}$.
  \item $\frac{-\sin x}{\cos x}=-\tan x$.
  \item $2\sin3x\cdot\cos3x\cdot3=3\sin6x$.
  \item $\frac{1}{1+x}\cdot\frac{1}{2\sqrt x}=\frac{1}{2\sqrt x(1+x)}$.
  \item $\frac{1}{x+\sqrt{x^{2}+1}}\Bigl(1+\frac{x}{\sqrt{x^{2}+1}}\Bigr)=\frac{1}{x+\sqrt{x^{2}+1}}\cdot\frac{\sqrt{x^{2}+1}+x}{\sqrt{x^{2}+1}}=\frac{1}{\sqrt{x^{2}+1}}$.
  \item Con $z=wx+b$: $\frac{\partial}{\partial x}\sigma(z)=w\,\sigma(z)(1-\sigma(z))$ y $\frac{\partial}{\partial w}\sigma(z)=x\,\sigma(z)(1-\sigma(z))$.
  \item $2x\bigl(1-\tanh^{2}(x^{2})\bigr)$.
  \item $\frac{\sigma'(x)}{\sigma(x)}=1-\sigma(x)=\sigma(-x)$.
\end{enumerate}
\resultado los indicados.
\comprobacion (e) es la derivada de $\arsinh$ (\cref{eje:1-arsinh}), y coincide con la que da el \cref{teo:inversa}: $1/\cosh(\arsinh x)=1/\sqrt{1+\sinh^{2}(\arsinh x)}=1/\sqrt{1+x^{2}}$. (h) es coherente con $\ln\sigma(x)=-\softplus(-x)$ y $\softplus'=\sigma$.
\end{solucion}

\begin{solucion}{eje:4-logaritmica}
\datos tres funciones positivas.
\planteamiento derivar $\ln f$ y multiplicar por $f$ (\cref{ej:logaritmica}).
\pasos
\begin{enumerate}[label=(\alph*)]
  \item $\ln f=\sin x\ln x$; $\frac{f'}{f}=\cos x\ln x+\frac{\sin x}{x}$; $f'=x^{\sin x}\bigl(\cos x\ln x+\frac{\sin x}{x}\bigr)$.
  \item $\ln f=x\bigl(\ln(x+1)-\ln x\bigr)$; $\frac{f'}{f}=\ln\frac{x+1}{x}+x\Bigl(\frac{1}{x+1}-\frac1x\Bigr)=\ln\Bigl(1+\frac1x\Bigr)-\frac{1}{x+1}$.
  \item $\ln f=3\ln(x+1)+\frac12\ln(x^{2}+4)-5\ln(2x+3)$; $\frac{f'}{f}=\frac{3}{x+1}+\frac{x}{x^{2}+4}-\frac{10}{2x+3}$, y $f'=f\cdot\frac{f'}{f}$.
\end{enumerate}
\resultado (a) $x^{\sin x}\bigl(\cos x\ln x+\frac{\sin x}{x}\bigr)$; (b) $\bigl(1+\frac1x\bigr)^{x}\bigl(\ln(1+\frac1x)-\frac{1}{x+1}\bigr)$; (c) $\frac{(x+1)^{3}\sqrt{x^{2}+4}}{(2x+3)^{5}}\bigl(\frac{3}{x+1}+\frac{x}{x^{2}+4}-\frac{10}{2x+3}\bigr)$.
\comprobacion (b): por la \cref{prop:notables-exp}, $(1+1/x)^{x}\to e$; la derivada debe tender a~0, y en efecto $\ln(1+\frac1x)-\frac{1}{x+1}=\frac1x-\frac{1}{x+1}+O(x^{-2})\to0$.
\end{solucion}

\begin{solucion}{eje:4-implicita}
\datos (a) $x^{2/3}+y^{2/3}=4$, punto $(-\sqrt{27},1)$; (b) $\sin x+\sin y=\frac12$.
\planteamiento derivación implícita y despeje de $y'$.
\pasos
\begin{enumerate}[label=\arabic*.]
  \item (a) El punto está en la curva: $(-\sqrt{27})^{2/3}=27^{1/3}=3$ y $1^{2/3}=1$.
  \item Derivando: $\frac23x^{-1/3}+\frac23y^{-1/3}y'=0$, luego $y'=-\bigl(\frac yx\bigr)^{1/3}$.
  \item En el punto: $\frac yx=-\frac{1}{\sqrt{27}}=-3^{-3/2}$, su raíz cúbica es $-3^{-1/2}$, y $y'=3^{-1/2}=\frac{1}{\sqrt3}$.
  \item Tangente: $y=1+\frac{1}{\sqrt3}(x+\sqrt{27})=\frac{x}{\sqrt3}+4$, porque $\sqrt{27}/\sqrt3=3$.
  \item (b) Derivando: $\cos x+\cos y\cdot y'=0$, $y'=-\frac{\cos x}{\cos y}$ donde $\cos y\ne0$. Tangente horizontal: $\cos x=0$, es decir, $x=\frac\pi2+k\pi$, con $\sin x=\pm1$.
  \item Si $\sin x=1$: $\sin y=-\frac12$, $y=-\frac\pi6+2m\pi$ o $y=\frac{7\pi}{6}+2m\pi$, y ahí $\cos y=\pm\frac{\sqrt3}{2}\ne0$. Si $\sin x=-1$: $\sin y=\frac32$, imposible.
\end{enumerate}
\resultado (a) $y=\frac{x}{\sqrt3}+4$; (b) los puntos $\bigl(\frac\pi2+2k\pi,\,-\frac\pi6+2m\pi\bigr)$ y $\bigl(\frac\pi2+2k\pi,\,\frac{7\pi}{6}+2m\pi\bigr)$, con $k,m\in\Z$.
\comprobacion (a): la tangente pasa por el punto, $\frac{-\sqrt{27}}{\sqrt3}+4=-3+4=1$. (b): $\sin\frac\pi2+\sin(-\frac\pi6)=1-\frac12=\frac12$.
\end{solucion}

\begin{solucion}{eje:4-tangentes-origen}
\datos $f(x)=1+(x-1)^{2}$, $f'(x)=2(x-1)$.
\planteamiento tangente en un punto genérico $t$ e imposición de que pase por $(0,0)$.
\pasos
\begin{enumerate}[label=\arabic*.]
  \item Tangente en $x=t$: $y=1+(t-1)^{2}+2(t-1)(x-t)$.
  \item Pasa por el origen si $0=1+(t-1)^{2}-2t(t-1)=2-t^{2}$, es decir, $t=\pm\sqrt2$.
  \item En el origen la ordenada es $0$, luego la tangente es $y=f'(t)x=2(t-1)x$: para $t=\sqrt2$, $y=2(\sqrt2-1)x$; para $t=-\sqrt2$, $y=-2(\sqrt2+1)x$.
\end{enumerate}
\resultado $y=2(\sqrt2-1)x$ e $y=-2(\sqrt2+1)x$.
\comprobacion en $t=\sqrt2$: $f(\sqrt2)=1+(\sqrt2-1)^{2}=4-2\sqrt2$ y $2(\sqrt2-1)\sqrt2=4-2\sqrt2$.
\end{solucion}

\begin{solucion}{eje:4-pegado}
\datos $f$ definida a trozos con $x^{2}+4x+1$ y $ax+b$.
\planteamiento derivable implica continua (\cref{teo:derivable-continua}): igualar valores y derivadas laterales en el punto de unión.
\pasos
\begin{enumerate}[label=\arabic*.]
  \item (a) En $x=0$: continuidad, $b=0^{2}+0+1=1$; derivadas laterales, $a=2\cdot0+4=4$.
  \item Con $a=4$, $b=1$, ambos trozos tienen en~0 el mismo valor y la misma pendiente, luego $f'_{-}(0)=f'_{+}(0)=4$ y $f$ es derivable; en los demás puntos lo es por ser polinómica en cada trozo.
  \item (b) En $x=1$: continuidad, $a+b=1+4+1=6$; derivadas, $a=2+4=6$. Luego $a=6$, $b=0$.
\end{enumerate}
\resultado (a) $a=4$, $b=1$; (b) $a=6$, $b=0$.
\comprobacion en ambos casos $ax+b$ es la recta tangente a la parábola en el punto de unión: $y=1+4x$ en $x=0$ e $y=6+6(x-1)=6x$ en $x=1$.
\end{solucion}

\begin{solucion}{eje:4-sucesivas}
\datos $f(x)=xe^{x}$; $\sigma'=\sigma(1-\sigma)$; $\tanh'=1-\tanh^{2}$.
\planteamiento derivar repetidamente y detectar el patrón; inducción.
\pasos
\begin{enumerate}[label=\arabic*.]
  \item (a) $f'=(x+1)e^{x}$, $f''=(x+2)e^{x}$, y en general se conjetura $f^{(n)}=(x+n)e^{x}$. Si vale para $n$, $f^{(n+1)}=e^{x}+(x+n)e^{x}=(x+n+1)e^{x}$; vale para $n=1$, luego para todo $n$. En particular, $f^{(5)}(x)=(x+5)e^{x}$.
  \item (b) Con $s=\sigma(x)$: $\sigma''=\frac{\dd}{\dd x}(s-s^{2})=(1-2s)s'=s(1-s)(1-2s)=s-3s^{2}+2s^{3}$.
  \item $\sigma'''=(1-6s+6s^{2})s'$. En $x=0$, $s=\frac12$ y $s'=\frac14$: $\sigma'''(0)=(1-3+\frac32)\cdot\frac14=-\frac18$.
  \item (c) Con $t=\tanh x$: $\tanh''=\frac{\dd}{\dd x}(1-t^{2})=-2tt'=-2t(1-t^{2})$.
\end{enumerate}
\resultado (a) $f^{(n)}=(x+n)e^{x}$; (b) $\sigma''=\sigma(1-\sigma)(1-2\sigma)$, $\sigma'''(0)=-\frac18$; (c) $\tanh''=-2t(1-t^{2})$.
\comprobacion por el \cref{eje:1-sigmoide-simetria}, $\sigma(x)=\frac12+\frac12\tanh\frac x2$, luego $\sigma'''(0)=\frac12\cdot\frac18\tanh'''(0)$; y $\tanh'''=-2(1-t^{2})(1-3t^{2})$ vale $-2$ en~0, lo que da $-\frac18$.
\end{solucion}

\begin{solucion}{eje:4-producto}
\datos $f$, $g$ derivables en $a$.
\planteamiento sumar y restar $f(a+h)g(a)$; para Leibniz, inducción sobre $n$ con la fórmula de Pascal.
\pasos
\begin{enumerate}[label=\arabic*.]
  \item \emph{Enunciado.} Si $f$ y $g$ son derivables en $a$, $fg$ es derivable en $a$ y $(fg)'(a)=f'(a)g(a)+f(a)g'(a)$.
  \item Para $h\ne0$ pequeño, $\frac{f(a+h)g(a+h)-f(a)g(a)}{h}=f(a+h)\frac{g(a+h)-g(a)}{h}+g(a)\frac{f(a+h)-f(a)}{h}$ (se ha sumado y restado $f(a+h)g(a)$).
  \item Cuando $h\to0$: $\frac{g(a+h)-g(a)}{h}\to g'(a)$ y $\frac{f(a+h)-f(a)}{h}\to f'(a)$ por definición; $f(a+h)\to f(a)$ porque $f$, derivable en $a$, es continua en $a$ (\cref{teo:derivable-continua}). Este es el único lugar donde se usa la continuidad.
  \item Por el álgebra de límites, el cociente tiende a $f(a)g'(a)+g(a)f'(a)$.
  \item \emph{Leibniz.} Para $n=1$ es la regla del producto. Supóngase cierta para $n$ (con $f,g$ derivables $n+1$ veces). Derivando cada término con la regla del producto,
  \[
  (fg)^{(n+1)}=\sum_{k=0}^{n}\binom nk\bigl(f^{(k+1)}g^{(n-k)}+f^{(k)}g^{(n+1-k)}\bigr)=\sum_{j=0}^{n+1}\Bigl[\binom n{j-1}+\binom nj\Bigr]f^{(j)}g^{(n+1-j)},
  \]
  con el convenio $\binom n{-1}=\binom n{n+1}=0$, tras escribir $j=k+1$ en la primera suma y $j=k$ en la segunda.
  \item Por la fórmula de Pascal, $\binom n{j-1}+\binom nj=\binom{n+1}{j}$, que es la fórmula para $n+1$. Para $n=2$ resulta $f''g+2f'g'+fg''$. \hfill$\blacksquare$
\end{enumerate}
\resultado la regla del producto y la fórmula de Leibniz.
\comprobacion con $f=x$, $g=e^{x}$: $(xe^{x})^{(n)}=\binom n0xe^{x}+\binom n1\cdot1\cdot e^{x}=(x+n)e^{x}$, como en el \cref{eje:4-sucesivas}.
\end{solucion}

\begin{solucion}{eje:4-par-impar}
\datos $f$ derivable y par (o impar).
\planteamiento derivar la identidad $f(-x)=\pm f(x)$ con la regla de la cadena.
\pasos
\begin{enumerate}[label=\arabic*.]
  \item Si $f$ es par, $f(-x)=f(x)$ para todo $x$. Ambos miembros son funciones derivables de $x$; derivando el izquierdo con la regla de la cadena, $-f'(-x)=f'(x)$, luego $f'(-x)=-f'(x)$: $f'$ es impar.
  \item Si $f$ es impar, $f(-x)=-f(x)$; derivando, $-f'(-x)=-f'(x)$, luego $f'(-x)=f'(x)$: $f'$ es par.
  \item Comprobaciones: $\cosh$ es par y $\cosh'=\sinh$ es impar; $\sinh$ es impar y $\sinh'=\cosh$ es par; $\sigma-\frac12$ es impar (\cref{prop:sigmoide}(ii)) y su derivada $\sigma'(x)=\sigma(x)\sigma(-x)$ es par. \hfill$\blacksquare$
\end{enumerate}
\resultado par $\Rightarrow$ derivada impar; impar $\Rightarrow$ derivada par.
\comprobacion consecuencia: si $f$ es par y derivable en~0, $f'(0)=-f'(0)$, luego $f'(0)=0$ (el vértice de una función par derivable es un punto crítico).
\end{solucion}

\begin{solucion}{eje:4-funcional}
\datos $f(x+y)=f(x)+f(y)+x^{2}y+xy^{2}$; $f(x)/x\to1$ cuando $x\to0$.
\planteamiento sustituir valores particulares y escribir el cociente incremental con la ecuación.
\pasos
\begin{enumerate}[label=\arabic*.]
  \item Con $x=y=0$: $f(0)=2f(0)$, luego $f(0)=0$.
  \item $f'(0)=\lim_{h\to0}\frac{f(h)-f(0)}{h}=\lim\frac{f(h)}{h}=1$.
  \item Sea $x\in\R$. Con $y=h$: $f(x+h)-f(x)=f(h)+x^{2}h+xh^{2}$, luego $\frac{f(x+h)-f(x)}{h}=\frac{f(h)}{h}+x^{2}+xh\to1+x^{2}$.
  \item Por tanto $f$ es derivable en todo punto, con $f'(x)=1+x^{2}$. \hfill$\blacksquare$
\end{enumerate}
\resultado $f(0)=0$, $f'(0)=1$, $f'(x)=1+x^{2}$.
\comprobacion la función $f(x)=x+\frac{x^{3}}{3}$ cumple la ecuación: $\frac{(x+y)^{3}-x^{3}-y^{3}}{3}=x^{2}y+xy^{2}$. Por el \cref{cor:tvm}(ii), es la única solución, pues cualquier otra difiere de ella en una constante, que es nula porque $f(0)=0$.
\end{solucion}

\begin{solucion}{eje:4-raiz-doble}
\datos $p$ polinomio, $a\in\R$.
\planteamiento teorema del factor: $p(a)=0$ si y solo si $(x-a)$ divide a $p$.
\pasos
\begin{enumerate}[label=\arabic*.]
  \item ($\Rightarrow$) Si $p=(x-a)^{2}q$, entonces $p(a)=0$ y $p'=2(x-a)q+(x-a)^{2}q'$, luego $p'(a)=0$.
  \item ($\Leftarrow$) Si $p(a)=0$, $p=(x-a)q$ con $q$ polinomio. Derivando, $p'=q+(x-a)q'$, y $p'(a)=q(a)$. Si además $p'(a)=0$, entonces $q(a)=0$ y $q=(x-a)r$, de modo que $p=(x-a)^{2}r$.
  \item Aplicación: $p(x)=x^{3}-3x+c$, $p'(x)=3x^{2}-3=0\iff x=\pm1$. $p(1)=c-2=0$ da $c=2$; $p(-1)=c+2=0$ da $c=-2$. \hfill$\blacksquare$
\end{enumerate}
\resultado $a$ es raíz doble si y solo si $p(a)=p'(a)=0$; $c=2$ (raíz doble $1$) o $c=-2$ (raíz doble $-1$).
\comprobacion $x^{3}-3x+2=(x-1)^{2}(x+2)$ y $x^{3}-3x-2=(x+1)^{2}(x-2)$.
\end{solucion}

\begin{solucion}{eje:4-x2sen}
\datos $f(x)=x^{2}\sin(1/x)$, $g(x)=x\sin(1/x)$, ambas nulas en~0.
\planteamiento fuera de~0, reglas de derivación; en~0, la definición; para la discontinuidad de $f'$, una sucesión.
\pasos
\begin{enumerate}[label=\arabic*.]
  \item (a) Para $x\ne0$: $f'(x)=2x\sin\frac1x+x^{2}\cos\frac1x\cdot\bigl(-\frac{1}{x^{2}}\bigr)=2x\sin\frac1x-\cos\frac1x$.
  \item En~0: $\frac{f(h)-f(0)}{h}=h\sin\frac1h$, y $\abs{h\sin\frac1h}\le\abs h\to0$. Luego $f'(0)=0$ y $f$ es derivable en todo $\R$.
  \item (b) Sea $x_n=\frac{1}{2\pi n}\to0$. Entonces $f'(x_n)=0-\cos(2\pi n)=-1$, que no tiende a $f'(0)=0$. Por el \cref{teo:heine}, $f'$ no es continua en~0.
  \item (c) $\abs{g(x)}\le\abs x\to0=g(0)$: $g$ es continua en~0. Pero $\frac{g(h)-g(0)}{h}=\sin\frac1h$, que no tiene límite en~0 (\cref{ej:sin1x}): $g$ no es derivable en~0. \hfill$\blacksquare$
\end{enumerate}
\resultado $f$ es derivable en $\R$ con $f'(0)=0$, pero $f'$ es discontinua en~0; $g$ es continua y no derivable en~0.
\comprobacion con $y_n=\frac{1}{2\pi n+\pi}$, $f'(y_n)=-\cos(2\pi n+\pi)=1$: $f'$ toma los valores $\pm1$ arbitrariamente cerca de~0, luego ni siquiera tiene límite allí. Aun así, $f'$ cumple la propiedad de los valores intermedios (\cref{eje:5-darboux}).
\end{solucion}

\begin{solucion}{eje:4-cadena-n}
\datos $F=f_n\circ\cdots\circ f_1$, $x_k=f_k(x_{k-1})$.
\planteamiento inducción sobre $n$ con el \cref{teo:cadena}; cota con la \cref{prop:derivadas-activaciones}.
\pasos
\begin{enumerate}[label=\arabic*.]
  \item Sea $F_n=f_n\circ\cdots\circ f_1$. Para $n=1$, $F_1'(x)=f_1'(x_0)$.
  \item Supóngase la fórmula para $F_n$. Como $F_{n+1}=f_{n+1}\circ F_n$ y $F_n(x)=x_n$, la regla de la cadena da $F_{n+1}'(x)=f_{n+1}'(x_n)\,F_n'(x)=f_{n+1}'(x_n)\prod_{k=1}^{n}f_k'(x_{k-1})=\prod_{k=1}^{n+1}f_k'(x_{k-1})$.
  \item Si todas las $f_k$ son $\sigma$, cada factor es $\sigma'(x_{k-1})\in(0,\frac14]$, luego $0<F'(x)\le(\frac14)^{n}=4^{-n}$. \hfill$\blacksquare$
\end{enumerate}
\resultado $F'(x)=\prod_kf_k'(x_{k-1})$; para $n$ sigmoides, $0<F'\le4^{-n}$.
\comprobacion e interpretación: con $n=10$, $F'\le4^{-10}\approx9{,}5\cdot10^{-7}$. En una red profunda, el gradiente respecto de los parámetros de las primeras capas contiene un producto así y se hace minúsculo: es el desvanecimiento del gradiente que motiva las activaciones ReLU y las conexiones residuales (módulo~11).
\end{solucion}
% ---------------------------------------------------------------------
\section{Tema 5}

\begin{solucion}{eje:5-hipotesis}
\datos cuatro funciones en intervalos.
\planteamiento comprobar continuidad en $[a,b]$ y derivabilidad en $(a,b)$; calcular la pendiente de la cuerda y resolver $f'(c)=m$.
\pasos
\begin{enumerate}[label=(\alph*)]
  \item No es derivable en $\frac12\in(-1,1)$: \textbf{no} se cumplen las hipótesis. Cuerda: $m=\frac{\frac12-\frac32}{2}=-\frac12$; pero $f'=\pm1$ donde existe. \textbf{No hay $c$}.
  \item Discontinua en $0\in[-1,1]$: \textbf{no} se cumplen. $m=\frac{1-1}{2}=0$ y $f'(x)=-2/x^{3}\ne0$. \textbf{No hay $c$}.
  \item \textbf{Sí} se cumplen. $m=e-1$ y $e^{c}=e-1$ da $c=\ln(e-1)\approx0{,}541\in(0,1)$.
  \item Continua, pero no derivable en $0$ (tangente vertical): \textbf{no} se cumplen. Sin embargo, $m=\frac{1-(-1)}{2}=1$ y $\frac13c^{-2/3}=1$ da $c=\pm3^{-3/2}\approx\pm0{,}192$: \textbf{sí hay} $c$.
\end{enumerate}
\resultado hipótesis solo en (c); hay $c$ en (c) y en (d), no en (a) ni en (b).
\comprobacion (d) muestra que las hipótesis son suficientes, no necesarias: el teorema no dice nada cuando fallan, y la conclusión puede cumplirse o no.
\end{solucion}

\begin{solucion}{eje:5-valores-c}
\datos cuatro funciones derivables en sus intervalos.
\planteamiento resolver $f'(c)=\frac{f(b)-f(a)}{b-a}$ y quedarse con las soluciones en $(a,b)$.
\pasos
\begin{enumerate}[label=(\alph*)]
  \item $m=1$, $2c=1$, $c=\frac12$.
  \item $m=\ln2$, $\frac1c=\ln2$, $c=\frac{1}{\ln2}\approx1{,}443\in(1,2)$.
  \item $m=\frac{6-(-6)}{4}=3$, $3c^{2}-1=3$, $c=\pm\frac{2}{\sqrt3}\approx\pm1{,}155$; ambos en $(-2,2)$.
  \item $m=\frac82=4$, $3c^{2}=4$, $c=\frac{2}{\sqrt3}\approx1{,}155$ (el valor negativo no está en $(0,2)$).
\end{enumerate}
\resultado (a) $\frac12$; (b) $\frac{1}{\ln2}$; (c) $\pm\frac{2}{\sqrt3}$; (d) $\frac{2}{\sqrt3}$.
\comprobacion (b): $\ln2\approx0{,}693$ y $1/1{,}443\approx0{,}693$.
\end{solucion}

\begin{solucion}{eje:5-lhopital}
\datos ocho límites indeterminados.
\planteamiento identificar la forma, reducirla a $0/0$ o $\infty/\infty$ y aplicar el \cref{teo:lhopital} (varias veces si hace falta).
\pasos
\begin{enumerate}[label=(\alph*)]
  \item $0/0$: $\frac{e^{x}-1}{2x}\to\frac12$.
  \item $0/0$: $\frac{\cos x-\cos^{-2}x}{3x^{2}}$, de nuevo $0/0$: $\frac{-\sin x-2\cos^{-3}x\sin x}{6x}=-\frac{\sin x}{x}\cdot\frac{1+2\cos^{-3}x}{6}\to-\frac{3}{6}=-\frac12$.
  \item $0/0$: $\frac{1}{1/x}=x\to1$.
  \item $\infty/\infty$: $\frac{x^{2}}{e^{x}}\to\frac{2x}{e^{x}}\to\frac{2}{e^{x}}\to0$.
  \item $0^{0}$: $x^{2x}=e^{2x\ln x}$ y $x\ln x\to0$ (\cref{ej:lhopital}(c)); el límite es $e^{0}=1$.
  \item $1^{\infty}$: $(1-\frac1x)^{x}=\exp\bigl(x\ln(1-\frac1x)\bigr)$. Con $t=1/x\to0^{+}$, $\frac{\ln(1-t)}{t}\overset{0/0}{\to}\frac{-1/(1-t)}{1}\to-1$. El límite es $e^{-1}$.
  \item $\infty-\infty$: $\frac{x-\sin x}{x\sin x}\overset{0/0}{\to}\frac{1-\cos x}{\sin x+x\cos x}\overset{0/0}{\to}\frac{\sin x}{2\cos x-x\sin x}\to\frac02=0$.
  \item $0/0$: $\frac{3^{x}\ln3-2^{x}\ln2}{1}\to\ln3-\ln2=\ln\frac32$.
\end{enumerate}
\resultado (a) $\frac12$; (b) $-\frac12$; (c) $1$; (d) $0$; (e) $1$; (f) $e^{-1}$; (g) $0$; (h) $\ln\frac32\approx0{,}405$.
\comprobacion (a), (b), (g) con desarrollos de Taylor (\cref{tema:6}): $e^{x}-x-1=\frac{x^{2}}{2}+o(x^{2})$; $\sin x-\tan x=-\frac{x^{3}}{2}+o(x^{3})$; $\frac1{\sin x}-\frac1x=\frac x6+o(x)$.
\end{solucion}

\begin{solucion}{eje:5-grafica}
\datos (a) $f(x)=\frac{x^{2}}{x-1}$; (b) $g(x)=x-3x^{1/3}$; (c) $h(x)=e^{-x^{2}}$.
\planteamiento dominio, asíntotas (\cref{def:asintotas}), signo de la primera derivada (crecimiento y extremos) y de la segunda (convexidad).
\pasos
\begin{enumerate}[label=(\alph*)]
  \item Dominio $\R\setminus\{1\}$. $f(x)=x+1+\frac{1}{x-1}$: asíntota vertical $x=1$ y oblicua $y=x+1$. $f'(x)=\frac{x(x-2)}{(x-1)^{2}}$: $f$ crece en $(-\infty,0]$ y en $[2,\infty)$ y decrece en $[0,1)$ y en $(1,2]$; máximo local $f(0)=0$ y mínimo local $f(2)=4$. $f''(x)=\frac{2}{(x-1)^{3}}$: cóncava en $(-\infty,1)$ y convexa en $(1,\infty)$; sin inflexiones (el cambio se produce en la asíntota).
  \item Dominio $\R$; $g$ es impar. $g'(x)=1-x^{-2/3}$ para $x\ne0$, negativa si $0<\abs x<1$ y positiva si $\abs x>1$: crece en $(-\infty,-1]$ y $[1,\infty)$, decrece en $[-1,1]$; máximo local $g(-1)=2$, mínimo local $g(1)=-2$. En $0$, $g'(x)\to-\infty$: tangente vertical. $g''(x)=\frac23x^{-5/3}$: cóncava para $x<0$, convexa para $x>0$; inflexión en $(0,0)$. Ceros: $x=0$ y $x^{2}=27$, $x=\pm3\sqrt3\approx\pm5{,}20$. Sin asíntotas: $g(x)/x\to1$ pero $g(x)-x=-3x^{1/3}$ no tiene límite finito.
  \item Dominio $\R$, par, $0<h\le1$. Asíntota horizontal $y=0$. $h'(x)=-2xe^{-x^{2}}$: crece en $(-\infty,0]$, decrece en $[0,\infty)$; máximo absoluto $h(0)=1$. $h''(x)=(4x^{2}-2)e^{-x^{2}}$: convexa si $\abs x>\frac{1}{\sqrt2}$ y cóncava si $\abs x<\frac{1}{\sqrt2}$; inflexiones en $\bigl(\pm\frac{1}{\sqrt2},e^{-1/2}\bigr)\approx(\pm0{,}707;\,0{,}607)$.
\end{enumerate}
\resultado los datos anteriores determinan las tres gráficas: (a) dos ramas, con un máximo local en el origen y un mínimo local en $(2,4)$; (b) una «S» tumbada con extremos en $(\mp1,\pm2)$ y tangente vertical en el origen; (c) la campana de Gauss.
\comprobacion (b): $g(-1)=-1-3(-1)=2$ y $g(1)=1-3=-2$, coherente con la imparidad. (c): $h$ es, salvo constantes, la densidad normal del módulo~5.
\end{solucion}

\begin{solucion}{eje:5-optimizacion}
\datos cuatro problemas de extremos.
\planteamiento reducir a una variable, puntos críticos, y justificar que son extremos absolutos (signo de la derivada o \cref{prop:intervalo-cerrado}).
\pasos
\begin{enumerate}[label=(\alph*)]
  \item $x=a-2y$, $P(y)=(a-2y)y=ay-2y^{2}$; $P'(y)=a-4y=0$ en $y=\frac a4$, y $P''=-4<0$; como $P$ es una parábola cóncava, es el máximo absoluto. $x=\frac a2$, $P=\frac{a^{2}}{8}$.
  \item $f(x)=-x^{2}\ln x$; $f'(x)=-2x\ln x-x=-x(2\ln x+1)$, positiva si $x<e^{-1/2}$ y negativa si $x>e^{-1/2}$: máximo absoluto en $x=e^{-1/2}$, con $f=-e^{-1}\cdot(-\frac12)=\frac{1}{2e}\approx0{,}184$.
  \item $k(x)=(x+1)e^{-x}$; $k'(x)=e^{-x}-(x+1)e^{-x}=-xe^{-x}<0$ para $x>0$: $k$ es estrictamente decreciente en $[0,\infty)$. Máximo absoluto $k(0)=1$. El ínfimo es $\lim_{x\to\infty}k(x)=0$ (\cref{prop:crecimientos}), pero $k(x)>0$ siempre: \textbf{no hay mínimo}.
  \item Con el diámetro sobre el eje $x$ y centro en el origen, el rectángulo tiene base $2x$ y altura $\sqrt{a^{2}-x^{2}}$, $x\in[0,a]$. Se maximiza $A^{2}=4x^{2}(a^{2}-x^{2})$, función de $u=x^{2}\in[0,a^{2}]$: $4u(a^{2}-u)$, máxima en $u=\frac{a^{2}}{2}$. Luego $x=\frac{a}{\sqrt2}$, base $a\sqrt2$, altura $\frac{a}{\sqrt2}$, área $a^{2}$; en los extremos $x=0$, $x=a$ el área es nula.
\end{enumerate}
\resultado (a) $x=\frac a2$, $y=\frac a4$, máximo $\frac{a^{2}}{8}$; (b) $\frac{1}{2e}$ en $x=e^{-1/2}$; (c) máximo $1$ en $x=0$, sin mínimo; (d) base $a\sqrt2$, altura $a/\sqrt2$, área $a^{2}$.
\comprobacion (d) con $a=1$: $x=0{,}7$ da $2\cdot0{,}7\cdot\sqrt{0{,}51}\approx0{,}9998<1$; $x=0{,}72$ da $0{,}9993<1$.
\end{solucion}

\begin{solucion}{eje:5-cubica}
\datos $y=x^{3}+ax^{2}+bx+c$, $y'=3x^{2}+2ax+b$.
\planteamiento la cúbica tiene un máximo y un mínimo locales si y solo si $y'$ cambia de signo dos veces, es decir, si tiene dos raíces reales distintas.
\pasos
\begin{enumerate}[label=\arabic*.]
  \item Si el discriminante $4a^{2}-12b>0$, $y'$ tiene raíces $x_1<x_2$ y su signo es $+,-,+$: por el criterio de la primera derivada, máximo local en $x_1$ y mínimo local en $x_2$.
  \item Si $4a^{2}-12b\le0$, $y'\ge0$ en todo $\R$ (con a lo sumo un cero), $y$ es estrictamente creciente (\cref{cor:tvm}(iii), aplicado a cada lado del cero) y no tiene extremos.
  \item La condición es, pues, $a^{2}>3b$; $c$ no interviene (traslada la gráfica en vertical).
  \item Ejemplos: $x^{3}-3x+1$ ($a=0$, $b=-3$): $0>-9$, extremos en $x=\pm1$. $x^{3}+x$ ($a=0$, $b=1$): $0>3$ es falso; $y'=3x^{2}+1>0$, sin extremos.
\end{enumerate}
\resultado tiene máximo y mínimo locales si y solo si $a^{2}>3b$.
\comprobacion criterio equivalente: el punto de inflexión está en $x_0=-\frac a3$, y la condición $a^{2}>3b$ equivale a $y'(x_0)=b-\frac{a^{2}}{3}<0$: la pendiente en la inflexión es negativa.
\end{solucion}

\begin{solucion}{eje:5-una-raiz}
\datos $p(x)=x^{3}+3x^{2}+16$; $h(x)=e^{x}+x-2$.
\planteamiento existencia con Bolzano; unicidad con la monotonía (\cref{cor:tvm}(iii)).
\pasos
\begin{enumerate}[label=\arabic*.]
  \item $p'(x)=3x(x+2)$: $p$ crece en $(-\infty,-2]$, decrece en $[-2,0]$ y crece en $[0,\infty)$. Por tanto, en $[-2,\infty)$ se tiene $p(x)\ge p(0)=16>0$, y no hay raíces en ese intervalo.
  \item En $(-\infty,-2]$, $p$ es estrictamente creciente, $p(-2)=20>0$ y $p(x)\to-\infty$: hay exactamente una raíz. Como $p(-4)=-64+48+16=0$, es $x=-4$.
  \item $h'(x)=e^{x}+1>0$: $h$ es estrictamente creciente y tiene a lo sumo un cero. $h(0)=-1<0<e-1=h(1)$: por Bolzano, exactamente uno, en $(0,1)$. \hfill$\blacksquare$
\end{enumerate}
\resultado $x^{3}+3x^{2}+16=0$ tiene la única raíz real $-4$; $e^{x}=2-x$ tiene una única solución, en $(0,1)$.
\comprobacion $p(x)=(x+4)(x^{2}-x+4)$, y $x^{2}-x+4$ tiene discriminante $1-16<0$. La solución de $e^{x}=2-x$ es $x\approx0{,}4429$.
\end{solucion}

\begin{solucion}{eje:5-desigualdades}
\datos cuatro desigualdades.
\planteamiento \cref{teo:tvm} en $[0,x]$ (o entre $x$ e $y$) y acotación de $f'(c)$.
\pasos
\begin{enumerate}[label=\arabic*.]
  \item (a) Si $x>1$, $\sin x\le1<x$. Si $0<x\le1$, el TVM da $\sin x-\sin0=x\cos c$ con $c\in(0,x)\subseteq(0,1)$; ahí $\cos c<1$, luego $\sin x<x$.
  \item (b) $\sqrt{1+x}-1=\frac{x}{2\sqrt{1+c}}$ con $c\in(0,x)$; como $\sqrt{1+c}>1$, $\sqrt{1+x}-1<\frac x2$.
  \item (c) $\ln(1+x)-\ln1=\frac{x}{1+c}$ con $c\in(0,x)$, y $\frac{1}{1+x}<\frac{1}{1+c}<1$; multiplicando por $x>0$, $\frac{x}{1+x}<\ln(1+x)<x$.
  \item (d) $\abs{\arctan't}=\frac{1}{1+t^{2}}\le1$, y el \cref{cor:tvm}(iv) da $\abs{\arctan x-\arctan y}\le\abs{x-y}$. \hfill$\blacksquare$
\end{enumerate}
\resultado las cuatro desigualdades.
\comprobacion con $x=0{,}5$: $\sin0{,}5=0{,}479<0{,}5$; $\sqrt{1{,}5}=1{,}2247<1{,}25$; $0{,}3333<\ln1{,}5=0{,}4055<0{,}5$.
\end{solucion}

\begin{solucion}{eje:5-comparacion}
\datos $f(a)\ge g(a)$, $f'>g'$ en $(a,\infty)$.
\planteamiento aplicar el \cref{cor:tvm}(iii) a $h=f-g$; después, inducción.
\pasos
\begin{enumerate}[label=\arabic*.]
  \item (a) $h=f-g$ es continua en $[a,\infty)$ y $h'>0$ en $(a,\infty)$, luego es estrictamente creciente en $[a,\infty)$. Para $x>a$, $h(x)>h(a)\ge0$, es decir, $f(x)>g(x)$.
  \item (b) Primero, $e^{x}>1=e^{0}$ para $x>0$, porque $\exp$ es estrictamente creciente. Con $f=e^{x}$ y $g=1+x$: $f(0)=g(0)=1$ y $f'=e^{x}>1=g'$ para $x>0$; por (a), $e^{x}>1+x$.
  \item Con $f=e^{x}$ y $g=1+x+\frac{x^{2}}{2}$: $f(0)=g(0)$ y $f'=e^{x}>1+x=g'$ para $x>0$ por el paso anterior; por (a), $e^{x}>1+x+\frac{x^{2}}{2}$.
  \item (c) Sea $P_n(x)=\sum_{k=0}^{n}\frac{x^{k}}{k!}$; nótese que $P_{n+1}'=P_n$ y $P_{n+1}(0)=1$. Si $e^{x}>P_n(x)$ para $x>0$, (a) con $f=e^{x}$, $g=P_{n+1}$ da $e^{x}>P_{n+1}(x)$ para $x>0$. El caso $n=0$ es el primer paso de (b). \hfill$\blacksquare$
\end{enumerate}
\resultado $e^{x}>\sum_{k=0}^{n}x^{k}/k!$ para todo $x>0$ y todo $n$.
\comprobacion con $x=1$, $n=3$: $e=2{,}718>1+1+0{,}5+0{,}1667=2{,}667$. El \cref{teo:taylor} da la misma conclusión: el resto es $\frac{e^{\xi}x^{n+1}}{(n+1)!}>0$.
\end{solucion}

\begin{solucion}{eje:5-raices-polinomio}
\datos $p$ de grado $n$ con raíces $x_1<\dots<x_n$; $f$ dos veces derivable con ceros $a<b<c$.
\planteamiento Rolle entre ceros consecutivos.
\pasos
\begin{enumerate}[label=\arabic*.]
  \item (a) Para $i=1,\dots,n-1$, $p(x_i)=p(x_{i+1})=0$ y Rolle da $c_i\in(x_i,x_{i+1})$ con $p'(c_i)=0$. Los intervalos son disjuntos, luego los $c_i$ son distintos: $p'$ tiene al menos $n-1$ raíces reales distintas, y no más, por ser de grado $n-1$.
  \item (b) Rolle en $[a,b]$ y en $[b,c]$ da $p_1\in(a,b)$ y $p_2\in(b,c)$ con $f'(p_1)=f'(p_2)=0$, y $p_1<p_2$. Rolle aplicado a $f'$ (derivable) en $[p_1,p_2]$ da $p\in(p_1,p_2)\subseteq(a,c)$ con $f''(p)=0$.
  \item (c) $q(x)=x^{5}+x-1$ tiene $q'(x)=5x^{4}+1>0$. Si tuviera dos raíces, Rolle daría un cero de $q'$; luego tiene a lo sumo una. Como $q(0)=-1<0<1=q(1)$, Bolzano da exactamente una, en $(0,1)$. \hfill$\blacksquare$
\end{enumerate}
\resultado (a) $p'$ tiene $n-1$ raíces reales distintas; (b) existe $p$; (c) una única raíz real.
\comprobacion (b) con $f=(x-a)(x-b)(x-c)$: $f''(x)=6x-2(a+b+c)$ se anula en la media $\frac{a+b+c}{3}\in(a,c)$. (c): la raíz es $x\approx0{,}7549$.
\end{solucion}

\begin{solucion}{eje:5-tvm-demo}
\datos $f$ dos veces derivable en $\R$.
\planteamiento (a) como en el \cref{teo:tvm}; (b) regla de L'Hôpital en la variable $h$, una sola vez, y la definición de $f''(a)$.
\pasos
\begin{enumerate}[label=\arabic*.]
  \item (a) \emph{Rolle:} si $f$ es continua en $[a,b]$, derivable en $(a,b)$ y $f(a)=f(b)$, existe $c\in(a,b)$ con $f'(c)=0$. \emph{TVM:} con $m=\frac{f(b)-f(a)}{b-a}$, la función $g(x)=f(x)-m(x-a)$ cumple las hipótesis de Rolle, pues $g(a)=g(b)=f(a)$; luego existe $c$ con $g'(c)=f'(c)-m=0$.
  \item (b) Sean $N(h)=f(a+h)+f(a-h)-2f(a)$ y $D(h)=h^{2}$. Ambas son derivables, $N(h)\to0$ y $D(h)\to0$ cuando $h\to0$, y $D'(h)=2h\ne0$ para $h\ne0$.
  \item $\frac{N'(h)}{D'(h)}=\frac{f'(a+h)-f'(a-h)}{2h}=\frac12\Bigl[\frac{f'(a+h)-f'(a)}{h}+\frac{f'(a-h)-f'(a)}{-h}\Bigr]\to\frac12\bigl[f''(a)+f''(a)\bigr]=f''(a)$, por la definición de $f''(a)$ como derivada de $f'$ (en el segundo cociente, con incremento $-h\to0$).
  \item Por el \cref{teo:lhopital}, $\frac{N(h)}{D(h)}\to f''(a)$. \hfill$\blacksquare$
\end{enumerate}
\resultado $f''(a)=\lim_{h\to0}\frac{f(a+h)+f(a-h)-2f(a)}{h^{2}}$.
\comprobacion con $f=\sin$, $a=1$, $h=0{,}01$: el cociente vale $-0{,}841464$, frente a $f''(1)=-\sin1=-0{,}841471$. Aplicar L'Hôpital dos veces exigiría que $f''$ fuera continua en $a$, hipótesis que no se tiene.
\end{solucion}

\begin{solucion}{eje:5-cauchy}
\datos $h(x)=f(x)-kg(x)$ con $k=\frac{f(b)-f(a)}{g(b)-g(a)}$.
\planteamiento (a) cálculo directo; (b) pendiente de la tangente a una curva paramétrica; (c) como en la demostración del \cref{teo:lhopital}.
\pasos
\begin{enumerate}[label=\arabic*.]
  \item (a) $h(b)-h(a)=\bigl(f(b)-f(a)\bigr)-k\bigl(g(b)-g(a)\bigr)=0$, por la definición de $k$.
  \item (b) La curva $t\mapsto(g(t),f(t))$, $t\in[a,b]$, une $(g(a),f(a))$ con $(g(b),f(b))$; la cuerda tiene pendiente $k$, y la tangente en $t=c$ tiene vector director $(g'(c),f'(c))$ y pendiente $f'(c)/g'(c)$. El teorema afirma que hay un punto de la curva con tangente paralela a la cuerda; con $g(t)=t$ es el TVM.
  \item (c) Sean $f(x),g(x)\to0$ cuando $x\to a^{+}$ y $\frac{f'}{g'}\to\ell$. Se define $f(a)=g(a)=0$. Para $x>a$ próximo, $g(x)\ne0$: si fuera $g(x)=0=g(a)$, Rolle daría un cero de $g'$ en $(a,x)$. Por el \cref{teo:cauchy} en $[a,x]$, $\frac{f(x)}{g(x)}=\frac{f'(c_x)}{g'(c_x)}$ con $a<c_x<x$, y cuando $x\to a^{+}$, $c_x\to a^{+}$ y el cociente tiende a $\ell$. \hfill$\blacksquare$
\end{enumerate}
\resultado $h(a)=h(b)$; interpretación geométrica; regla de L'Hôpital.
\comprobacion la hipótesis $g'\ne0$ se usa dos veces: para que $g(x)\ne0$ (y el cociente $f/g$ tenga sentido) y para que $f'(c_x)/g'(c_x)$ esté definido.
\end{solucion}

\begin{solucion}{eje:5-darboux}
\datos $g$ derivable cerca de $[a,b]$, $g'(a)<k<g'(b)$.
\planteamiento Weierstrass para $\varphi(x)=g(x)-kx$ y Fermat en el mínimo.
\pasos
\begin{enumerate}[label=\arabic*.]
  \item $\varphi$ es continua en $[a,b]$; por Weierstrass alcanza su mínimo en algún $c\in[a,b]$.
  \item $\varphi'(a)=g'(a)-k<0$. Como $\frac{\varphi(a+h)-\varphi(a)}{h}\to\varphi'(a)<0$, para $h>0$ pequeño $\varphi(a+h)<\varphi(a)$: el mínimo no está en $a$.
  \item $\varphi'(b)=g'(b)-k>0$. Como $\frac{\varphi(b-h)-\varphi(b)}{-h}\to\varphi'(b)>0$, para $h>0$ pequeño $\varphi(b-h)<\varphi(b)$: el mínimo no está en $b$.
  \item Luego $c\in(a,b)$ y, por el \cref{teo:fermat}, $\varphi'(c)=0$, es decir, $g'(c)=k$.
  \item \emph{Saltos.} Supóngase que $g'$ tiene en $x_0$ límites laterales $L_-<L_+$ (el caso $L_->L_+$ es análogo). Se eligen $k_1<k_2$ en $(L_-,L_+)$ tales que $g'(x_0)\notin[k_1,k_2]$, y $k\in(k_1,k_2)$. Existe $\delta>0$ con $g'<k_1$ en $(x_0-\delta,x_0)$ y $g'>k_2$ en $(x_0,x_0+\delta)$. Tomando $a$ y $b$ en esos intervalos, Darboux da $c\in(a,b)$ con $g'(c)=k$; pero $c=x_0$ es imposible ($g'(x_0)\ne k$), $c<x_0$ da $g'(c)<k$ y $c>x_0$ da $g'(c)>k$. Contradicción.
  \item La función escalón solo toma los valores $0$ y~$1$; si fuera $g'$, Darboux en $[-1,1]$ con $k=\frac12$ daría $g'(c)=\frac12$. \hfill$\blacksquare$
\end{enumerate}
\resultado las derivadas cumplen la propiedad de los valores intermedios, aunque no sean continuas, y no tienen discontinuidades de salto.
\comprobacion la derivada del \cref{eje:4-x2sen} es discontinua en~0, pero su discontinuidad es esencial (oscila entre $-1$ y $1$), no de salto, en coherencia con este resultado.
\end{solucion}

\begin{solucion}{eje:5-creciente-derivada}
\datos $f$ creciente y derivable en $a$.
\planteamiento signo de los cocientes incrementales y paso al límite.
\pasos
\begin{enumerate}[label=\arabic*.]
  \item (a) Si $h>0$, $f(a+h)\ge f(a)$ y el cociente $\frac{f(a+h)-f(a)}{h}$ es $\ge0$; si $h<0$, numerador $\le0$ y denominador $<0$, y el cociente también es $\ge0$. El límite de una función $\ge0$ es $\ge0$: $f'(a)\ge0$.
  \item (b) Si $f$ es estrictamente creciente, los cocientes son $>0$, pero el paso al límite solo conserva desigualdades no estrictas. Contraejemplo: $f(x)=x^{3}$ es estrictamente creciente y $f'(0)=0$.
  \item (c) $\sigma'=\sigma(1-\sigma)>0$ y $\tanh'=1-\tanh^{2}>0$ en todo $\R$, luego son estrictamente crecientes (\cref{cor:tvm}(iii)). Si $x<y$, $\ReLU(x)=\max(0,x)\le\max(0,y)=\ReLU(y)$, luego $\ReLU$ es creciente; no estrictamente, porque $\ReLU(-2)=\ReLU(-1)=0$. \hfill$\blacksquare$
\end{enumerate}
\resultado $f'(a)\ge0$; la desigualdad puede ser una igualdad aunque $f$ sea estrictamente creciente.
\comprobacion el recíproco parcial es el \cref{cor:tvm}(iii): $f'>0$ en un intervalo implica crecimiento estricto; $x^{3}$ muestra que la condición es suficiente pero no necesaria.
\end{solucion}

\begin{solucion}{eje:5-max-convexas}
\datos la \cref{def:convexa}.
\planteamiento encadenar desigualdades con $z=tx+(1-t)y$.
\pasos
\begin{enumerate}[label=\arabic*.]
  \item (a) Sea $h=\max(f,g)$ con $f,g$ convexas, $x,y\in I$, $t\in[0,1]$, $z=tx+(1-t)y$. Entonces $f(z)\le tf(x)+(1-t)f(y)\le th(x)+(1-t)h(y)$, y lo mismo para $g(z)$; luego $h(z)\le th(x)+(1-t)h(y)$.
  \item Las funciones afines $\ell(x)=mx+n$ son convexas (la desigualdad es una igualdad). Como $\ReLU(x)=\max(0,x)$ y $\abs x=\max(x,-x)$, ambas son convexas.
  \item (b) Sumando las desigualdades de $f$ y $g$ en $z$ se obtiene la de $f+g$.
  \item Si $f$ es convexa y $g$ convexa y creciente: $g(f(z))\le g\bigl(tf(x)+(1-t)f(y)\bigr)\le tg(f(x))+(1-t)g(f(y))$; la primera desigualdad usa que $g$ es creciente y $f$ convexa, y la segunda, que $g$ es convexa.
  \item (c) $x^{2}-1$ es convexa (derivada segunda $2>0$) y $\softplus$ es convexa ($\softplus''=\sigma'>0$) y creciente ($\softplus'=\sigma>0$). Por (b), $\softplus(x^{2}-1)$ es convexa. \hfill$\blacksquare$
\end{enumerate}
\resultado el máximo y la suma de convexas son convexas, y la composición con una convexa creciente también; $\ReLU$, $\abs x$ y $\softplus(x^{2}-1)$ son convexas.
\comprobacion la hipótesis de crecimiento es necesaria: $f(x)=x^{2}$ y $g(u)=-u$ (convexa, decreciente) dan $-x^{2}$, que no es convexa.
\end{solucion}

\begin{solucion}{eje:5-jensen}
\datos $f$ convexa en un intervalo $I$, $x_i\in I$, $t_i\ge0$, $\sum t_i=1$.
\planteamiento inducción sobre $n$, agrupando los $n$ primeros puntos.
\pasos
\begin{enumerate}[label=\arabic*.]
  \item Para $n=1$ es una igualdad, y para $n=2$ es la \cref{def:convexa}. Supóngase cierta para $n$ y sean $t_1,\dots,t_{n+1}$.
  \item Si $t_{n+1}=1$, todos los demás son nulos y la desigualdad es una igualdad. Si no, sea $s=1-t_{n+1}=\sum_{i\le n}t_i>0$ e $y=\sum_{i\le n}\frac{t_i}{s}x_i$, que es una combinación convexa de puntos de $I$ y, por tanto, está en $I$ (está entre el mínimo y el máximo de los $x_i$).
  \item Por la definición con dos puntos, $f\bigl(\sum_{i\le n+1}t_ix_i\bigr)=f(sy+t_{n+1}x_{n+1})\le sf(y)+t_{n+1}f(x_{n+1})$.
  \item Por la hipótesis de inducción con los pesos $t_i/s$, que suman~1, $f(y)\le\sum_{i\le n}\frac{t_i}{s}f(x_i)$. Sustituyendo, $f\bigl(\sum t_ix_i\bigr)\le\sum_{i\le n+1}t_if(x_i)$.
  \item (b) $-\ln$ es convexa en $(0,\infty)$, pues $(-\ln)''=1/x^{2}>0$. Con $t_i=\frac1n$: $-\ln\frac{\sum x_i}{n}\le-\frac1n\sum\ln x_i=-\ln\sqrt[n]{x_1\cdots x_n}$. Cambiando de signo y aplicando la exponencial, que es creciente, $\sqrt[n]{x_1\cdots x_n}\le\frac{x_1+\dots+x_n}{n}$. \hfill$\blacksquare$
\end{enumerate}
\resultado la desigualdad de Jensen y la desigualdad entre las medias aritmética y geométrica.
\comprobacion $x=(1,4,16)$: media geométrica $\sqrt[3]{64}=4$, aritmética $7$. Esta desigualdad, con $-\ln$, es la que en el módulo~1 relaciona la perplejidad con una media geométrica.
\end{solucion}
% ---------------------------------------------------------------------
\section{Tema 6}

\begin{solucion}{eje:6-coeficientes}
\datos $P_{2,0}f(x)=3-x+5x^{2}$; $p(x)=x^{3}$.
\planteamiento los coeficientes del polinomio de Taylor son $f^{(k)}(a)/k!$ (\cref{def:taylor}).
\pasos
\begin{enumerate}[label=\arabic*.]
  \item (a) $f(0)=3$; $f'(0)=-1$; $\frac{f''(0)}{2}=5$, luego $f''(0)=10$.
  \item Como $f'(0)=-1\ne0$, por el \cref{teo:fermat} $f$ no tiene extremo local en~0.
  \item (b) $p(1)=1$, $p'(1)=3$, $p''(1)=6$, $p'''(1)=6$: $P_{3,1}p(x)=1+3(x-1)+3(x-1)^{2}+(x-1)^{3}$.
  \item Por el binomio de Newton, $(1+(x-1))^{3}=1+3(x-1)+3(x-1)^{2}+(x-1)^{3}$, que es $x^{3}$.
\end{enumerate}
\resultado (a) $f(0)=3$, $f'(0)=-1$, $f''(0)=10$; no hay extremo en~0. (b) $1+3(x-1)+3(x-1)^{2}+(x-1)^{3}=x^{3}$.
\comprobacion en (b), con $x=2$: $1+3+3+1=8=2^{3}$.
\end{solucion}

\begin{solucion}{eje:6-por-que}
\datos la \cref{prop:taylor-unico} y el \cref{eje:4-par-impar}.
\planteamiento unicidad del polinomio de Taylor y paridad de las derivadas sucesivas.
\pasos
\begin{enumerate}[label=\arabic*.]
  \item Si $p$ es un polinomio de grado $\le n$, el propio $p$ es un polinomio de grado $\le n$ con las mismas derivadas que $p$ en $a$ hasta el orden $n$; por la unicidad, $P_{n,a}p=p$.
  \item Si $f$ es par, $f'$ es impar, $f''$ par, y en general $f^{(k)}$ es impar para $k$ impar. Una función impar definida en~0 se anula en~0, luego $f^{(k)}(0)=0$ para $k$ impar y solo quedan potencias pares.
  \item $\sin$ es impar, luego su polinomio de Maclaurin solo tiene potencias impares: el coeficiente de $x^{4}$ es $\sin^{(4)}(0)/4!=\sin0/24=0$, y $P_{4,0}\sin=P_{3,0}\sin$.
\end{enumerate}
\resultado las tres explicaciones anteriores.
\comprobacion $P_{4,0}\cos=1-\frac{x^{2}}{2}+\frac{x^{4}}{24}$ (solo potencias pares) y $P_{5,0}\tanh=x-\frac{x^{3}}{3}+\frac{2x^{5}}{15}$ (solo impares).
\end{solucion}

\begin{solucion}{eje:6-maclaurin}
\datos cinco funciones, $a=0$.
\planteamiento calcular $f^{(k)}(0)$ para $k\le n$ y dividir por $k!$.
\pasos
\begin{enumerate}[label=(\alph*)]
  \item $\cos$: derivadas en~0: $1,0,-1,0,1$. $P_4=1-\frac{x^{2}}{2}+\frac{x^{4}}{24}$.
  \item $e^{-x}$: $f^{(k)}(0)=(-1)^{k}$. $P_3=1-x+\frac{x^{2}}{2}-\frac{x^{3}}{6}$.
  \item $\frac{1}{1-x}$: $f^{(k)}(x)=\frac{k!}{(1-x)^{k+1}}$, $f^{(k)}(0)/k!=1$. $P_3=1+x+x^{2}+x^{3}$.
  \item $\sqrt{1+x}$: $f(0)=1$, $f'(0)=\frac12$, $f''(0)=-\frac14$. $P_2=1+\frac x2-\frac{x^{2}}{8}$.
  \item $\arctan$: $f'=\frac{1}{1+x^{2}}$, $f''=\frac{-2x}{(1+x^{2})^{2}}$, $f'''=\frac{6x^{2}-2}{(1+x^{2})^{3}}$; en~0: $0,1,0,-2$. $P_3=x-\frac{x^{3}}{3}$.
\end{enumerate}
\resultado los polinomios anteriores.
\comprobacion (b) es (a) del \cref{ej:maclaurin} con $x$ cambiado por $-x$; (d) es el binomio generalizado con $\alpha=\frac12$: $\binom{1/2}{2}=\frac{\frac12\cdot(-\frac12)}{2}=-\frac18$.
\end{solucion}

\begin{solucion}{eje:6-otros-puntos}
\datos $\ln x$ y $1/x$ en $a=1$; $\sqrt x$ en $a=4$.
\planteamiento derivadas en el punto y \cref{def:taylor}, en potencias de $(x-a)$.
\pasos
\begin{enumerate}[label=(\alph*)]
  \item $\ln$: en $x=1$, $f=0$, $f'=1$, $f''=-1$, $f'''=2$. $P_{3,1}=(x-1)-\frac{(x-1)^{2}}{2}+\frac{(x-1)^{3}}{3}$.
  \item $\sqrt x$: en $x=4$, $f=2$, $f'=\frac{1}{2\sqrt4}=\frac14$, $f''=-\frac{1}{4\cdot4^{3/2}}=-\frac{1}{32}$. $P_{2,4}=2+\frac{x-4}{4}-\frac{(x-4)^{2}}{64}$.
  \item $1/x$: $f^{(k)}(1)=(-1)^{k}k!$. $P_{3,1}=1-(x-1)+(x-1)^{2}-(x-1)^{3}$.
  \item $\ln1{,}1\approx0{,}1-0{,}005+0{,}000333=0{,}095333$, frente a $\ln1{,}1=0{,}095310$: error $2{,}3\cdot10^{-5}$.
\end{enumerate}
\resultado (a) $(x-1)-\frac{(x-1)^{2}}{2}+\frac{(x-1)^{3}}{3}$; (b) $2+\frac{x-4}{4}-\frac{(x-4)^{2}}{64}$; (c) $1-(x-1)+(x-1)^{2}-(x-1)^{3}$; $\ln1{,}1\approx0{,}09533$.
\comprobacion el resto de Lagrange de (a) es $\frac{f^{(4)}(\xi)}{24}(0{,}1)^{4}$ con $f^{(4)}(\xi)=-6/\xi^{4}$ y $\xi\in(1;1{,}1)$; su valor absoluto es menor que $\frac{6}{24}\cdot10^{-4}=2{,}5\cdot10^{-5}$, y el error real, $2{,}3\cdot10^{-5}$, lo respeta.
\end{solucion}

\begin{solucion}{eje:6-ln12}
\datos $f$ dos veces derivable; $f=\ln$, $a=1$, $x=1{,}2$.
\planteamiento $P_{2,a}$ y resto de orden~3 (\cref{teo:taylor} con $n=2$).
\pasos
\begin{enumerate}[label=\arabic*.]
  \item (a) $f(x)\approx f(a)+f'(a)(x-a)+\frac{f''(a)}{2}(x-a)^{2}$.
  \item (b) Con $\ln$ en $a=1$: $\ln(1+h)\approx h-\frac{h^{2}}{2}$; para $h=0{,}2$, $\ln1{,}2\approx0{,}2-0{,}02=0{,}18$.
  \item (c) $R_2=\frac{f'''(\xi)}{6}h^{3}$ con $f'''(\xi)=\frac{2}{\xi^{3}}$, $\xi\in(1;1{,}2)$. Como $\xi>1$, $0<R_2<\frac{2}{6}(0{,}2)^{3}=0{,}00267$.
  \item El resto es positivo: $0{,}18$ aproxima por defecto, y $0{,}18<\ln1{,}2<0{,}18267$.
\end{enumerate}
\resultado $\ln1{,}2\approx0{,}18$, con error menor que $0{,}0027$ y por defecto.
\comprobacion $\ln1{,}2=0{,}18232$; el error real es $0{,}00232$, dentro de la cota y del signo previstos.
\end{solucion}

\begin{solucion}{eje:6-cuadraticas}
\datos seis funciones cerca de~0.
\planteamiento $P_{2,0}$, o composición de desarrollos conocidos (\cref{tab:maclaurin}) truncando en grado~2.
\pasos
\begin{enumerate}[label=(\alph*)]
  \item $\tan0=0$, $\tan'0=1$, $\tan''x=2\tan x(1+\tan^{2}x)$ vale~0 en~0: $\tan x\approx x$ (el término cuadrático es nulo, por ser $\tan$ impar).
  \item $\ln\cos x=\ln\bigl(1-\frac{x^{2}}{2}+o(x^{2})\bigr)=-\frac{x^{2}}{2}+o(x^{2})$, usando $\ln(1+u)=u+o(u)$.
  \item $\sin2x\approx2x$; $\cos2x\approx1-\frac{(2x)^{2}}{2}=1-2x^{2}$.
  \item $\frac{1}{\cos x}=\frac{1}{1-x^{2}/2+o(x^{2})}=1+\frac{x^{2}}{2}+o(x^{2})$, usando $\frac{1}{1-u}=1+u+o(u)$.
  \item $e^{x^{2}}=1+x^{2}+o(x^{2})$.
  \item $\frac{1}{a+bx}=\frac1a\cdot\frac{1}{1+\frac bax}=\frac1a\Bigl(1-\frac bax+\frac{b^{2}}{a^{2}}x^{2}\Bigr)+o(x^{2})=\frac1a-\frac{b}{a^{2}}x+\frac{b^{2}}{a^{3}}x^{2}+o(x^{2})$.
\end{enumerate}
\resultado (a) $x$; (b) $-\frac{x^{2}}{2}$; (c) $2x$ y $1-2x^{2}$; (d) $1+\frac{x^{2}}{2}$; (e) $1+x^{2}$; (f) $\frac1a-\frac{b}{a^{2}}x+\frac{b^{2}}{a^{3}}x^{2}$.
\comprobacion (b) derivando: $(\ln\cos x)'=-\tan x$, $(\ln\cos x)''=-(1+\tan^{2}x)$, que vale $-1$ en~0: coeficiente $-\frac12$.
\end{solucion}

\begin{solucion}{eje:6-error}
\datos (a) $x=0{,}1$; (b) $x=0{,}5$, error $<10^{-5}$; (c) $\sqrt[3]{8{,}1}$, $a=8$.
\planteamiento \cref{cor:cota} con una cota de la derivada siguiente en el intervalo.
\pasos
\begin{enumerate}[label=(\alph*)]
  \item $\sin0{,}1\approx0{,}1-\frac{0{,}001}{6}=0{,}0998333$. Como $P_{3,0}\sin=P_{4,0}\sin$, se puede usar el resto de orden~5: $\abs R\le\frac{(0{,}1)^{5}}{120}\approx8{,}3\cdot10^{-8}$.
  \item $\abs{R_n}=\frac{e^{\xi}(0{,}5)^{n+1}}{(n+1)!}<\frac{2\cdot(0{,}5)^{n+1}}{(n+1)!}$, pues $e^{\xi}<e^{0{,}5}<2$. Para $n=5$ la cota es $4{,}3\cdot10^{-5}$; para $n=6$, $3{,}1\cdot10^{-6}<10^{-5}$. Basta $n=6$.
  \item $f(x)=x^{1/3}$, $f(8)=2$, $f'(8)=\frac13\cdot8^{-2/3}=\frac{1}{12}$: $\sqrt[3]{8{,}1}\approx2+\frac{0{,}1}{12}=2{,}008333$. $f''(t)=-\frac29t^{-5/3}$, y para $t\ge8$, $\abs{f''(t)}\le\frac29\cdot\frac{1}{32}=\frac{1}{144}$. Luego $\abs R\le\frac{1}{2}\cdot\frac{1}{144}\cdot(0{,}1)^{2}\approx3{,}47\cdot10^{-5}$.
\end{enumerate}
\resultado (a) $0{,}0998333$ con error $<8{,}4\cdot10^{-8}$; (b) $n=6$; (c) $2{,}008333$ con error $<3{,}5\cdot10^{-5}$ (por exceso, porque $f''<0$).
\comprobacion valores reales: $\sin0{,}1=0{,}09983342$ (error $8{,}3\cdot10^{-8}$); $P_{6,0}(0{,}5)=1{,}6487196$ frente a $e^{0{,}5}=1{,}6487213$ (error $1{,}7\cdot10^{-6}$); $\sqrt[3]{8{,}1}=2{,}0082989$ (error $3{,}45\cdot10^{-5}$).
\end{solucion}

\begin{solucion}{eje:6-limites}
\datos cinco límites en~0.
\planteamiento sustituir desarrollos hasta el primer término no nulo del numerador (\cref{tab:maclaurin}, con un término más cuando haga falta).
\pasos
\begin{enumerate}[label=(\alph*)]
  \item $e^{x}-1-x-\frac{x^{2}}{2}=\frac{x^{3}}{6}+o(x^{3})$: límite $\frac16$.
  \item $\cos x-1+\frac{x^{2}}{2}=\frac{x^{4}}{24}+o(x^{5})$: límite $\frac{1}{24}$.
  \item $\sin x-x\cos x=\bigl(x-\frac{x^{3}}{6}\bigr)-x\bigl(1-\frac{x^{2}}{2}\bigr)+o(x^{4})=\frac{x^{3}}{3}+o(x^{4})$: límite $\frac13$.
  \item $\tanh x-x=-\frac{x^{3}}{3}+o(x^{4})$: límite $-\frac13$.
  \item $\sigma(x)-\frac12-\frac x4=-\frac{x^{3}}{48}+o(x^{4})$ (\cref{ej:taylor-activaciones}): límite $-\frac{1}{48}$.
\end{enumerate}
\resultado (a) $\frac16$; (b) $\frac{1}{24}$; (c) $\frac13$; (d) $-\frac13$; (e) $-\frac{1}{48}$.
\comprobacion (c) con L'Hôpital: $\frac{\cos x-\cos x+x\sin x}{3x^{2}}=\frac{\sin x}{3x}\to\frac13$. (d) y (e) son coherentes entre sí: $\sigma(x)-\frac12=\frac12\tanh\frac x2$, y $\frac12\cdot\bigl(-\frac13\bigr)\cdot\frac18=-\frac{1}{48}$.
\end{solucion}

\begin{solucion}{eje:6-poligono}
\datos polígono regular de $n$ lados inscrito en la circunferencia unidad.
\planteamiento teorema del coseno en el triángulo formado por el centro y dos vértices consecutivos; aproximación $\cos\theta=1-\frac{\theta^{2}}{2}+o(\theta^{3})$.
\pasos
\begin{enumerate}[label=\arabic*.]
  \item (a) Dos radios de longitud~1 forman un ángulo central $\theta=\frac{2\pi}{n}$; el lado opuesto mide $\ell$ con $\ell^{2}=1+1-2\cos\theta$, es decir, $\ell=\sqrt{2-2\cos\frac{2\pi}{n}}$.
  \item (b) $2-2\cos\theta=\theta^{2}+o(\theta^{3})$, luego $\ell=\theta\sqrt{1+o(\theta)}$ y el perímetro es $n\ell=n\theta\sqrt{1+o(\theta)}=2\pi\sqrt{1+o(1/n)}\to2\pi$.
  \item (c) $n=6$: $\ell=\sqrt{2-2\cdot\frac12}=1$, perímetro $6$. $n=100$: perímetro $\approx6{,}28215$.
\end{enumerate}
\resultado el perímetro tiende a $2\pi$, la longitud de la circunferencia.
\comprobacion $2\pi\approx6{,}28319$; para $n=100$ la diferencia es $0{,}00103$, del orden de $\frac{2\pi\cdot\theta^{2}}{24}=\frac{(2\pi)^{3}}{24n^{2}}\approx0{,}00103$, que es lo que predice el término siguiente del desarrollo.
\end{solucion}

\begin{solucion}{eje:6-taylor-orden1}
\datos $F(t)=f(t)-f(a)-f'(a)(t-a)$, $G(t)=(t-a)^{2}$.
\planteamiento dos aplicaciones del \cref{teo:cauchy}.
\pasos
\begin{enumerate}[label=\arabic*.]
  \item Sea $x\ne a$ (si $x=a$ es trivial). $F(a)=G(a)=0$, y $G'(t)=2(t-a)\ne0$ para $t$ estrictamente entre $a$ y $x$.
  \item Por Cauchy entre $a$ y $x$: $\frac{F(x)}{G(x)}=\frac{F(x)-F(a)}{G(x)-G(a)}=\frac{F'(c_1)}{G'(c_1)}=\frac{f'(c_1)-f'(a)}{2(c_1-a)}$, con $c_1$ entre $a$ y $x$.
  \item Sean $F_1(t)=f'(t)-f'(a)$ y $G_1(t)=2(t-a)$, derivables con $G_1'=2\ne0$ y $F_1(a)=G_1(a)=0$. Por Cauchy entre $a$ y $c_1$: $\frac{F_1(c_1)}{G_1(c_1)}=\frac{f''(\xi)}{2}$, con $\xi$ entre $a$ y $c_1$, y por tanto entre $a$ y $x$.
  \item Luego $F(x)=\frac{f''(\xi)}{2}G(x)$, es decir, $f(x)=f(a)+f'(a)(x-a)+\frac12f''(\xi)(x-a)^{2}$.
  \item Si $f''\ge0$, el último término es $\ge0$ y $f(x)\ge f(a)+f'(a)(x-a)$: la gráfica está por encima de la tangente en $a$. \hfill$\blacksquare$
\end{enumerate}
\resultado el teorema de Taylor de orden~1 con resto de Lagrange, y la desigualdad de la tangente para $f''\ge0$.
\comprobacion con $f(x)=x^{3}$, $a=0$: $x^{3}=0+0+\frac12\cdot6\xi\cdot x^{2}$ da $\xi=\frac x3$, que está entre $0$ y $x$. La conclusión final reproduce el \cref{teo:convexidad}(i) por otra vía.
\end{solucion}

\begin{solucion}{eje:6-binomica}
\datos $f(x)=(1+x)^{1/2}$, $c_n=\binom{1/2}{n}$.
\planteamiento (b) cociente de coeficientes consecutivos e inducción; (c) acotar el resto de Lagrange.
\pasos
\begin{enumerate}[label=\arabic*.]
  \item (a) Si $f$ es $n+1$ veces derivable en un intervalo abierto que contiene a $a$ y $x$, existe $\xi$ entre $a$ y $x$ con $f(x)=\sum_{k=0}^{n}\frac{f^{(k)}(a)}{k!}(x-a)^{k}+\frac{f^{(n+1)}(\xi)}{(n+1)!}(x-a)^{n+1}$ (\cref{teo:taylor}).
  \item (b) $c_0=1$ y $c_{n+1}=c_n\cdot\frac{\frac12-n}{n+1}$. Como $\abs{\frac12-n}\le n+\frac12<n+1$, $\abs{c_{n+1}}\le\abs{c_n}$; por inducción, $\abs{c_n}\le1$.
  \item (c) $f^{(k)}(t)=\frac12\bigl(\frac12-1\bigr)\cdots\bigl(\frac12-k+1\bigr)(1+t)^{\frac12-k}=k!\,c_k(1+t)^{\frac12-k}$. Luego $f^{(k)}(0)/k!=c_k$ y, para $x\in(0,1)$, el resto es $R_n(x)=c_{n+1}(1+\xi)^{\frac12-n-1}x^{n+1}$ con $0<\xi<x$.
  \item Como $\frac12-n-1<0$ y $1+\xi>1$, $(1+\xi)^{\frac12-n-1}<1$; por (b), $\abs{R_n(x)}\le x^{n+1}$, que tiende a~0 porque $0<x<1$.
  \item Por tanto, $\sum_{k=0}^{n}c_kx^{k}=f(x)-R_n(x)\to\sqrt{1+x}$. \hfill$\blacksquare$
\end{enumerate}
\resultado $\sqrt{1+x}=\sum_{n\ge0}\binom{1/2}{n}x^{n}$ para $x\in(0,1)$.
\comprobacion los primeros coeficientes son $1,\frac12,-\frac18,\frac1{16}$; con $x=0{,}21$, $1+0{,}105-0{,}0055125+0{,}0005788=1{,}1000663$, y $\sqrt{1{,}21}=1{,}1$, con error del orden de $\abs{c_4}x^{4}\approx7{,}6\cdot10^{-5}$.
\end{solucion}

\begin{solucion}{eje:6-ln2}
\datos $f(x)=\ln(1+x)$, $x\in[0,1]$.
\planteamiento resto de Lagrange con $f^{(n+1)}(t)=\frac{(-1)^{n}n!}{(1+t)^{n+1}}$ (\cref{ej:maclaurin}(c)).
\pasos
\begin{enumerate}[label=\arabic*.]
  \item $R_n(x)=\frac{f^{(n+1)}(\xi)}{(n+1)!}x^{n+1}=\frac{(-1)^{n}}{(n+1)(1+\xi)^{n+1}}x^{n+1}$ con $0\le\xi\le x$.
  \item Como $(1+\xi)^{n+1}\ge1$, $\abs{R_n(x)}\le\frac{x^{n+1}}{n+1}$.
  \item Con $x=1$: $\bigl|\ln2-\sum_{k=1}^{n}\frac{(-1)^{k-1}}{k}\bigr|\le\frac{1}{n+1}\to0$, luego la suma parcial tiende a $\ln2$.
  \item Error menor que $10^{-3}$: basta $\frac{1}{n+1}<10^{-3}$, es decir, $n\ge1000$. \hfill$\blacksquare$
\end{enumerate}
\resultado $\ln2=1-\frac12+\frac13-\cdots$; la cota garantiza error $<10^{-3}$ con $1000$ términos.
\comprobacion la convergencia es muy lenta: con $n=10$, la suma es $0{,}6456$ y $\ln2=0{,}6931$, error $0{,}048<\frac1{11}$. Para calcular $\ln2$ en la práctica se usan otras series, como la de $\artanh$: $\ln2=2\artanh\frac13$.
\end{solucion}

\begin{solucion}{eje:6-e-irracional}
\datos $e=\sum_{k=0}^{n}\frac1{k!}+\frac{e^{\xi}}{(n+1)!}$ con $0<\xi<1$ (\cref{ej:calculo-e}); $2<e<3$.
\planteamiento reducción al absurdo: multiplicar por $q!$ y obtener un entero entre $0$ y $1$.
\pasos
\begin{enumerate}[label=\arabic*.]
  \item Supóngase $e=p/q$ con $p,q\in\N$; multiplicando numerador y denominador por~2 si hace falta, puede suponerse $q\ge2$.
  \item Con $n=q$: $q!\,e-\sum_{k=0}^{q}\frac{q!}{k!}=\frac{e^{\xi}}{q+1}$.
  \item El miembro izquierdo es entero: $q!\,e=(q-1)!\,p$, y $q!/k!$ es entero para $k\le q$.
  \item El derecho cumple $0<\frac{e^{\xi}}{q+1}<\frac{3}{q+1}\le1$, porque $q\ge2$.
  \item No hay enteros estrictamente entre $0$ y $1$: contradicción. Luego $e$ es irracional. \hfill$\blacksquare$
\end{enumerate}
\resultado $e\notin\Q$.
\comprobacion la cota $e<3$ se demostró en el \cref{teo:e}; la cota $e^{\xi}>0$ es la que impide que el resto sea nulo.
\end{solucion}

\begin{solucion}{eje:6-criterio-n}
\datos cuatro funciones con punto crítico en~0; dos desigualdades.
\planteamiento (a) primera derivada no nula en~0 y \cref{prop:criterio-n}; (b) resto de Lagrange de orden adecuado con signo conocido.
\pasos
\begin{enumerate}[label=\arabic*.]
  \item $x^{4}$: $f^{(4)}(0)=24>0$, $n=4$ par: mínimo estricto. $x^{5}$: $f^{(5)}(0)=120$, $n=5$ impar: no hay extremo.
  \item $f(x)=\cos x-1+\frac{x^{2}}{2}$: $f'=-\sin x+x$, $f''=1-\cos x$, $f'''=\sin x$ se anulan en~0, y $f^{(4)}(0)=\cos0=1>0$: mínimo estricto.
  \item $f(x)=x\sin x$: $f'=\sin x+x\cos x$ se anula en~0 y $f''=2\cos x-x\sin x$ vale $2>0$: mínimo estricto.
  \item (b) Para $x\ge0$, Taylor de orden~2 de $\sin$ en~0: $\sin x=x-\frac{\cos\xi}{6}x^{3}$, y como $\cos\xi\le1$ y $x^{3}\ge0$, $\sin x\ge x-\frac{x^{3}}{6}$.
  \item Para todo $x$, Taylor de orden~1 de $\cos$ en~0: $\cos x=1-\frac{\cos\xi}{2}x^{2}\ge1-\frac{x^{2}}{2}$. \hfill$\blacksquare$
\end{enumerate}
\resultado mínimos estrictos en~0 para $x^{4}$, $\cos x-1+\frac{x^{2}}{2}$ y $x\sin x$; ningún extremo para $x^{5}$; las dos desigualdades.
\comprobacion la desigualdad $\cos x\ge1-\frac{x^{2}}{2}$ dice que $\cos x-1+\frac{x^{2}}{2}\ge0=f(0)$: el mínimo del paso~2 es, de hecho, absoluto.
\end{solucion}

\begin{solucion}{eje:6-diferencias}
\datos $f\in C^{3}$, $\abs{f''}\le M_2$, $\abs{f'''}\le M_3$.
\planteamiento \cref{teo:taylor} de orden~1 para la diferencia progresiva y de orden~2 en $a+h$ y $a-h$ para la centrada.
\pasos
\begin{enumerate}[label=\arabic*.]
  \item $f(a+h)=f(a)+f'(a)h+\frac{f''(\xi)}{2}h^{2}$, luego $\frac{f(a+h)-f(a)}{h}-f'(a)=\frac{f''(\xi)}{2}h$ y su valor absoluto es $\le\frac{M_2}{2}\abs h$.
  \item $f(a\pm h)=f(a)\pm f'(a)h+\frac{f''(a)}{2}h^{2}\pm\frac{f'''(\xi_\pm)}{6}h^{3}$. Restando, los términos de orden~0 y~2 se cancelan: $f(a+h)-f(a-h)=2f'(a)h+\frac{f'''(\xi_+)+f'''(\xi_-)}{6}h^{3}$.
  \item Dividiendo por $2h$: $\frac{f(a+h)-f(a-h)}{2h}-f'(a)=\frac{f'''(\xi_+)+f'''(\xi_-)}{12}h^{2}$, cuyo valor absoluto es $\le\frac{2M_3}{12}h^{2}=\frac{M_3}{6}h^{2}$.
  \item Numéricamente, con $f=\sin$, $a=1$, $h=0{,}01$ y $M_2=M_3=1$: el error de la progresiva es $-4{,}22\cdot10^{-3}$ (cota $5\cdot10^{-3}$) y el de la centrada, $-9{,}0\cdot10^{-6}$ (cota $1{,}67\cdot10^{-5}$). \hfill$\blacksquare$
\end{enumerate}
\resultado errores $O(h)$ y $O(h^{2})$; la diferencia centrada es preferible: con el mismo $h$, es unas 500 veces más precisa en el ejemplo.
\comprobacion al dividir $h$ por~10, el error de la progresiva se divide por unos 10 y el de la centrada por unos 100; es el análisis de la comprobación de gradientes del módulo~4, donde además interviene el error de redondeo, que crece cuando $h$ disminuye.
\end{solucion}


\backmatter
% =====================================================================
\begin{thebibliography}{99}
\addcontentsline{toc}{chapter}{Bibliografía}
\markboth{Bibliografía}{}
\small

\bibitem{apostol}
T.~M. Apostol.
\newblock \emph{Calculus}, vol.~1: \emph{One-Variable Calculus, with an Introduction to Linear Algebra}, 2.ª ed.
\newblock Wiley, 1967.
\newblock Cap.~3 «Continuous Functions», cap.~4 «Differential Calculus», cap.~6 «The Logarithm, the Exponential, and the Inverse Trigonometric Functions» y cap.~7 «Polynomial Approximations to Functions».

\bibitem{cambA1}
W.~T. Gowers (profesor); D.~Chua (transcripción).
\newblock \emph{Part IA Analysis~I}, apuntes de clase.
\newblock University of Cambridge, Mathematical Tripos, Lent 2015.
\newblock Cap.~2 «Convergence of sequences» (§2.1--2.3), cap.~4 «Continuous functions» (§4.1) y cap.~5 «Differentiability» (§5.1--5.3).
\newblock \url{https://dec41.user.srcf.net/h/IA_L/analysis_i}

\bibitem{tripos2013}
University of Cambridge.
\newblock \emph{Mathematical Tripos Part IA}, Paper~1, 30 de mayo de 2013.
\newblock Preguntas 9D, 10E y 11E (Analysis~I).
\newblock \url{https://www.maths.cam.ac.uk/undergrad/pastpapers/files/2013/PaperIA_1.pdf}

\bibitem{tripos2019}
University of Cambridge.
\newblock \emph{Mathematical Tripos Part IA}, Paper~1, 2019.
\newblock Preguntas 3E, 9D y 10D (Analysis~I).
\newblock \url{https://www.maths.cam.ac.uk/undergrad/pastpapers/files/2019/paperia_1_2019.pdf}

\bibitem{mit1801}
D.~Jerison.
\newblock 18.01SC: Single Variable Calculus.
\newblock Massachusetts Institute of Technology, MIT OpenCourseWare, otoño de 2010. Unidad~1 «Differentiation» (sesiones~1--20 y \emph{Problem Sets}~1--2) y unidad~2 «Applications of Differentiation» (sesiones~23--36 y \emph{Problem Sets}~3--5).
\newblock \url{https://ocw.mit.edu/courses/18-01sc-single-variable-calculus-fall-2010/}

\bibitem{mit1801ex}
D.~Jerison y A.~Mattuck.
\newblock \emph{18.01 Exercises}: secciones 1C--1G y~1J («Differentiation») y 2A, 2B, 2C y~2G («Applications of Differentiation»); exámenes~1 y~2 del curso 18.01SC con sus soluciones.
\newblock MIT OpenCourseWare, 2010.
\newblock \url{https://ocw.mit.edu/courses/18-01sc-single-variable-calculus-fall-2010/}

\bibitem{osc1}
G.~Strang y E.~Herman.
\newblock \emph{Calculus}, vol.~1.
\newblock OpenStax, Rice University, 2016.
\newblock Cap.~1 «Functions and Graphs» (§1.1--1.5), cap.~2 «Limits» (§2.1--2.5), cap.~3 «Derivatives» (§3.1--3.9) y cap.~4 «Applications of Derivatives» (§4.1--4.8).
\newblock \url{https://openstax.org/details/books/calculus-volume-1}

\bibitem{osc2}
G.~Strang y E.~Herman.
\newblock \emph{Calculus}, vol.~2.
\newblock OpenStax, Rice University, 2016.
\newblock §6.3 «Taylor and Maclaurin Series».
\newblock \url{https://openstax.org/details/books/calculus-volume-2}

\bibitem{spivak}
M.~Spivak.
\newblock \emph{Calculus}, 4.ª ed.
\newblock Publish or Perish, 2008.
\newblock Cap.~3--12 (de «Functions» a «Inverse Functions»), cap.~18 «The Logarithm and Exponential Functions», cap.~20 «Approximation by Polynomial Functions» y cap.~22 «Infinite Sequences».

\bibitem{stewart}
J.~Stewart.
\newblock \emph{Calculus: Early Transcendentals}, 8.ª ed.
\newblock Cengage Learning, 2016.
\newblock Cap.~1 «Functions and Models» (§1.1--1.5), cap.~2 «Limits and Derivatives» (§2.1--2.8), cap.~3 «Differentiation Rules» (§3.1--3.11), cap.~4 «Applications of Differentiation» (§4.1--4.5 y~4.7) y cap.~11 «Infinite Sequences and Series».

\end{thebibliography}


\end{document}
